% IV : rédiger un résumé de texte — Français 1re Tech — Cours signé OnBuch+
% Programme officiel MINESEC. Compiler avec : tectonic N20-iv-rediger-resume-texte.tex
\newcommand{\DOCMATIERE}{Français}
\newcommand{\DOCNIVEAU}{1re Tech}
\newcommand{\DOCMODULE}{Français – classe de Première technique}
\newcommand{\DOCLECON}{N20}
\newcommand{\DOCTITRE}{IV : rédiger un résumé de texte}
% OnBuch+ — préambule « cours signés » ENRICHI (compilé avec tectonic / XeLaTeX)
% Système visuel « Pop » d'OnBuch+ : fond crème (jamais blanc), bordures encre
% épaisses, ombres « dures » décalées, titres Archivo Black en capitales.
% Version MATHÉMATIQUES : graphes (pgfplots/tikz), boîtes pédagogiques
% (définition, propriété, méthode, attention, le savais-tu, exercice résolu,
% exercice, TP, bilan), page de garde et signature OnBuch+ ancrée en bas.
%
% Macros attendues dans main.tex AVANT \input{preamble} :
%   \DOCMATIERE  \DOCNIVEAU  \DOCMODULE  \DOCLECON (numéro)  \DOCTITRE
\documentclass[11pt]{article}

\usepackage{fontspec}
\usepackage[french]{babel}
\usepackage[a4paper,margin=2cm,top=3.4cm,bottom=2.8cm,headheight=1.4cm,footskip=1.3cm]{geometry}
\usepackage{amsmath,amssymb,amsfonts}
\usepackage{mathtools} % \xmapsto, \coloneqq... souvent utilisés par les modèles
\usepackage{cancel} % simplifications barrées (bilans de forces)
% \abs{z} (module, valeur absolue) et \norm{u} (norme), utilisés par les modèles
\providecommand{\abs}{}\let\abs\relax\DeclarePairedDelimiter{\abs}{\lvert}{\rvert}
\providecommand{\norm}{}\let\norm\relax\DeclarePairedDelimiter{\norm}{\lVert}{\rVert}
\usepackage[table]{xcolor} % [table] : fournit aussi \rowcolors (alternance de lignes)
\usepackage{pagecolor}
\usepackage{tikz}
\usetikzlibrary{arrows.meta,positioning,calc,shapes.geometric,shapes.misc,shapes.arrows,decorations.pathmorphing,decorations.pathreplacing,decorations.markings,patterns,shadows,mindmap,trees,fit,backgrounds,matrix,intersections,angles,quotes}
\usepackage{pgfplots}
\usepackage[version=4]{mhchem} % équations nucléaires \ce{^{235}_{92}U}
\usepackage[european,siunitx]{circuitikz} % schémas électriques
\pgfplotsset{compat=1.18}
\usepgfplotslibrary{fillbetween,groupplots} % aires entre courbes (calcul intégral)
\usepackage{enumitem}
\usepackage{fancyhdr}
% [most] charge toutes les bibliothèques tcolorbox (listings, minted, raster,
% documentation...) et gonfle inutilement la mémoire TeX sur un document long
% (~300 boîtes) : on ne charge que ce qui est réellement utilisé.
\usepackage[skins,breakable]{tcolorbox}
\usepackage{graphicx}
\usepackage{colortbl}
\usepackage{booktabs}
\usepackage{tabularx}
\usepackage{multirow}
\usepackage{diagbox} % case d'en-tête diagonale des tableaux à double entrée
% Colonnes extensibles centrée / gauche / droite (C, L, R), souvent utilisées par les modèles
\newcolumntype{C}{>{\centering\arraybackslash}X}
\newcolumntype{L}{>{\raggedright\arraybackslash}X}
\newcolumntype{R}{>{\raggedleft\arraybackslash}X}
\usepackage{array}
\usepackage{multicol}
\usepackage{siunitx}
% parse-numbers=false : accepte aussi \SI{2,5 \times 10^{-3}}{...}
\sisetup{locale=FR,output-decimal-marker={,},group-minimum-digits=4,parse-numbers=false}
\DeclareSIUnit\ohm{\ensuremath{\symup{\Omega}}} % U+2126 absent de la police
\DeclareSIUnit\division{div} % réglages d'oscilloscope (ms/div, V/div)
\DeclareSIUnit\tr{tr} % tours (vitesse de rotation, ex. \SI{1500}{\tr\per\minute})
\DeclareSIUnit\molar{M} % concentration molaire (ex. \SI{3}{\molar}, \SI{1}{\micro\molar})
\DeclareSIUnit\an{an} % année (le modèle confond parfois avec \year, un compteur TeX) — ex. \SI{5}{\kilo\meter\per\an}
\DeclareSIUnit\FCFA{FCFA} % franc CFA (ex. \SI{500}{\FCFA\per\kilo\gram})
\DeclareSIUnit\degreeBrix{\textdegree Brix} % teneur en sucre d'un sirop/jus (réfractométrie)
\DeclareSIUnit\month{mois} % le modèle utilise parfois \month, un compteur TeX, comme unité de durée
\DeclareSIUnit\microliter{\micro\liter} % le modèle écrit parfois \microliter en un seul token
\DeclareSIUnit\million{millions} % dénombrement (ex. \SI{15}{\million\per\mL} spermatozoïdes)
\usepackage{qrcode}
% \degree : commande fréquemment utilisée par les modèles pour un angle ou une
% température en dehors de \SI{}{} (non fournie par les paquets chargés).
\providecommand{\degree}{^{\circ}}
% \qed (fin de démonstration), utilisé par les modèles sans amsthm
\providecommand{\qed}{\hfill$\square$}
% \barre : double barre des valeurs interdites dans un tableau de variations (tkz-tab)
\providecommand{\barre}{\ensuremath{\|}}
% environnement proof (amsthm non chargé) utilisé par les modèles
\ifdefined\proof\else\newenvironment{proof}[1][Démonstration]{\par\noindent\textit{#1.}\ }{\qed\par}\fi
\usepackage{lastpage}
\usepackage{hyperref}
\hypersetup{hidelinks}

% ============================================================
% Palette « Pop » OnBuch+ — mêmes hex que la classe P (plus_ui.dart)
% ============================================================
\definecolor{poporange}{HTML}{F59321}
\definecolor{popdark}{HTML}{E07A0C}
\definecolor{popcream}{HTML}{FFF6E8}
\definecolor{popink}{HTML}{1C1714}
\definecolor{poppurple}{HTML}{7A5AE0}
\definecolor{popgreen}{HTML}{2E9E6A}
\definecolor{poppink}{HTML}{E0457B}
\definecolor{popblue}{HTML}{2F7DE1}
\definecolor{popgold}{HTML}{C98A12}
\definecolor{popmuted}{HTML}{8A7F76}
\definecolor{popline}{HTML}{EADFD2}
\definecolor{popgray}{HTML}{8A7F76}
\colorlet{popgrey}{popgray}
% Alias tolérants (le modèle nomme parfois les couleurs différemment)
\colorlet{popwhite}{white}
\colorlet{popblack}{popink}
\colorlet{popred}{poppink}
\colorlet{popyellow}{popgold}
% Teintes claires (fonds de tableaux/encadrés) — le modèle nomme parfois
% les couleurs avec un suffixe L (ex. popblueL) sans les avoir définies.
\colorlet{poporangeL}{poporange!20}
\colorlet{popdarkL}{popdark!20}
\colorlet{poppurpleL}{poppurple!20}
\colorlet{popgreenL}{popgreen!20}
\colorlet{poppinkL}{poppink!20}
\colorlet{popblueL}{popblue!20}
\colorlet{popgoldL}{popgold!20}
\colorlet{popredL}{popred!20}
\colorlet{popgrayL}{popgray!20}
% Teintes claires pour fonds de tableaux / zones
\colorlet{popblueL}{popblue!10!white}
\colorlet{poporangeL}{poporange!14!white}
\colorlet{popgreenL}{popgreen!12!white}
\colorlet{poppurpleL}{poppurple!12!white}
\colorlet{poppinkL}{poppink!12!white}
\colorlet{popgoldL}{popgold!14!white}

\pagecolor{popcream}

% ============================================================
% Typographie
% ============================================================
\defaultfontfeatures{Ligatures=TeX}
\setmainfont{PlusJakartaSans-Medium.ttf}[Path=../../../fonts/,BoldFont=PlusJakartaSans-Bold.ttf]
% Maths en sans-serif (Fira Math) avec chiffres et lettres droites de Plus
% Jakarta, pour que les formules se fondent dans le texte.
\usepackage[math-style=ISO,bold-style=ISO]{unicode-math}
\setmathfont{FiraMath-Regular.otf}[Path=../../../fonts/]
\setmathfont{PlusJakartaSans-Medium.ttf}[Path=../../../fonts/,range={up/num,up/{Latin,latin},\mathup}]
\usepackage{icomma} % virgule décimale sans espace en mode math
% Caractères mathématiques Unicode absents des polices de texte (Plus Jakarta,
% Archivo Black) : rendus par leur équivalent mathématique, en texte comme en
% formule — sinon ils s'affichent en carré vide (ex. « Arithmétique dans ℤ »).
\usepackage{newunicodechar}
\newunicodechar{ℝ}{\ensuremath{\mathbb{R}}}
\newunicodechar{ℤ}{\ensuremath{\mathbb{Z}}}
\newunicodechar{ℂ}{\ensuremath{\mathbb{C}}}
\newunicodechar{ℕ}{\ensuremath{\mathbb{N}}}
\newunicodechar{ℚ}{\ensuremath{\mathbb{Q}}}
\newunicodechar{𝔻}{\ensuremath{\mathbb{D}}}
\newunicodechar{α}{\ensuremath{\alpha}}
\newunicodechar{β}{\ensuremath{\beta}}
\newunicodechar{γ}{\ensuremath{\gamma}}
\newunicodechar{δ}{\ensuremath{\delta}}
\newunicodechar{ε}{\ensuremath{\varepsilon}}
\newunicodechar{θ}{\ensuremath{\theta}}
\newunicodechar{λ}{\ensuremath{\lambda}}
\newunicodechar{μ}{\ensuremath{\mu}}
\newunicodechar{σ}{\ensuremath{\sigma}}
\newunicodechar{τ}{\ensuremath{\tau}}
\newunicodechar{φ}{\ensuremath{\varphi}}
\newunicodechar{ψ}{\ensuremath{\psi}}
\newunicodechar{ω}{\ensuremath{\omega}}
\newunicodechar{Σ}{\ensuremath{\Sigma}}
% majuscules produites par \MakeUppercase dans les titres (α-aminés -> Α)
\newunicodechar{Α}{\ensuremath{\alpha}}
\newunicodechar{Β}{\ensuremath{\beta}}
\newunicodechar{Γ}{\ensuremath{\Gamma}}
\newunicodechar{Δ}{\ensuremath{\Delta}}
\newunicodechar{Ε}{\ensuremath{\varepsilon}}
\newunicodechar{Θ}{\ensuremath{\Theta}}
\newunicodechar{Λ}{\ensuremath{\Lambda}}
\newunicodechar{Μ}{\ensuremath{\mu}}
\newunicodechar{Τ}{\ensuremath{\tau}}
\newunicodechar{Φ}{\ensuremath{\Phi}}
\newunicodechar{Ψ}{\ensuremath{\Psi}}
\newunicodechar{Ω}{\ensuremath{\Omega}}
\newunicodechar{↦}{\ensuremath{\mapsto}}
\newunicodechar{∈}{\ensuremath{\in}}
\newunicodechar{∉}{\ensuremath{\notin}}
\newunicodechar{∧}{\ensuremath{\wedge}}
\newunicodechar{⇔}{\ensuremath{\Leftrightarrow}}
\newunicodechar{⟺}{\ensuremath{\Longleftrightarrow}}
\newunicodechar{⇒}{\ensuremath{\Rightarrow}}
\newunicodechar{⟹}{\ensuremath{\Longrightarrow}}
\newunicodechar{∘}{\ensuremath{\circ}}
\newunicodechar{∪}{\ensuremath{\cup}}
\newunicodechar{∩}{\ensuremath{\cap}}
\newunicodechar{≡}{\ensuremath{\equiv}}
\newunicodechar{⊂}{\ensuremath{\subset}}
\newunicodechar{⊥}{\ensuremath{\perp}}
\newunicodechar{∃}{\ensuremath{\exists}}
\newunicodechar{∀}{\ensuremath{\forall}}
\newunicodechar{∛}{\ensuremath{\sqrt[3]{\,}}}
\newunicodechar{∜}{\ensuremath{\sqrt[4]{\,}}}
\newunicodechar{⃗}{\ensuremath{{}^{\to}}}
\newunicodechar{ⁿ}{\ifmmode^{n}\else\textsuperscript{\ensuremath{n}}\fi}
\newunicodechar{ˣ}{\ifmmode^{x}\else\textsuperscript{\ensuremath{x}}\fi}
\newunicodechar{ᵇ}{\ifmmode^{b}\else\textsuperscript{\ensuremath{b}}\fi}
\newunicodechar{ʳ}{\ifmmode^{r}\else\textsuperscript{\ensuremath{r}}\fi}
\newunicodechar{ᵘ}{\ifmmode^{u}\else\textsuperscript{\ensuremath{u}}\fi}
\newunicodechar{ʸ}{\ifmmode^{y}\else\textsuperscript{\ensuremath{y}}\fi}
\newunicodechar{ᵃ}{\ifmmode^{a}\else\textsuperscript{\ensuremath{a}}\fi}
\newunicodechar{ᵉ}{\ifmmode^{e}\else\textsuperscript{\ensuremath{e}}\fi}
\newunicodechar{ᵏ}{\ifmmode^{k}\else\textsuperscript{\ensuremath{k}}\fi}
\newunicodechar{ᵖ}{\ifmmode^{p}\else\textsuperscript{\ensuremath{p}}\fi}
\newunicodechar{ᴺ}{\ifmmode^{N}\else\textsuperscript{\ensuremath{N}}\fi}
\newunicodechar{ᶜ}{\ifmmode^{c}\else\textsuperscript{\ensuremath{c}}\fi}
\newunicodechar{ʰ}{\ifmmode^{h}\else\textsuperscript{\ensuremath{h}}\fi}
\newunicodechar{ᵠ}{\ifmmode^{q}\else\textsuperscript{\ensuremath{q}}\fi}
\newunicodechar{ₙ}{\ifmmode_{n}\else\textsubscript{\ensuremath{n}}\fi}
\newunicodechar{ₐ}{\ifmmode_{a}\else\textsubscript{\ensuremath{a}}\fi}
\newunicodechar{ᵢ}{\ifmmode_{i}\else\textsubscript{\ensuremath{i}}\fi}
\newunicodechar{ᵣ}{\ifmmode_{r}\else\textsubscript{\ensuremath{r}}\fi}
\newunicodechar{ₖ}{\ifmmode_{k}\else\textsubscript{\ensuremath{k}}\fi}
\newunicodechar{ᵦ}{\ifmmode_{\beta}\else\textsubscript{\ensuremath{\beta}}\fi}
\newunicodechar{ₓ}{\ifmmode_{x}\else\textsubscript{\ensuremath{x}}\fi}
\newfontfamily\popdisplay{ArchivoBlack-Regular.ttf}[Path=../../../fonts/]
\newfontfamily\poplogo{SpaceGrotesk-Bold.ttf}[Path=../../../fonts/]
\newfontfamily\popsemi{PlusJakartaSans-SemiBold.ttf}[Path=../../../fonts/]
\newfontfamily\popxbold{PlusJakartaSans-ExtraBold.ttf}[Path=../../../fonts/]
\newfontfamily\popxboldls{PlusJakartaSans-ExtraBold.ttf}[Path=../../../fonts/,LetterSpace=3.0]

\setlist[enumerate]{leftmargin=1.7em,itemsep=5pt,topsep=4pt}
\setlist[itemize]{leftmargin=1.7em,itemsep=4pt,topsep=4pt,label={\color{poporange}\textbullet}}
\setlength{\parindent}{0pt}
\setlength{\parskip}{5pt}
\renewcommand{\arraystretch}{1.25}

% pgfplots : style Pop par défaut
\pgfplotsset{
  every axis/.append style={axis line style={popink,line width=1pt},tick style={popink},
    grid style={popline,line width=0.5pt},label style={font=\small},tick label style={font=\footnotesize}},
  popcurve/.style={poporange,line width=1.8pt,no markers,smooth},
}
% Même style disponible dans un \draw TikZ simple (hors pgfplots)
\tikzset{popcurve/.style={poporange,line width=1.8pt}}
\tikzset{popgrid/.style={popline,very thin}}
\colorlet{popcurve}{poporange} % le nom du style est parfois utilisé comme couleur

% ============================================================
% Pill / squiggle
% ============================================================
\newcommand{\poppill}[3]{%
  \tikz[baseline=-0.55ex]\node[draw=popink,line width=1.2pt,rounded corners=2.6pt,%
    fill=#2,inner xsep=6.5pt,inner ysep=3pt]{%
    \popxboldls\fontsize{7}{7}\selectfont\color{#3}\MakeUppercase{#1}};%
}
\newcommand{\popsquiggle}[1]{%
  \tikz[baseline=-0.4ex]{
    \draw[#1,line width=2.6pt,line cap=round]
      plot [smooth] coordinates {(0,0.09) (0.34,0.01) (0.68,0.21) (1.02,0.01) (1.36,0.10)};
  }%
}
% Mot-clé surligné (marqueur orange sous le texte)
\newcommand{\cle}[1]{{\popxbold\color{popdark}#1}}

% ============================================================
% En-tête / pied
% ============================================================
\pagestyle{fancy}
\fancyhf{}
\renewcommand{\headrulewidth}{0pt}
\renewcommand{\footrulewidth}{0pt}
\lhead{\raisebox{-0.15ex}{\poplogo\fontsize{13}{13}\selectfont\color{popink}OnBuch\color{poporange}+}}
\rhead{\poppill{\DOCMATIERE\ \textbullet\ \DOCNIVEAU\ \textbullet\ Leçon \DOCLECON}{white}{popink}}
\lfoot{\popxbold\fontsize{7.6}{7.6}\selectfont\color{popmuted}\MakeUppercase{Cours signé OnBuch+}\ \textbullet\ {\popsemi\DOCTITRE}}
\rfoot{\tikz[baseline=-0.6ex]\node[rounded corners=8.5pt,fill=popink,minimum height=17pt,minimum width=17pt,inner xsep=5pt,inner sep=0]{\popxbold\fontsize{8}{8}\selectfont\color{popcream}\thepage\,/\,\pageref*{LastPage}};}

% ============================================================
% Titres
% ============================================================
\newcommand{\chaptitre}[2]{%
  \begin{tcolorbox}[enhanced,colback=poporange,colframe=popink,boxrule=2.6pt,arc=11pt,
      drop shadow=popink,left=15pt,right=15pt,top=12pt,bottom=11pt,before skip=3pt,after skip=18pt]
    \poppill{#2}{popink}{popcream}\par\vspace{9pt}
    {\raggedright\hyphenpenalty=10000\exhyphenpenalty=10000\popdisplay\fontsize{22}{24}\selectfont\color{popcream}\MakeUppercase{#1}\par}
  \end{tcolorbox}
}
% Section numérotée : pastille orange avec le numéro + titre Archivo
\newcounter{popsec}
\newcommand{\coursec}[1]{%
  \stepcounter{popsec}\setcounter{popsub}{0}%
  \par\vspace{16pt}\noindent
  \begin{minipage}{\linewidth}
  \tikz[baseline=-0.7ex]\node[draw=popink,line width=1.6pt,fill=poporange,rounded corners=4pt,minimum size=22pt,inner sep=0]{\popdisplay\fontsize{12}{12}\selectfont\color{popcream}\thepopsec};%
  \hspace{8pt}{\popdisplay\fontsize{15}{17}\selectfont\color{popink}\MakeUppercase{#1}}\par
  \vspace{4pt}\noindent{\color{poporange}\rule{36pt}{3.6pt}}
  \end{minipage}\par\nopagebreak\vspace{8pt}
}
\newcounter{popsub}
\newcommand{\courssub}[1]{%
  \stepcounter{popsub}%
  \par\vspace{9pt}\noindent{\popxbold\fontsize{11.5}{13}\selectfont\color{popink}{\color{poporange}\thepopsec.\thepopsub}\hspace{6pt}#1}\par\nopagebreak\vspace{4pt}
}

% ============================================================
% Boîtes pédagogiques — même recette : fond blanc, bordure épaisse colorée,
% ombre dure encre, badge plein. Titre optionnel [#1].
% ============================================================
\tcbset{popbase/.style={breakable,enhanced,colback=white,boxrule=2.3pt,arc=9pt,
  shadow={3pt}{-3pt}{0pt}{popink},left=14pt,right=14pt,top=11pt,bottom=11pt,before skip=11pt,after skip=15pt,
  fontupper=\popsemi\fontsize{10.6}{14.5}\selectfont\color{popink},
  fontlower=\popsemi\fontsize{10.6}{14.5}\selectfont\color{popink}}}
% Coche dessinée (le glyphe ✓ n'existe pas dans les polices du document)
\newcommand{\popcheck}[1]{\tikz[baseline=-0.5ex]\draw[#1,line width=1.6pt,line cap=round,line join=round](-0.13,0.0)--(-0.04,-0.09)--(0.13,0.1);}
\providecommand{\coche}{\popcheck{popgreen}}
\providecommand{\resultat}[1]{\par\smallskip\noindent\fcolorbox{popgreen}{popgreenL}{\parbox{\dimexpr\linewidth-2\fboxsep-2\fboxrule\relax}{\textbf{Résultat.} #1}}\par\smallskip}
\newcommand{\popboxhead}[3]{\poppill{#1}{#2}{white}\ifx\relax#3\relax\else\hspace{6pt}{\popxbold\fontsize{10.6}{12}\selectfont\color{popink}#3}\fi\par\vspace{7pt}}

\newtcolorbox{aretenir}[1][]{popbase,colframe=popgreen,before upper={\popboxhead{À retenir}{popgreen}{#1}}}
\newtcolorbox{exemplebox}[1][]{popbase,colframe=poporange,before upper={\popboxhead{Exemple}{poporange}{#1}}}
\newtcolorbox{definition}[1][]{popbase,colframe=popblue,before upper={\popboxhead{Définition}{popblue}{#1}}}
\newtcolorbox{propriete}[1][]{popbase,colframe=poppurple,before upper={\popboxhead{Propriété}{poppurple}{#1}}}
\newtcolorbox{methode}[1][]{popbase,colframe=popgold,before upper={\popboxhead{Méthode}{popgold}{#1}}}
\newtcolorbox{attention}[1][]{popbase,colframe=poppink,before upper={\popboxhead{Attention}{poppink}{#1}}}
\newtcolorbox{experience}[1][]{popbase,colframe=popblue,colback=popblueL,before upper={\popboxhead{Expérience}{popblue}{#1}}}
% « Le savais-tu ? » : carte violette pleine (contexte, culture, Cameroun)
\newtcolorbox{savaistu}[1][]{popbase,colback=poppurple,colframe=popink,
  fontupper=\popsemi\fontsize{10.4}{14.2}\selectfont\color{white},
  before upper={\poppill{Le savais-tu\,?}{white}{poppurple}\ifx\relax#1\relax\else\hspace{6pt}{\popxbold\color{white}#1}\fi\par\vspace{7pt}}}
% Exercice résolu : énoncé en haut, \tcblower puis la solution sur fond crème
\newcounter{popexo}\newcounter{popexr}
\newtcolorbox{exoresolu}[1][]{popbase,colframe=poporange,colbacklower=popcream,
  segmentation style={popink,line width=1.2pt,dashed},
  before upper={\stepcounter{popexr}\popboxhead{Exercice résolu \thepopexr}{poporange}{#1}},
  before lower={\poppill{Solution}{popink}{popcream}\par\vspace{6pt}}}
% Exercice (non résolu en place) — la correction est dans la partie Corrigés
\newtcolorbox{exercice}[1][]{popbase,colframe=popink,
  before upper={\stepcounter{popexo}\popboxhead{Exercice \thepopexo}{popink}{#1}}}
% Corrigé
\newtcolorbox{corrige}[1][]{popbase,colframe=popgreen,colback=popgreenL,
  before upper={\popboxhead{Corrigé}{popgreen}{#1}}}
% Objectifs / prérequis / bilan
\newtcolorbox{objectifs}{popbase,colframe=popink,colback=poporangeL,
  before upper={\popboxhead{Objectifs de la leçon}{popink}{}}}
\newtcolorbox{prerequis}{popbase,colframe=popink,colback=popblueL,
  before upper={\popboxhead{Prérequis}{popblue}{}}}
\newtcolorbox{bilan}{popbase,colframe=popink,colback=white,boxrule=2.8pt,
  before upper={\popboxhead{Fiche bilan}{poporange}{L'essentiel à connaître pour l'examen}}}
% Figure encadrée avec légende
\newtcolorbox{popfigure}[1][]{enhanced,breakable=false,colback=white,colframe=popline,boxrule=1.4pt,arc=8pt,
  left=10pt,right=10pt,top=10pt,bottom=8pt,before skip=10pt,after skip=12pt,halign=center,
  lower separated=false,halign lower=center,
  fontlower=\popsemi\fontsize{9}{11.5}\selectfont\color{popmuted},
  #1}
\newcommand{\legende}[1]{\par\vspace{6pt}{\popsemi\fontsize{9}{11.5}\selectfont\color{popmuted}#1}}

% ============================================================
% Signature OnBuch+
% ============================================================
\newcommand{\poplogoinline}[1]{{\poplogo\fontsize{#1}{#1}\selectfont\color{popink}OnBuch\color{poporange}+}}
\newcommand{\signatureonbuch}{%
  \begin{tcolorbox}[enhanced,colback=popink,colframe=popink,boxrule=2.2pt,arc=10pt,
      drop shadow=poporange,left=14pt,right=14pt,top=10pt,bottom=10pt,before skip=0pt,after skip=0pt]
    \tikz[baseline=-0.7ex]\node[circle,fill=popgreen,draw=popcream,line width=1.3pt,minimum size=22pt,inner sep=0]{\popcheck{white}};%
    \hspace{9pt}\begin{minipage}[c]{0.62\linewidth}
      {\popxbold\fontsize{12}{13}\selectfont\color{popcream}Signé\ \textbullet\ {\poplogo OnBuch\color{poporange}+}}\par\vspace{2pt}
      {\raggedright\popsemi\fontsize{8}{10.5}\selectfont\color{popline}\MakeUppercase{Équipe pédagogique OnBuch+}\\\MakeUppercase{Conforme au programme officiel MINESEC}\par}
    \end{minipage}\hfill
    {\poplogo\fontsize{22}{22}\selectfont\color{popcream}OnBuch\color{poporange}+}
  \end{tcolorbox}%
}

% ============================================================
% Page de garde
% #1 titre · #2 matière · #3 niveau · #4 module · #5 n° leçon · #6 durée ·
% #7 description / objectifs (texte ou itemize)
% ============================================================
\newcommand{\pagegarde}[7]{%
  \thispagestyle{empty}
  \newgeometry{a4paper,margin=1.7cm,top=1.6cm,bottom=1.4cm}%
  \begin{tikzpicture}[remember picture,overlay]
    \fill[poporange!18] ([xshift=-3.2cm,yshift=-3.6cm]current page.north east) circle (4.6cm);
    \fill[poppurple!14] ([xshift=2.4cm,yshift=7.6cm]current page.south west) circle (3.8cm);
  \end{tikzpicture}%
  {\poplogo\fontsize{30}{30}\selectfont\color{popink}OnBuch\color{poporange}+}\hfill
  \raisebox{0.6ex}{\poppill{Cours signé \textbullet\ Premium}{popink}{popcream}}\par
  \vspace{1pt}\popsquiggle{poppurple}\par
  \vspace{8pt}
  \begin{tcolorbox}[enhanced,colback=poporange,colframe=popink,boxrule=2.8pt,arc=13pt,
      drop shadow=popink,left=18pt,right=18pt,top=12pt,bottom=13pt,after skip=10pt]
    \poppill{#2 \textbullet\ #3}{popink}{popcream}\hspace{5pt}\poppill{#4}{white}{popink}\par\vspace{8pt}
    {\popdisplay\fontsize{14}{14}\selectfont\color{popink}LEÇON #5}\par\vspace{4pt}
    {\raggedright\hyphenpenalty=10000\exhyphenpenalty=10000\popdisplay\fontsize{25}{27}\selectfont\color{popcream}\MakeUppercase{#1}\par}
    \vspace{8pt}
    {\popxbold\fontsize{9.5}{9.5}\selectfont\color{popink}\MakeUppercase{Durée conseillée : #6}}
  \end{tcolorbox}
  \begin{tcolorbox}[enhanced,colback=white,colframe=popink,boxrule=2.2pt,arc=10pt,
      drop shadow=popink,left=16pt,right=16pt,top=10pt,bottom=10pt,after skip=8pt]
    \poppill{Dans cette leçon}{poppurple}{white}\par\vspace{5pt}
    \popsemi\fontsize{9.3}{12.6}\selectfont\color{popink}#7
  \end{tcolorbox}
  \noindent\begin{minipage}[t]{0.487\linewidth}
    \begin{tcolorbox}[enhanced,colback=popgreen,colframe=popink,boxrule=2.2pt,arc=10pt,
        drop shadow=popink,left=11pt,right=11pt,top=10pt,bottom=10pt,height=3.1cm,valign=center]
      \tcbox[colback=white,colframe=popink,boxrule=1.3pt,arc=5pt,boxsep=0pt,nobeforeafter,
        left=3pt,right=3pt,top=3pt,bottom=3pt,valign=center]{\qrcode[height=1.9cm]{https://play.google.com/store/apps/details?id=com.luvvix.onbuch&hl=fr}}%
      \hspace{8pt}\begin{minipage}[c]{0.5\linewidth}
        {\popdisplay\fontsize{11}{12.5}\selectfont\color{white}\MakeUppercase{Télécharge OnBuch}}\par\vspace{4pt}
        {\raggedright\popsemi\fontsize{8.4}{11}\selectfont\color{white}Gratuit \textbullet\ Google Play\\Tuteur IA, annales, concours, résultats\par}
      \end{minipage}
    \end{tcolorbox}
  \end{minipage}\hfill
  \begin{minipage}[t]{0.487\linewidth}
    \begin{tcolorbox}[enhanced,colback=poppurple,colframe=popink,boxrule=2.2pt,arc=10pt,
        drop shadow=popink,left=11pt,right=11pt,top=10pt,bottom=10pt,height=3.1cm,valign=center]
      \tikz[baseline=-0.7ex]\node[circle,fill=white,draw=popink,line width=1.3pt,minimum size=1.6cm,inner sep=0]{\popdisplay\fontsize{24}{24}\selectfont\color{poppurple}+};%
      \hspace{8pt}\begin{minipage}[c]{0.6\linewidth}
        {\popdisplay\fontsize{11}{12.5}\selectfont\color{white}\MakeUppercase{Passe à OnBuch+}}\par\vspace{4pt}
        {\raggedright\popsemi\fontsize{8.4}{11}\selectfont\color{white}Tous les cours signés, Tuteur IA illimité, fiches PDF — onglet OnBuch+ de l'app\par}
      \end{minipage}
    \end{tcolorbox}
  \end{minipage}
  \vspace{8pt}
  \popcontenu
  \vfill
  \signatureonbuch
  \restoregeometry
  \newpage
}

% Rangée « Ce cours contient » (page de garde)
\newcommand{\poptile}[3]{% #1 couleur #2 gros texte #3 légende
  \begin{tcolorbox}[enhanced,colback=#1,colframe=popink,boxrule=2pt,arc=9pt,shadow={2.5pt}{-2.5pt}{0pt}{popink},
      left=6pt,right=6pt,top=9pt,bottom=9pt,halign=center,valign=center,height=1.75cm,width=\linewidth,nobeforeafter]
    {\popdisplay\fontsize{12.5}{13}\selectfont\color{white}\MakeUppercase{#2}}\par\vspace{3pt}
    {\centering\popsemi\fontsize{7.8}{9.5}\selectfont\color{white}#3\par}
  \end{tcolorbox}}
\newcommand{\popcontenu}{%
  \par\noindent{\popxboldls\fontsize{7.5}{7.5}\selectfont\color{popmuted}CE COURS SIGNÉ CONTIENT}\par\vspace{7pt}
  \noindent\begin{minipage}[t]{0.235\linewidth}\poptile{popblue}{Cours}{complet, illustré, pas à pas}\end{minipage}\hfill
  \begin{minipage}[t]{0.235\linewidth}\poptile{poporange}{Méthodes}{et exercices résolus}\end{minipage}\hfill
  \begin{minipage}[t]{0.235\linewidth}\poptile{poppink}{Exercices}{type examen + corrigés}\end{minipage}\hfill
  \begin{minipage}[t]{0.235\linewidth}\poptile{popgreen}{Fiche bilan}{l'essentiel à retenir}\end{minipage}\par}

% Sommaire
\newtcolorbox{sommaire}{enhanced,colback=white,colframe=popink,boxrule=2.2pt,arc=10pt,
  drop shadow=popink,left=18pt,right=18pt,top=13pt,bottom=13pt,before skip=6pt,after skip=20pt,
  before upper={\poppill{Sommaire}{poppurple}{white}\par\vspace{9pt}\popsemi\fontsize{10.6}{14.5}\selectfont\color{popink}}}

% Signature de fin de cours — poussée tout en bas de la dernière page
\newcommand{\findecours}{%
  \par\vfill\nopagebreak
  \begin{center}\popsquiggle{poporange}\end{center}\vspace{4pt}
  \signatureonbuch
}
\AtBeginDocument{\def\llbracket{\mathopen{[\mkern-3.5mu[}}\def\rrbracket{\mathclose{]\mkern-3.5mu]}}} % intervalles d'entiers (glyphe ⟦ absent de la police)
% Glyphes absents de Fira Math : redéfinis à partir de glyphes disponibles
\AtBeginDocument{%
  \def\setminus{\mathbin{\backslash}}\let\smallsetminus\setminus
  \def\vdots{\mathord{\vbox{\baselineskip3.2pt\lineskiplimit0pt\kern4pt\hbox{.}\hbox{.}\hbox{.}}}}%
  \def\ddots{\mathinner{\mkern1mu\raise6.5pt\hbox{.}\mkern2mu\raise3.6pt\hbox{.}\mkern2mu\raise0.7pt\hbox{.}\mkern1mu}}%
  \def\triangle{\mathord{\scalebox{0.9}{$\symup{\Delta}$}}}%
  \def\nabla{\mathord{\raisebox{\depth}{\scalebox{1}[-1]{$\symup{\Delta}$}}}}%
  \def\checkmark{\popcheck{popdark}}%
  \def\star{\mathbin{\ast}}\def\bigstar{\mathord{\ast}}%
  \def\diamond{\mathbin{\scalebox{0.6}{$\Diamond$}}}%
  \def\bigcup{\mathop{\mathchoice{\raisebox{-0.3em}{\scalebox{1.7}{$\cup$}}}{\scalebox{1.25}{$\cup$}}{\cup}{\cup}}\displaylimits}%
  \def\bigcap{\mathop{\mathchoice{\raisebox{-0.3em}{\scalebox{1.7}{$\cap$}}}{\scalebox{1.25}{$\cap$}}{\cap}{\cap}}\displaylimits}%
  \def\lesseqqgtr{\mathrel{\lessgtr}}\def\bigoplus{\mathop{\oplus}\displaylimits}%
}
\providecommand{\ssi}{si et seulement si} % abréviation fréquente
\AtBeginDocument{\providecommand{\wideparen}{\overparen}} % arc de cercle

%% FIGFIX-BEGIN v4 — correctifs globaux des figures (docs/onbuch-generation/tools/figfix)
\usepackage{adjustbox}\usepackage{etoolbox}
% toute figure TikZ/pgfplots plus large que la ligne est réduite à la largeur disponible
\BeforeBeginEnvironment{tikzpicture}{\begin{adjustbox}{max width=\linewidth}}
\AfterEndEnvironment{tikzpicture}{\end{adjustbox}}
% légendes pgfplots lisibles par-dessus les courbes
\pgfplotsset{every axis/.append style={legend style={fill=white,fill opacity=0.92,text opacity=1,draw=black!15,inner sep=3pt}}}
% étiquettes d'arêtes (syntaxe edge["..."]) avec fond blanc
\tikzset{every edge quotes/.append style={fill=white,inner sep=1.2pt}}
%% FIGFIX-END
\begin{document}

\pagegarde{IV : rédiger un résumé de texte}{Français}{1re Tech}{Français – classe de Première technique}{N20}{10 séances (fiche de progression harmonisée)}{Cette leçon outille l'élève de Première technique pour maîtriser la contraction de texte : repérage des idées essentielles, techniques de réduction (suppression, généralisation, reformulation) et rédaction d'un résumé fidèle, fluide et personnel. Elle s'appuie sur l'étude de l'œuvre 'Le Lion et la perle' de Wole Soyinka et consolide les notions de grammaire de la phrase (types et structures) ainsi que les relations hypéronymie/hyponymie.
\vspace{4pt}
\begin{multicols}{2}\small
\begin{itemize}
\item À la fin de cette leçon, je sais présenter un exposé sur 'Le Lion et la perle' (auteur, contexte, thèmes).
\item À la fin de cette leçon, je sais définir le résumé et distinguer résumé, compte rendu et synthèse.
\item À la fin de cette leçon, je sais repérer les idées essentielles d'un texte et hiérarchiser l'information.
\item À la fin de cette leçon, je sais appliquer les techniques de réduction : suppression, généralisation (hyponyme → hypéronyme), reformulation.
\item À la fin de cette leçon, je sais identifier et produire les types de phrases (déclarative, interrogative, impérative, exclamative).
\item À la fin de cette leçon, je sais analyser la structure d'une phrase (simple, composée, complexe) pour la condenser correctement.
\end{itemize}\end{multicols}}

\begin{sommaire}
\begin{multicols}{2}
\begin{enumerate}
\item L'œuvre support : Le Lion et la perle (exposé et contexte)
\item Le résumé : définition, objectifs et contraintes
\item Grammaire outil 1 : les types de phrases
\item Grammaire outil 2 : structure de la phrase et relations lexicales
\item Les techniques de réduction du texte
\item Rédaction, amélioration et évaluation du résumé
\item Activité d'intégration
\item Méthodes et exercices résolus
\item Exercices
\item Corrigés des exercices
\item Fiche bilan
\end{enumerate}
\end{multicols}
\end{sommaire}

\begin{objectifs}
\begin{itemize}
\item À la fin de cette leçon, je sais présenter un exposé sur 'Le Lion et la perle' (auteur, contexte, thèmes).
\item À la fin de cette leçon, je sais définir le résumé et distinguer résumé, compte rendu et synthèse.
\item À la fin de cette leçon, je sais repérer les idées essentielles d'un texte et hiérarchiser l'information.
\item À la fin de cette leçon, je sais appliquer les techniques de réduction : suppression, généralisation (hyponyme → hypéronyme), reformulation.
\item À la fin de cette leçon, je sais identifier et produire les types de phrases (déclarative, interrogative, impérative, exclamative).
\item À la fin de cette leçon, je sais analyser la structure d'une phrase (simple, composée, complexe) pour la condenser correctement.
\item À la fin de cette leçon, je sais rédiger un résumé respectant la consigne de longueur, la fidélité au texte et la correction de la langue.
\item À la fin de cette leçon, je sais réviser et améliorer mon résumé (relecture, comptage des mots, connecteurs).
\item À la fin de cette leçon, je sais justifier ma démarche de contraction devant un texte donné.
\end{itemize}
\end{objectifs}

\begin{prerequis}
\begin{itemize}
\item Lecture et compréhension globale d'un texte narratif et argumentatif (acquis de Seconde)
\item Notion de phrase : sujet, verbe, compléments (classes de collège)
\item Champs lexicaux et relations de sens (synonymie, antonymie) vus en Seconde
\item Lecture suivie de l'œuvre 'Le Lion et la perle' (leçons précédentes de l'unité)
\item Paragraphe et connecteurs logiques (production d'écrits, Seconde)
\end{itemize}
\end{prerequis}

\newpage

\coursec{L'œuvre support : Le Lion et la perle (exposé et contexte)}

\textbf{Situation d'accroche.} Le chef de village de Nkongsamba doit transmettre au sous-préfet, en dix lignes maximum, le contenu d'un rapport de quatre pages sur l'état de la route rurale : comment dire l'essentiel sans rien trahir ? De même, au lycée, un élève qui a lu \emph{Le Lion et la perle} de Wole Soyinka doit pouvoir en présenter l'essentiel à un camarade absent. Au Baccalauréat technique, l'épreuve de contraction de texte exige exactement cette compétence : réduire un texte au quart de sa longueur tout en respectant les idées de l'auteur. Mais avant de réduire un texte, il faut le connaître parfaitement : c'est pourquoi notre leçon commence par la présentation complète de l'œuvre qui nous servira de support.

\textbf{Problématique.} Comment présenter de manière objective et organisée une œuvre littéraire — son auteur, son intrigue, ses thèmes et son contexte — afin de disposer d'une base solide pour apprendre ensuite à la résumer méthodiquement ?

\courssub{Présentation de l'auteur : Wole Soyinka}

\begin{definition}[L'exposé]
Un \cle{exposé} est une présentation orale ou écrite, \textbf{objective}, \textbf{structurée} et \textbf{organisée}, d'une œuvre, d'un auteur ou d'un sujet. Il informe sans donner d'opinion personnelle : il suit un plan clair (introduction, développement en parties, conclusion) et s'appuie sur des faits vérifiables.
\end{definition}

Commençons par l'intuition : pour comprendre une pièce de théâtre, il faut d'abord savoir \emph{qui} l'a écrite et \emph{d'où} parle l'auteur. Un texte n'est jamais neutre : il porte l'expérience de son créateur.

\cle{Wole Soyinka} est un dramaturge, poète et romancier \textbf{nigérian}, né en 1934 à Abeokuta, dans l'ouest du Nigeria, en pays yoruba. Formé à l'University College d'Ibadan (Nigeria) puis à l'université de Leeds (Grande-Bretagne), il fonde des troupes de théâtre (notamment le « 1960 Masks ») et écrit des pièces qui mêlent la tradition orale yoruba (masques, chants, danse, proverbes) et les techniques du théâtre moderne. Engagé politiquement, il est emprisonné pendant la guerre du Biafra (1967-1970) pour avoir dénoncé les violences. En 1986, il reçoit le \textbf{prix Nobel de littérature}.

\begin{savaistu}[Un Nobel africain]
Wole Soyinka est le \cle{premier Africain} à avoir reçu le prix Nobel de littérature, en 1986. Le jury a salué un écrivain qui, « dans une perspective culturelle large et avec des nuances poétiques, dessine le drame de l'existence ». Cette consécration mondiale prouve que les littératures africaines ont leur place au sommet des lettres mondiales.
\end{savaistu}

\begin{popfigure}
\begin{center}
\begin{tikzpicture}[scale=1]
\draw[popline, line width=1.5pt, -{Stealth[length=3mm]}] (0,0) -- (10.5,0);
\foreach \x/\y/\titre/\txt in {1/1934/{Naissance à Abeokuta}, 5/1959/{Création de \emph{Le Lion et la perle}}, 9/1986/{Prix Nobel de littérature}}{
  \fill[popblue] (\x,0) circle (0.12);
  \node[below, font=\small\bfseries, text=popdark] at (\x,-0.15) {\y};
  \node[above, font=\small, text width=3cm, align=center] at (\x,0.25) {\titre};
}
\end{tikzpicture}
\end{center}
\legende{Frise chronologique : trois dates clés de la vie de Wole Soyinka et de son œuvre.}
\end{popfigure}

\courssub{Genre de l'œuvre : une comédie satirique}

\emph{Le Lion et la perle} (titre original : \emph{The Lion and the Jewel}, 1959) est une \cle{pièce de théâtre} en trois actes, intitulés symboliquement « Matin », « Midi » et « Nuit ». C'est une \textbf{comédie} : elle fait rire par les situations, les quiproquos et les jeux de langage. Mais c'est aussi une \textbf{satire} : derrière le rire, Soyinka se moque à la fois du traditionalisme borné et du modernisme superficiel. Le rire devient alors une arme critique : c'est le propre du registre \cle{satirique}.

\courssub{Résumé de l'intrigue : le Lion, la Perle et l'instituteur}

L'action se déroule à Ilujinle, un village yoruba imaginaire.

\begin{itemize}
  \item \textbf{Matin.} \cle{Lakunle}, jeune instituteur occidentalisé, courtise \cle{Sidi}, la plus belle fille du village, surnommée « la perle ». Il refuse de payer le prix de la fiancée (la dot), qu'il juge arriéré. Or un photographe a publié les portraits de Sidi dans un magazine : sa beauté est désormais célèbre.
  \item \textbf{Midi.} \cle{Baroka}, le vieux chef du village (le Bale), surnommé « le Lion », voit les photos et décide d'épouser Sidi. Il fait courir le bruit qu'il est devenu impuissant pour l'attirer sans méfiance chez lui.
  \item \textbf{Nuit.} Sidi, venue se moquer du vieillard, tombe dans son piège : Baroka, qui n'est pas impuissant, la séduit. Lakunle croit encore pouvoir l'épouser « souillée » et sans dot, mais Sidi choisit Baroka. La pièce se clôt sur la danse de la mariée.
\end{itemize}

\begin{popfigure}
\begin{center}
\begin{tikzpicture}[scale=1, every node/.style={font=\small}]
\node[draw=popblue, fill=popblueL, rounded corners, thick, minimum width=3.2cm, minimum height=1cm, align=center] (baroka) at (0,3){\textbf{Baroka} (le Lion)\\ Chef traditionnel (Bale)};
\node[draw=poppink, fill=poppinkL, rounded corners, thick, minimum width=3.2cm, minimum height=1cm, align=center] (sidi) at (5,3){\textbf{Sidi} (la Perle)\\ Belle du village};
\node[draw=popgreen, fill=popgreenL, rounded corners, thick, minimum width=3.2cm, minimum height=1cm, align=center] (lakunle) at (2.5,0){\textbf{Lakunle}\\ Instituteur moderniste};
\draw[-{Stealth}, thick, popblue] (baroka) -- node[above, sloped, pos=0.5]{désire épouser} (sidi);
\draw[-{Stealth}, thick, popgreen] (lakunle) -- node[below, sloped, pos=0.5]{courtise (sans dot)} (sidi);
\draw[{Stealth}-{Stealth}, thick, poporange, dashed] (baroka) -- node[left, text width=2.4cm, align=center]{Conflit : tradition vs modernité} (lakunle);
\draw[-{Stealth}, thick, poppink] (sidi) -- node[above, sloped, pos=0.5]{choisit} (baroka);
\draw[-{Stealth}, thick, popmuted, densely dotted] (sidi) -- node[below, sloped, pos=0.5]{rejet} (lakunle);
\end{tikzpicture}
\end{center}
\legende{Triangle relationnel : Baroka et Lakunle s'affrontent pour Sidi ; elle choisit le Lion (flèche rose) et rejette l'instituteur (flèche pointillée grise).}
\end{popfigure}

\courssub{Les thèmes majeurs}

\begin{itemize}
  \item \textbf{Tradition contre modernité.} Baroka incarne la coutume (dot, polygamie, autorité du chef), Lakunle le progrès importé (école, machines, rejet de la dot). Soyinka ne donne raison à aucun des deux : Baroka est rusé et manipulateur, Lakunle est ridicule et prétentieux.
  \item \textbf{Ruse et pouvoir.} Le Lion gagne par l'intelligence et le mensonge, non par la force : le pouvoir traditionnel sait s'adapter pour survivre.
  \item \textbf{Condition de la femme africaine.} Sidi est d'abord un « objet » disputé entre deux hommes ; mais son choix final montre qu'elle n'est pas un simple trophée : la pièce interroge la place de la femme dans une société en mutation.
\end{itemize}

\courssub{Inscription de l'œuvre dans son contexte}

La pièce naît dans le \textbf{Nigeria des années 1950-1960}, à la veille de l'indépendance (1er octobre 1960). Le pays vit alors une tension profonde entre les \cle{coutumes} locales (chefferies, religions traditionnelles, oralité) et l'\cle{occidentalisation} héritée de la colonisation britannique (école, administration, christianisme, techniques modernes). Soyinka met en scène ce débat à l'échelle d'un village : la photographie, le magazine, le projet de chemin de fer symbolisent la modernité qui pénètre Ilujinle. La victoire de Baroka suggère que la tradition nigériane, loin d'être figée, sait absorber le moderne à son profit.

\courssub{Bilan de l'étude de l'œuvre}

\begin{center}
\begin{tabularx}{\linewidth}{|l|X|}\hline
\rowcolor{popblueL} \textbf{Aspect} & \textbf{Bilan} \\ \hline
Personnages & Trio central fortement typé : le vieux chef rusé, la belle vaniteuse mais lucide, l'instituteur pédant. \\ \hline
Structure & Trois actes (Matin, Midi, Nuit) : progression dramatique claire, coup de théâtre final. \\ \hline
Registre & \textbf{Satirique} : humour, ironie, caricature pour critiquer les deux camps. \\ \hline
Langue et humour & Mélange de prose et de vers, chants et danses yoruba, comique de mots et de situation. \\ \hline
Portée & Satire équilibrée : ni rejet de la tradition, ni culte aveugle du progrès. \\ \hline
\end{tabularx}
\end{center}

\begin{methode}[Construire un exposé littéraire]
\begin{enumerate}
  \item Présenter l'\textbf{auteur} (nom, nationalité, dates, place dans la littérature).
  \item Présenter l'\textbf{œuvre} (titre, date, genre, registre).
  \item Donner un \textbf{résumé} bref de l'intrigue (sans dévoiler la fin si exposé oral de découverte).
  \item Dégager les \textbf{thèmes} majeurs.
  \item Situer l'œuvre dans son \textbf{contexte} historique et culturel.
  \item Conclure par une \textbf{appréciation} mesurée (intérêt, portée de l'œuvre).
\end{enumerate}
\end{methode}

\begin{attention}[Résumé ou analyse ?]
Ne pas confondre le \textbf{résumé} d'une œuvre (reformulation brève et fidèle des événements, sans commentaire) et son \textbf{analyse critique} (jugement argumenté sur la valeur, le style, les intentions). Dans un exposé, les deux ont leur place mais à des moments distincts ; dans une contraction de texte, tout jugement personnel est interdit.
\end{attention}

\begin{exoresolu}[Rédiger le paragraphe « auteur » d'un exposé]
Énoncé : à partir des informations du cours, rédigez en quatre ou cinq lignes le paragraphe de présentation de Wole Soyinka pour un exposé oral.
\tcblower
Solution : « Wole Soyinka est un écrivain nigérian né en 1934 à Abeokuta, en pays yoruba. Dramaturge, poète et romancier, il mêle dans ses œuvres la tradition orale africaine et le théâtre moderne. Engagé pour la justice, il fut emprisonné pendant la guerre du Biafra. En 1986, il devient le premier Africain lauréat du prix Nobel de littérature. » Ce paragraphe est objectif, organisé (identité, œuvre, engagement, consécration) et sans opinion personnelle : c'est un modèle d'exposé.
\end{exoresolu}

\textbf{Lien avec la suite de la leçon.} Désormais, l'œuvre est connue : auteur, intrigue, thèmes, contexte. Les extraits de \emph{Le Lion et la perle} nous serviront de \cle{textes d'application} pour apprendre, dans les prochaines séances, à réduire un texte puis à rédiger un résumé conforme aux exigences du Baccalauréat technique.

\coursec{Le résumé : définition, objectifs et contraintes}

Après avoir découvert \emph{Le Lion et la perle} de Wole Soyinka à travers un exposé structuré et une mise en contexte historique et culturelle (Nigéria postcolonial, conflit tradition/modernité), nous disposons désormais d'un texte support riche et complexe. C'est précisément sur ce type de texte — littéraire, argumentatif ou documentaire — que s'exerce l'épreuve de contraction. Savoir résumer, ce n'est pas seulement « raccourcir » : c'est opérer un travail intellectuel rigoureux de sélection, de restructuration et de reformulation. Cette section pose les fondations théoriques et méthodologiques indispensables avant d'aborder les techniques de réduction proprement dites.

\courssub{Définition et nature de la contraction de texte}

\begin{definition}[Contraction de texte (ou résumé)]
La \cle{contraction de texte} est un exercice de réécriture qui consiste à réduire un texte source à une fraction imposée de sa longueur (généralement le quart, parfois le tiers ou le cinquième) tout en conservant \textbf{uniquement les idées essentielles}, dans l'ordre logique de l'auteur, sans en trahir le sens, ni y ajouter d'éléments personnels. Le produit final est le \cle{résumé}.
\end{definition}

La contraction n'est pas une simple suppression de mots au hasard. C'est une \textbf{opération de transformation} qui fait passer le texte d'un état « développé » (explicatif, illustré, redondant) à un état « concentré » (dense, informatif, autonome). Le résumé doit pouvoir se lire indépendamment du texte original : le correcteur (ou le lecteur) doit en comprendre la teneur sans avoir lu la source.

\begin{popfigure}
\begin{tikzpicture}[
    node distance=1.2cm and 1.5cm,
    box/.style={draw=popblue, thick, rounded corners, fill=popblue!10, minimum width=3.8cm, minimum height=1.8cm, align=center, font=\small},
    arrow/.style={-Stealth, thick, popdark},
    labelbox/.style={font=\small\bfseries, popdark, text width=3cm, align=center}
]
% Nœuds
\node[box, align=center] (source){Texte source\\ \textbf{400 mots}\\(exposé, exemples,\\ détails, transitions)};
\node[box, below=of source, align=center] (idees){Repérage des\\ \textbf{idées essentielles}\\(sélection, hiérarchisation,\\ élimination du superflu)};
\node[box, below=of idees, align=center] (resume){Résumé final\\ \textbf{100 mots}\\(idées reformulées,\\ enchaînées logiquement,\\ autonomie syntaxique)};

% Flèches
\draw[arrow] (source.south) -- (idees.north) node[midway, right=0.3cm, labelbox, align=center]{Lecture analytique \\ \& Réduction};
\draw[arrow] (idees.south) -- (resume.north) node[midway, right=0.3cm, labelbox, align=center]{Rédaction \\ \& Reformulation};

% Entonnoir visuel (décoration)
\draw[poporange, thick, decorate, decoration={zigzag, segment length=4mm, amplitude=1.5mm}] 
    ([xshift=-2.5cm]source.north west) -- ([xshift=-2.5cm]resume.south west);
\draw[poporange, thick, decorate, decoration={zigzag, segment length=4mm, amplitude=1.5mm}] 
    ([xshift=2.5cm]source.north east) -- ([xshift=2.5cm]resume.south east);
\end{tikzpicture}
\legende{Schéma du processus de contraction : du texte source (400 mots) vers le résumé (100 mots) par l'entonnoir de la sélection et de la reformulation.}
\end{popfigure}

\courssub{Distinction entre résumé, compte rendu, synthèse et commentaire}

Ces quatre exercices sont souvent confondus. Pourtant, chacun répond à une commande, une finalité et des contraintes formelles distinctes. Le tableau ci-dessous en fixe les frontières nettes.

\begin{popfigure}
\begin{center}
\begin{tabularx}{\linewidth}{|>{\bfseries}X|X|X|X|X|}
\hline
\rowcolor{popblueL}
\textbf{Critère} & \textbf{Résumé (Contraction)} & \textbf{Compte rendu} & \textbf{Synthèse} & \textbf{Commentaire} \\
\hline
\textbf{But} & Restituer fidèlement la pensée de l'auteur dans une longueur imposée. & Rendre compte d'un ouvrage, d'une conférence, d'un film (structure + contenu + appréciation). & Fusionner plusieurs textes sur un même thème pour en dégager les convergences/divergences. & Analyser, interpréter, expliquer le fonctionnement d'un texte (style, arguments, effets). \\
\hline
\textbf{Source} & Un \textbf{seul} texte. & Une \textbf{seule} œuvre/événement. & \textbf{Plusieurs} textes (dossier). & Un \textbf{seul} texte (souvent court). \\
\hline
\textbf{Longueur} & \textbf{Imposée} (ex: 1/4, $\pm 10\%$). & Libre (souvent 1/3 à 1/2). & Imposée ou libre selon consigne. & Libre (développement argumenté). \\
\hline
\textbf{Point de vue} & \textbf{Strictement celui de l'auteur} (disparition du « je » résumeur). & Celui du rédacteur (il peut donner son avis en conclusion). & Neutre, comparatif (pas d'avis personnel). & Celui de l'analyste (thèse personnelle argumentée). \\
\hline
\textbf{Structure} & Suit l'ordre logique du texte source. & Suit la structure de l'œuvre + partie critique. & Thématique (par idées, pas par textes). & Sujet imposé ou problème posé (plan dialectique/thématique). \\
\hline
\textbf{Citations} & \textbf{Interdites} (reformulation obligatoire). & Autorisées (courtes, pour illustrer). & Autorisées (courtes, pour étayer). & \textbf{Obligatoires} (preuves textuelles). \\
\hline
\end{tabularx}
\end{center}
\legende{Tableau comparatif : résumé, compte rendu, synthèse et commentaire. Au Probatoire technique, seule la contraction (résumé) est évaluée à l'écrit.}
\end{popfigure}

\begin{attention}[Piège fréquent : confusion résumé / commentaire]
Ne cherchez \textbf{jamais} à « expliquer » le texte dans un résumé. Pas de « L'auteur montre que... », « Cette image signifie... », « On comprend que... ». Le résumé \textbf{restitue} : « L'auteur affirme que... », « Le texte développe l'idée que... ». L'analyse appartient au commentaire composé.
\end{attention}

\courssub{Les trois qualités d'un bon résumé : fidélité, brièveté, clarté}

Un résumé réussi repose sur un triangle d'exigences indissociables. L'absence de l'une entraîne l'échec de l'exercice.

\begin{aretenir}[Les trois piliers du résumé]
\begin{itemize}
    \item \textbf{Fidélité (Respect du sens) :} Aucune idée essentielle omise, aucune idée secondaire promue au rang d'essentielle, aucune déformation du point de vue de l'auteur. C'est la condition \emph{sine qua non}.
    \item \textbf{Brièveté (Respect de la contrainte quantitative) :} Le nombre de mots doit tomber dans l'intervalle toléré (généralement $\pm 10\%$ de la cible). Un résumé trop long sanctionne l'incapacité à sélectionner ; trop court, l'incapacité à développer l'essentiel.
    \item \textbf{Clarté (Qualité de la rédaction) :} Le résumé est un \textbf{texte à part entière} : phrases complètes, connecteurs logiques, cohérence interne, orthographe/grammaire irréprochables. Ce n'est pas une liste de mots-clés ni un télégraphe.
\end{itemize}
\end{aretenir}

\courssub{Outils grammaticaux mobilisés pour la réduction (Rappel du programme)}

La reformulation et la condensation s'appuient sur des manipulations syntaxiques et lexicales précises, étudiées dans les séquences « Langue française » de cette unité. Maîtriser ces outils, c'est gagner en concision sans perdre en sens.

\begin{aretenir}[Leviers grammaticaux de la contraction]
\begin{itemize}
    \item \textbf{Types de phrases (Morphosyntaxe) :} Passage de la phrase verbale à la \textbf{phrase nominale} (nominalisation : « Il décide de partir » $\rightarrow$ « Son départ ») ; utilisation de la \textbf{phrase elliptique} pour les enchaînements évidents.
    \item \textbf{Structure de la phrase :} Transformation de \textbf{phrases simples} juxtaposées en \textbf{phrase complexe} (subordination : relative, conjonctive, participiale) pour condenser l'information (« Il arrive. Il est fatigué. Il dort » $\rightarrow$ « Arrivé fatigué, il dort »).
    \item \textbf{Hypéronymie / Hyponymie (Sémantique) :} Remplacer une énumération d'\textbf{hyponymes} (mots spécifiques : \emph{chien, chat, cheval}) par leur \textbf{hyperonyme} (mot générique : \emph{animaux domestiques}) pour généraliser et gagner des mots.
\end{itemize}
\end{aretenir}

\courssub{Ce qu'il faut bannir : les « déchets » du texte source}

La réduction opère un tri radical. Tout ce qui n'est pas « idée force » ou « articulation logique » doit disparaître.

\begin{exemplebox}[Éléments à éliminer systématiquement]
\begin{enumerate}
    \item \textbf{Les exemples et illustrations :} « Par exemple », « c'est le cas de... », « ainsi que le montre l'histoire de... » $\rightarrow$ \textbf{Supprimer}.
    \item \textbf{Les répétitions et reformulations internes :} L'auteur dit souvent la même chose avec des mots différents pour insister. $\rightarrow$ \textbf{Garder une seule occurrence}.
    \item \textbf{Les citations longues :} Remplacer par une reformulation indirecte (style indirect libre). $\rightarrow$ \textbf{Intégrer l'idée, jeter les guillemets}.
    \item \textbf{Les détails secondaires :} Dates précises, noms propres anecdotiques, chiffres illustratifs, descriptions pittoresques. $\rightarrow$ \textbf{Supprimer ou généraliser} (« au XXe siècle » au lieu de « en 1962 »).
    \item \textbf{Les opinions personnelles du résumeur :} « Je pense que... », « C'est intéressant de noter que... », « Cependant, on peut contester... » $\rightarrow$ \textbf{Interdit absolu}.
    \item \textbf{Les connecteurs de pure politesse ou de transition lourde :} « En ce qui concerne... », « Il est important de souligner que... » $\rightarrow$ \textbf{Remplacer par des connecteurs forts} (« Or », « Donc », « Par ailleurs »).
\end{enumerate}
\end{exemplebox}

\begin{attention}[Erreur fatale : le « copier-coller » masqué]
Recopier des phrases entières du texte original en changeant juste deux mots (ex: « Il est évident que » $\rightarrow$ « C'est évident que ») n'est \textbf{pas} une reformulation. C'est du plagiat scolaire. Le correcteur le repère instantanément (même structure syntaxique, même vocabulaire). \textbf{Sanction : note sévère, voire zéro à l'exercice.} La reformulation impose de changer \textbf{la structure de la phrase} (subordination, nominalisation) et \textbf{le lexique} (synonymes, hyperonymes, verbalisation/nominalisation).
\end{attention}

\courssub{Les étapes de la démarche : un processus en cinq temps}

La contraction ne s'improvise pas. Elle suit un protocole rigoureux que tout candidat doit automatiser.

\begin{methode}[Démarche type de la contraction de texte]
\begin{enumerate}
    \item \textbf{Lecture analytique (15--20\% du temps) :} Lire le texte 2 à 3 fois. Identifier : le thème, la thèse, le plan (grandes parties), le ton, le public visé. Annoter dans la marge : idée force par paragraphe.
    \item \textbf{Repérage et sélection des idées (25--30\% du temps) :} Lister les idées essentielles (une par paragraphe ou groupe de paragraphes). \textbf{Barrez} au brouillon exemples, répétitions, détails. Hiérarchisez : idée principale $>$ arguments majeurs $>$ articulations logiques.
    \item \textbf{Réduction / Réécriture au brouillon (30--35\% du temps) :} Reformulez chaque idée sélectionnée \textbf{avec vos mots}. \textbf{Appliquez les outils grammaticaux} : nominalisez les verbes, subordonnez les phrases simples, généralisez par hyperonymie, passez au style indirect. Visez 80--90\% de la cible au brouillon (marge de manœuvre pour l'ajustement final).
    \item \textbf{Rédaction finale soignée (15--20\% du temps) :} Recopiez proprement sur la copie d'examen. Assurez la fluidité (connecteurs), l'autonomie (pas de « il » ambigu), la correction linguistique.
    \item \textbf{Relecture et comptage (5--10\% du temps) :} \textbf{Comptez les mots} (voir méthode ci-dessous). Vérifiez : fidélité ? brièveté (dans la fourchette) ? clarté (sens, grammaire, orthographe) ? Ajustez si nécessaire (couper un adjectif, fusionner deux phrases).
\end{enumerate}
\end{methode}

\courssub{Le comptage des mots : méthode et tolérance}

Au Probatoire et au Baccalauréat technique, la contrainte de longueur est \textbf{stricte}. Le non-respect de la fourchette entraîne une pénalité lourde (souvent 2 à 4 points sur 20, voire note éliminatoire si hors tolérance).

\begin{methode}[Compter ses mots rapidement et exactement]
\begin{enumerate}
    \item \textbf{Convention officielle :} Un mot = toute suite de caractères séparée par un espace. Les contractions comptent pour un mot (\emph{l'homme}, \emph{qu'il}, \emph{au} = 1 mot). Les nombres écrits en chiffres (\emph{2026}) comptent pour 1 mot. La ponctuation ne compte pas.
    \item \textbf{Méthode de la ligne moyenne (estimation rapide) :} Comptez les mots de la 3e ligne rédigée $\times$ nombre total de lignes écrites. Donne une estimation $\pm 5\%$.
    \item \textbf{Méthode exacte (obligatoire en fin d'épreuve) :} Comptez \textbf{chaque mot} en pointant du stylo sur la copie finale. Notez le total en marge. Recomptez une seconde fois dans le sens inverse pour vérifier.
    \item \textbf{Calcul de la cible et de la tolérance :}
    \begin{itemize}
        \item Cible = (Nombre de mots du texte source) $\div$ 4 (ou diviseur indiqué par la consigne).
        \item Tolérance $\pm 10\%$ : Intervalle = [Cible $\times$ 0,9 ; Cible $\times$ 1,1].
    \end{itemize}
\end{enumerate}
\end{methode}

\begin{exoresolu}[Calcul de la fourchette autorisée]
\textbf{Énoncé :} Un texte source comporte \textbf{320 mots}. La consigne demande un résumé au \textbf{quart} avec une tolérance de \textbf{$\pm 10\%$}. Déterminer le nombre de mots cible et l'intervalle acceptable.

\textbf{Solution détaillée :}
\begin{align*}
\text{Cible (quart)} &= 320 \div 4 = \textbf{80 mots}. \\
\text{Marge basse (-10\%)} &= 80 \times 0,9 = \textbf{72 mots}. \\
\text{Marge haute (+10\%)} &= 80 \times 1,1 = \textbf{88 mots}. \\
\text{Intervalle acceptable} &= \textbf{[72 mots ; 88 mots]}.
\end{align*}
\textbf{Conséquence :} Tout résumé de 71 mots ou moins, ou de 89 mots ou plus, est \textbf{hors tolérance} et pénalisé. Visez \textbf{80 mots exactement} pour être au centre de la cible.
\tcblower
\textbf{Astuce :} Si votre brouillon fait 85 mots, coupez 5 mots (adjectifs, adverbes, connecteurs lourds). S'il fait 75 mots, développez légèrement une idée force (ajoutez un complément de nom, précisez un verbe).
\end{exoresolu}

\courssub{Enjeu pour l'examen : la contraction au Probatoire et au Baccalauréat technique}

\begin{savaistu}[L'épreuve écrite de Français — Série Technique]
Au Cameroun, l'épreuve de Français au \textbf{Probatoire technique} (fin de Première) et au \textbf{Baccalauréat technique} (fin de Terminale) comporte obligatoirement une \textbf{contraction de texte} (coefficient variable selon les séries, souvent 2 à 4 sur la note totale de l'écrit).
\begin{itemize}
    \item \textbf{Nature du texte support :} Texte argumentatif, documentaire ou littéraire (extrait de roman, pièce, essai), longueur 300--450 mots.
    \item \textbf{Consigne type :} « Résumez le texte ci-dessus en 80 mots environ (quart du texte), tolérance $\pm 10\%$. »
    \item \textbf{Notation type (sur 20 pts) :} Fidélité (6--8 pts) + Brièveté/Respect contrainte (4--6 pts) + Clarté/Langue (6--8 pts).
    \item \textbf{Compétences transverses mobilisées :} Morphosyntaxe (types de phrases, subordination pour condenser), Lexique (hyperonymie/hyponymie pour généraliser), Cohérence textuelle (connecteurs, référents).
\end{itemize}
Maîtriser la contraction, c'est sécuriser \textbf{15 à 20\% de la note finale} de l'écrit. C'est aussi un outil professionnel : compte rendu de stage, synthèse de notes techniques, rapport d'intervention.
\end{savaistu}

\begin{aretenir}[Check-list « Prêt pour l'examen »]
Avant de rendre votre copie, cochez mentalement :
\begin{itemize}
    \item [ ] Nombre de mots compté \textbf{deux fois} et dans l'intervalle [Cible $\pm 10\%$] ?
    \item [ ] Aucune phrase recopiée du texte ? (Reformulation totale : structure + lexique)
    \item [ ] Aucune opinion personnelle ? (Disparition du « je »)
    \item [ ] Toutes les idées forces présentes ? (Fidélité)
    \item [ ] Texte fluide, connecté, sans fautes graves ? (Clarté)
    \item [ ] Titre ? (Souvent non demandé, mais si consigne « Donnez un titre », le faire).
\end{itemize}
\end{aretenir}

\coursec{Grammaire outil 1 : les types de phrases}

Dans la section précédente, nous avons défini le résumé : un texte bref, fidèle et neutre, qui restitue les idées essentielles d'un texte sans les déformer. Mais pour réduire un texte, encore faut-il savoir en manipuler les phrases. Or toutes les phrases ne se ressemblent pas : certaines affirment, d'autres questionnent, ordonnent ou expriment une émotion. Reconnaître ces formes, c'est posséder un premier outil de grammaire indispensable : c'est ce que nous apprenons ici, avant de voir, dans la section suivante, comment les phrases sont construites de l'intérieur.

\courssub{Intuition : quatre attitudes du locuteur}

Imaginez quatre personnages du village d'Ilujinle, dans \emph{Le Lion et la perle} de Wole Soyinka. Lakunle, amoureux, dit : \og Sidi est la plus belle fille du village. \fg{} Baroka, intrigué, demande : \og Qui est cette jeune fille ? \fg{} Sadiku, la messagère, ordonne : \og Viens avec moi, Sidi. \fg{} Et la foule s'écrie : \og Quelle perle magnifique ! \fg{}

Ces quatre phrases parlent toutes de Sidi, mais elles ne font pas la même chose. La première \cle{affirme} un fait, la deuxième \cle{pose une question}, la troisième \cle{donne un ordre}, la quatrième \cle{exprime une émotion}. C'est précisément ce que les grammairiens appellent les \cle{types de phrases} : la forme que prend un énoncé selon l'attitude du locuteur envers ce qu'il dit.

\begin{definition}[Type de phrase]
Le \cle{type de phrase} est la forme de l'énoncé déterminée par l'attitude du locuteur : \textbf{affirmer} (phrase déclarative), \textbf{questionner} (phrase interrogative), \textbf{ordonner} (phrase impérative) ou \textbf{s'émouvoir} (phrase exclamative). On distingue ainsi quatre types : \textbf{déclaratif}, \textbf{interrogatif}, \textbf{impératif}, \textbf{exclamatif}.
\end{definition}

\begin{attention}[Ne pas confondre type et forme]
Le \cle{type de phrase} (déclaratif, interrogatif, impératif, exclamatif) renvoie à l'attitude du locuteur. Il ne faut pas le confondre avec la \cle{forme de la phrase} (affirmative ou négative), qui sera étudiée plus tard. Une phrase peut être déclarative et négative : \og Sidi ne refuse pas le cadeau. \fg{} est de type déclaratif, de forme négative.
\end{attention}

\courssub{Les quatre types de phrases et leurs marques}

\textbf{La phrase déclarative}\par

La phrase \cle{déclarative} énonce un fait, une opinion, un jugement. C'est le type le plus fréquent, celui du récit et de l'argumentation. Ses marques de reconnaissance sont :

\begin{itemize}
\item le \textbf{point} (.) en fin de phrase ;
\item une \textbf{intonation descendante} à l'oral ;
\item l'ordre habituel des mots : sujet puis verbe ;
\item le verbe à un \textbf{mode personnel} (indicatif le plus souvent).
\end{itemize}

\begin{exemplebox}[Phrase déclarative]
\og Lakunle est l'instituteur du village d'Ilujinle. \fg{} — Le locuteur affirme un fait, sans émotion particulière ni question.
\end{exemplebox}

\textbf{La phrase interrogative}\par

La phrase \cle{interrogative} pose une question directe. Ses marques sont :

\begin{itemize}
\item le \textbf{point d'interrogation} (?) ;
\item une \textbf{intonation montante} à l'oral ;
\item souvent un \textbf{mot interrogatif} : pronoms (\emph{qui, que, quoi, lequel}), adjectifs (\emph{quel, quelle}), adverbes (\emph{où, quand, comment, pourquoi, combien}) ;
\item parfois l'\textbf{inversion du sujet} (\og Viens-tu ? \fg) ou la tournure \emph{est-ce que}.
\end{itemize}

\begin{exemplebox}[Phrase interrogative]
\og Pourquoi refuses-tu de payer le prix de la dot, Lakunle ? \fg{} — Sidi interroge directement son prétendant.
\end{exemplebox}

\textbf{La phrase impérative}\par

La phrase \cle{impérative} exprime un ordre, un conseil, une prière ou une défense. Ses marques sont :

\begin{itemize}
\item le \textbf{point} (.) ou le \textbf{point d'exclamation} (!) selon l'intensité ;
\item le verbe au \textbf{mode impératif}, sans sujet exprimé : \emph{viens, écoutez, partons} ;
\item parfois l'infinitif dans les consignes : \og Ne pas toucher. \fg
\end{itemize}

\begin{exemplebox}[Phrase impérative]
\og Écoute-moi, Sidi, et laisse là ton seau d'eau ! \fg{} — Lakunle donne un ordre à Sidi. Le verbe \emph{écoute} est à l'impératif présent, sans sujet.
\end{exemplebox}

\textbf{La phrase exclamative}\par

La phrase \cle{exclamative} exprime une émotion vive : admiration, colère, surprise, joie, douleur. Ses marques sont :

\begin{itemize}
\item le \textbf{point d'exclamation} (!) ;
\item une \textbf{intonation forte} ;
\item souvent des mots exclamatifs : \emph{comme, que, quel, quelle, combien, ce que} (\og Quelle beauté ! \fg, \og Comme il danse bien ! \fg).
\end{itemize}

\begin{exemplebox}[Phrase exclamative]
\og Quel honneur pour notre village ! \fg{} — Les habitants d'Ilujinle expriment leur fierté devant la renommée de Sidi.
\end{exemplebox}

Le tableau suivant synthétise les quatre types, leurs marques, un exemple tiré de l'œuvre et l'effet produit.

\begin{center}
\begin{tabularx}{\linewidth}{|l|X|X|X|}
\hline
\rowcolor{popblueL} \textbf{Type de phrase} & \textbf{Marques} & \textbf{Exemple (l'œuvre)} & \textbf{Effet de sens} \\
\hline
Déclarative & point ; intonation descendante ; sujet + verbe & \og Lakunle est l'instituteur du village. \fg & informer, raconter, argumenter \\
\hline
Interrogative & ? ; mot interrogatif ; inversion ou \emph{est-ce que} & \og Pourquoi refuses-tu la dot ? \fg & dialoguer, douter, provoquer \\
\hline
Impérative & verbe à l'impératif ; pas de sujet ; . ou ! & \og Sidi, épouse-moi ! \fg & ordonner, prier, conseiller \\
\hline
Exclamative & ! ; \emph{quel, comme, que} ; intonation forte & \og Quelle perle magnifique ! \fg & exprimer une émotion, dramatiser \\
\hline
\end{tabularx}
\end{center}

\courssub{Types obligatoires et types facultatifs}

Toutes les formes ne se valent pas. Les grammairiens distinguent les \cle{types obligatoires} et les \cle{types facultatifs}.

\begin{definition}[Types obligatoires et facultatifs]
Les \cle{types obligatoires} sont les formes de base, marquées par la grammaire elle-même : la phrase \textbf{déclarative} (ordre sujet--verbe, point) et la phrase \textbf{interrogative} (mot interrogatif, inversion, point d'interrogation). Les \cle{types facultatifs} sont des nuances que le contexte énonciatif ajoute à une forme de base : une phrase déclarative peut devenir \textbf{impérative} ou \textbf{exclamative} par l'intonation, la ponctuation ou la situation de communication.
\end{definition}

Ainsi, la phrase \og Tu viens avec moi. \fg{} est grammaticalement déclarative. Mais dite sur un ton sec par Sadiku à Sidi, elle vaut un ordre : elle prend une \textbf{valeur impérative}. De même, \og Tu as vu cette photo ! \fg{}, dite avec enthousiasme, prend une \textbf{valeur exclamative}. Le contexte énonciatif (qui parle, à qui, dans quelle situation, avec quel ton) transforme la valeur de l'énoncé sans en changer la structure.

\courssub{Les énoncés détournés : quand la forme trompe}

Un piège classique : une phrase peut avoir la \emph{forme} d'un type et la \emph{valeur} d'un autre. On parle d'\cle{énoncé détourné}.

\begin{exemplebox}[Énoncés détournés]
\begin{itemize}
\item \og Tu vas fermer cette porte. \fg{} : forme déclarative, mais valeur \textbf{impérative} (c'est un ordre).
\item \og Il fait chaud ici ! \fg{} : forme déclarative, mais valeur \textbf{exclamative} (émotion) ou même impérative (sous-entendu : ouvrez la fenêtre).
\item \og Est-ce que je t'ai déjà menti ? \fg{} : forme interrogative, mais le locuteur n'attend pas de réponse : c'est une \textbf{question rhétorique} qui affirme \og je ne t'ai jamais menti \fg.
\end{itemize}
\end{exemplebox}

\begin{attention}[Interrogation indirecte : un faux type interrogatif]
Une phrase contenant une \cle{interrogation indirecte} n'est PAS de type interrogatif. Dans \og Baroka demande si Sidi acceptera sa proposition. \fg{}, la phrase entière est \textbf{déclarative} : elle se termine par un point et affirme un fait (Baroka pose une question). Seule la question rapportée (\emph{si Sidi acceptera}) est interrogative, et elle est subordonnée. Retenez : le type se juge sur la phrase entière et sa ponctuation finale, pas sur une partie de la phrase.
\end{attention}

\courssub{Le rôle des types de phrases dans un texte}

Pourquoi un auteur choisit-il un type plutôt qu'un autre ? Parce que chaque type produit un \cle{effet de sens}.

\begin{itemize}
\item La phrase \textbf{déclarative} construit le récit et l'argumentation : elle expose les faits et les idées avec calme et autorité.
\item La phrase \textbf{interrogative} anime le \textbf{dialogue}, exprime le doute ou provoque l'interlocuteur. Dans \emph{Le Lion et la perle}, les échanges entre Sidi et Lakunle fourmillent de questions qui traduisent le conflit entre tradition et modernité.
\item La phrase \textbf{impérative} marque un rapport de force : celui qui ordonne occupe une position dominante (Baroka envers ses sujets, Sadiku envers Sidi).
\item La phrase \textbf{exclamative} fait entrer l'\textbf{émotion} dans le texte : admiration, moquerie, triomphe. Elle donne au texte théâtral sa vivacité.
\end{itemize}

L'alternance des types crée le rythme d'un texte. Un passage tout en déclaratives est monotone ; un dialogue mêlant questions, ordres et exclamations est vivant et dramatique.

\begin{savaistu}[Le théâtre, art de la parole vive]
\emph{Le Lion et la perle} est une pièce de théâtre : elle est écrite pour être dite et jouée, non pour être lue en silence. C'est pourquoi Wole Soyinka, dramaturge nigérian et prix Nobel de littérature en 1986, y multiplie les questions, les ordres et les exclamations : sur scène, l'intonation des acteurs donne à ces phrases toute leur force. Au Cameroun, la tradition orale (contes, palabres, épopées) repose sur le même principe : la manière de dire compte autant que ce qui est dit.
\end{savaistu}

\courssub{Utilité pour le résumé : neutraliser les énoncés}

Voici l'application directe de cette leçon à notre objectif : le résumé. Un résumé doit être \textbf{neutre} : il rapporte les idées sans reproduire les émotions, les ordres ni les questions du texte original. Il faut donc transformer les phrases interrogatives, impératives et exclamatives en phrases \textbf{déclaratives neutres}.

\begin{methode}[Neutraliser une phrase exclamative ou impérative dans un résumé]
\begin{enumerate}
\item \textbf{Identifier} le type de la phrase originale (ponctuation, mode du verbe, mots exclamatifs).
\item \textbf{Dégager l'idée} que la phrase exprime sous l'émotion ou l'ordre.
\item \textbf{Choisir un verbe déclaratif} qui rend compte de l'attitude du personnage : \emph{admirer, s'étonner, ordonner, demander, prier, refuser}.
\item \textbf{Rédiger} une phrase déclarative au présent de vérité générale, sans point d'exclamation ni d'interrogation.
\end{enumerate}
\end{methode}

\begin{exemplebox}[Transformations pour le résumé]
\begin{itemize}
\item \og Quelle beauté ! \fg{} $\rightarrow$ \og Il admire sa beauté. \fg
\item \og Sidi, épouse-moi ! \fg{} (impérative) $\rightarrow$ \og Lakunle demande à Sidi de l'épouser. \fg
\item \og Crois-tu vraiment que la modernité consiste à refuser la dot ? \fg{} (question rhétorique) $\rightarrow$ \og Sidi reproche à Lakunle de confondre modernité et refus des traditions. \fg
\end{itemize}
\end{exemplebox}

\begin{exoresolu}[Identifier et neutraliser des phrases]
\textbf{Énoncé.} Pour chacune des phrases suivantes, inspirées de \emph{Le Lion et la perle}, indiquez son type, puis transformez-la en phrase déclarative neutre utilisable dans un résumé.
\begin{enumerate}
\item \og Sidi, laisse ce seau et écoute-moi ! \fg
\item \og Quel vieux renard que ce Baroka ! \fg
\item \og Sadiku demande si Sidi viendra au palais. \fg
\end{enumerate}
\tcblower
\textbf{Solution détaillée.}
\begin{enumerate}
\item \textbf{Type : impérative} (verbes \emph{laisse} et \emph{écoute} à l'impératif, sans sujet, point d'exclamation). Résumé : \og Lakunle ordonne à Sidi de poser son seau et de l'écouter. \fg{} On a dégagé l'idée (un ordre) et choisi le verbe déclaratif \emph{ordonner}.
\item \textbf{Type : exclamative} (point d'exclamation, mot exclamatif \emph{quel}). Résumé : \og Le locuteur juge Baroka très rusé. \fg{} L'image du renard est rendue par l'idée de ruse, sans émotion.
\item \textbf{Type : déclarative}, malgré l'interrogation indirecte \emph{si Sidi viendra} (point final, verbe principal \emph{demande} à l'indicatif). La phrase est déjà neutre : elle peut entrer telle quelle dans le résumé.
\end{enumerate}
\end{exoresolu}

\begin{exercice}[Entraînement sur un extrait de \emph{Le Lion et la perle}]
Relevez les phrases suivantes, inspirées de la scène où Sadiku annonce à Sidi la proposition de Baroka. Pour chacune : a) indiquez le type ; b) justifiez par une marque précise ; c) proposez, si nécessaire, une formulation déclarative neutre pour un résumé.
\begin{enumerate}
\item \og Sidi, le Baale te demande en mariage. \fg
\item \og Refuserais-tu l'honneur d'entrer au palais du Lion ? \fg
\item \og Viens célébrer cette nouvelle avec moi ! \fg
\item \og Quelle gloire pour notre famille ! \fg
\item \og Sidi se demande si elle doit accepter. \fg
\end{enumerate}
\end{exercice}

\begin{corrige}[Exercice : types de phrases dans \emph{Le Lion et la perle}]
\begin{enumerate}
\item \textbf{Déclarative} : point final, ordre sujet--verbe, verbe \emph{demande} à l'indicatif. Déjà neutre, utilisable telle quelle dans le résumé.
\item \textbf{Interrogative} : point d'interrogation, inversion du sujet (\emph{refuserais-tu}). C'est une question rhétorique qui vaut un reproche. Résumé : \og Sadiku reproche à Sidi d'hésiter devant l'honneur que lui fait le Baale. \fg
\item \textbf{Impérative} : verbe \emph{viens} à l'impératif, sans sujet, point d'exclamation. Résumé : \og Sadiku invite Sidi à célébrer la nouvelle avec elle. \fg
\item \textbf{Exclamative} : point d'exclamation, mot exclamatif \emph{quelle}. Résumé : \og Sadiku estime que ce mariage honore leur famille. \fg
\item \textbf{Déclarative} : l'interrogation indirecte (\emph{si elle doit accepter}) ne change pas le type de la phrase, qui affirme un fait et se termine par un point. Utilisable telle quelle.
\end{enumerate}
\end{corrige}

\begin{aretenir}[L'essentiel]
\begin{itemize}
\item Il existe \textbf{quatre types de phrases} : déclarative (affirmer), interrogative (questionner), impérative (ordonner), exclamative (s'émouvoir).
\item On les reconnaît à la \textbf{ponctuation}, à l'\textbf{intonation}, aux \textbf{mots interrogatifs ou exclamatifs} et au \textbf{mode du verbe}.
\item Les types \textbf{obligatoires} (déclaratif, interrogatif) sont marqués par la grammaire ; les types \textbf{facultatifs} dépendent du contexte énonciatif.
\item Une phrase peut être \textbf{détournée} : forme déclarative, valeur impérative ou exclamative. L'interrogation indirecte ne rend pas une phrase interrogative.
\item Dans un \textbf{résumé}, toutes les phrases doivent devenir des énoncés \textbf{déclaratifs neutres} : on remplace l'émotion, l'ordre ou la question par un verbe qui rapporte l'attitude du personnage.
\end{itemize}
\end{aretenir}

Maintenant que nous savons reconnaître la forme extérieure des phrases, il nous reste à ouvrir le capot : comment une phrase est-elle construite de l'intérieur ? C'est l'objet de la section suivante, consacrée à la structure de la phrase (simple, composée, complexe) et aux relations entre les mots.

\coursec{Grammaire outil 2 : structure de la phrase et relations lexicales}

Dans la section précédente, nous avons étudié les \cle{types de phrases} (déclarative, interrogative, impérative, exclamative), c'est-à-dire la \emph{fonction} de la phrase : ce qu'elle cherche à faire (informer, questionner, ordonner, exprimer une émotion). Nous abordons maintenant la \cle{structure} de la phrase : comment elle est \emph{fabriquée} de l'intérieur, combien de propositions elle contient et comment ces propositions sont reliées entre elles. Cette compétence est l'outil principal du résumé : pour contracter un texte, il faut savoir démonter une longue phrase, identifier l'idée essentielle et la reformuler en moins de mots. Nous verrons ensuite comment les relations entre les mots (hypéronymie et hyponymie) permettent de condenser les énumérations.

\courssub{La phrase simple : une seule proposition}

\begin{definition}[La phrase simple]
Une \cle{phrase simple} est une phrase qui ne contient qu'\textbf{un seul verbe conjugué} (à un mode personnel : indicatif, subjonctif, impératif, conditionnel). Elle forme donc une seule \cle{proposition}, c'est-à-dire un énoncé organisé autour d'un verbe conjugué.
\end{definition}

L'intuition est simple : le verbe conjugué est le « moteur » de la phrase. Une phrase simple n'a qu'un seul moteur. Attention : elle peut contenir des infinitifs ou des participes, qui ne sont \emph{pas} des verbes conjugués.

\begin{exemplebox}[Phrases simples]
\begin{itemize}
\item « Lakunle court vers le marché. » (un seul verbe conjugué : \emph{court})
\item « Sidi, la perle du village, refusa de payer le prix de la dot. » (un seul verbe conjugué : \emph{refusa} ; la phrase reste simple malgré sa longueur)
\item « Le photographe voulait prendre Sidi en photo. » (\emph{voulait} est conjugué ; \emph{prendre} est un infinitif : la phrase reste simple)
\end{itemize}
\end{exemplebox}

\begin{attention}[Piège fréquent]
La longueur d'une phrase ne dit rien de sa structure. Une phrase de vingt mots peut être simple s'il n'y a qu'un verbe conjugué ; une phrase de huit mots peut être complexe s'il y en a deux. \textbf{Comptez les verbes conjugués, pas les mots.}
\end{attention}

\courssub{La phrase composée : juxtaposition et coordination}

\begin{definition}[La phrase composée]
Une \cle{phrase composée} contient \textbf{plusieurs verbes conjugués} répartis en plusieurs propositions, mais ces propositions sont \emph{indépendantes} les unes des autres : aucune n'est à l'intérieur d'une autre. On distingue deux cas :
\begin{itemize}
\item la \cle{juxtaposition} : les propositions sont séparées par une virgule ou un point-virgule, sans mot de liaison ;
\item la \cle{coordination} : les propositions sont reliées par une \cle{conjonction de coordination} : \emph{mais, ou, et, donc, or, ni, car}.
\end{itemize}
\end{definition}

\begin{exemplebox}[Juxtaposition et coordination]
\begin{itemize}
\item Juxtaposition : « Baroka est vieux, il est rusé, il connaît les coutumes. » (trois propositions séparées par des virgules)
\item Coordination : « Sidi est belle, \textbf{mais} elle est vaniteuse. » / « Lakunle refusa de payer la dot, \textbf{car} il jugeait cette coutume dépassée. »
\end{itemize}
\end{exemplebox}

Pour mémoriser les sept conjonctions de coordination, on utilise souvent la phrase mnémotechnique : « \emph{Mais où est donc Ornicar ?} » (\textbf{m}ais, \textbf{ou}, \textbf{et}, \textbf{d}onc, \textbf{or}, \textbf{ni}, \textbf{car}).

\courssub{La phrase complexe : la subordination}

\begin{definition}[La phrase complexe]
Une \cle{phrase complexe} contient plusieurs propositions dont l'une, la \cle{proposition principale}, commande les autres, appelées \cle{propositions subordonnées}. La subordonnée dépend de la principale : elle joue un rôle \emph{à l'intérieur} de celle-ci (complément, épithète, etc.). On distingue principalement :
\begin{itemize}
\item la \cle{subordonnée relative}, introduite par un pronom relatif (\emph{qui, que, où, dont}) : « Le chef \textbf{qui était rusé} organisa un festin. »
\item la \cle{subordonnée conjonctive} (ou complétive), introduite par une conjonction de subordination (\emph{que, parce que, quand, si, puisque, comme}) : « Lakunle pense \textbf{que les coutumes sont dépassées}. »
\item la \cle{subordonnée participiale}, construite avec un participe présent ou passé : « \textbf{Portant ses livres sous le bras}, Lakunle traversa le village. »
\end{itemize}
\end{definition}

\begin{popfigure}
\begin{center}
\begin{tikzpicture}[
  pprinc/.style={draw=popblue, fill=popblueL, rounded corners=3pt, inner sep=5pt, font=\small\bfseries},
  psub/.style={draw=poporange, fill=poporangeL, rounded corners=3pt, inner sep=4pt, font=\small},
  lbl/.style={font=\footnotesize\itshape, text=popmuted},
  arr/.style={-{Stealth}, thick, poporange}
]
% Proposition principale (découpée en deux parties autour de la subordonnée)
\node[pprinc] (p1a) at (0,0) {Le chef};
\node[psub] (psub) at (3,0) {qui était rusé};
\node[pprinc] (p1b) at (6.5,0) {organisa un festin.};
% Flèche d'encadrement
\draw[arr] (psub.south) to[bend right=30] node[lbl, below=2pt] {épithète de « chef »} (p1a.south);
% Étiquettes
\node[lbl, popblue] at (0,0.65) {P1 (principale)};
\node[lbl, popblue] at (6.5,0.65) {P1 (suite)};
\node[lbl, poporange] at (3,0.65) {P2 (sub. relative)};
% Lignes de séparation visuelles
\draw[popline, dashed] (1.5,-0.8) -- (1.5,0.9);
\draw[popline, dashed] (4.75,-0.8) -- (4.75,0.9);
\end{tikzpicture}
\end{center}
\legende{Découpe d'une phrase complexe : la proposition principale P1 (« Le chef... organisa un festin ») encadre la subordonnée relative P2 (« qui était rusé »), qui complète le nom « chef ».}
\end{popfigure}

\courssub{Transformer une phrase complexe : l'opération clé de la contraction}

Pourquoi cette grammaire est-elle un \emph{outil} ? Parce que résumer, c'est souvent \textbf{démonter} les phrases complexes pour n'en garder que l'ossature, ou \textbf{condenser} les subordonnées en simples groupes de mots. Deux mouvements sont possibles :

\begin{enumerate}
\item \textbf{Complexe $\rightarrow$ simples} : on découpe la phrase complexe en plusieurs phrases simples pour bien voir les idées. « Le chef qui était rusé organisa un festin » devient : « Le chef était rusé. Il organisa un festin. » On compte ainsi les idées du texte.
\item \textbf{Complexe $\rightarrow$ condensée} : on garde la proposition principale et on \emph{réduit} la subordonnée à un groupe nominal ou à un adjectif : « le chef \emph{qui était rusé} » devient « le chef \emph{rusé} ». L'idée est conservée, la phrase raccourcit.
\end{enumerate}

\begin{methode}[Condenser une phrase complexe]
\begin{enumerate}
\item Repérer tous les verbes conjugués et découper la phrase en propositions.
\item Identifier la \textbf{proposition principale} : elle porte l'idée essentielle, on la conserve.
\item Pour chaque subordonnée, demander : apporte-t-elle une idée \emph{essentielle} ou un simple détail ?
\item Si c'est un détail : la supprimer ou la réduire en groupe nominal (« le chef qui était rusé » $\rightarrow$ « le chef rusé » ; « le livre qui fut publié en 1963 » $\rightarrow$ « le livre publié en 1963 »).
\item Si c'est une idée essentielle : la reformuler brièvement, sans la supprimer.
\item Vérifier que la phrase obtenue reste grammaticalement correcte et fidèle au sens.
\end{enumerate}
\end{methode}

\begin{attention}[Coordination : on condense, on ne supprime pas]
Dans une phrase \emph{composée}, les propositions coordonnées expriment souvent des idées \textbf{différentes et également importantes}. On ne peut pas en supprimer une sans trahir le texte. « Sidi est belle, mais elle est vaniteuse » contient deux idées (la beauté, la vanité) : le résumé doit garder les deux, quitte à condenser la formulation (« Sidi, belle mais vaniteuse »). La suppression brutale n'est permise que pour les répétitions et les détails secondaires.
\end{attention}

\courssub{Hypéronymie et hyponymie : les relations entre les mots}

Passons de la structure de la phrase aux relations entre les mots. Observez cette liste : \emph{mangue, orange, ananas, papaye}. Tous ces mots désignent des... fruits. Le mot « fruit » est plus général ; il « contient » le sens de chacun des autres. C'est cette relation que nous allons exploiter.

\begin{definition}[Hypéronyme et hyponyme]
Un \cle{hypéronyme} (ou \emph{mot générique}) est un terme dont le sens \textbf{inclut} celui de plusieurs termes plus spécifiques, appelés \cle{hyponymes}. Ainsi, \emph{fruit} est l'hypéronyme de \emph{mangue, orange, ananas} ; inversement, \emph{mangue} est un hyponyme de \emph{fruit}. La relation est hiérarchique : \emph{animal} $>$ \emph{bovin} $>$ \emph{vache}.
\end{definition}

\begin{propriete}[Condensation par l'hypéronyme]
Toute énumération d'hyponymes peut être condensée par son hypéronyme commun \textbf{sans perte de l'idée essentielle}, à condition que le détail des termes énumérés ne soit pas important pour le sens du texte. C'est l'une des techniques de réduction les plus rentables dans un résumé.
\end{propriete}

\begin{exemplebox}[Un exemple chiffré]
« Il acheta des mangues, des oranges, des ananas et des papayes » (\textbf{11 mots}) $\rightarrow$ « Il acheta des fruits » (\textbf{4 mots}). Gain : 7 mots, soit plus de 60 \% de la phrase, sans perdre l'idée essentielle (un achat de fruits).
\end{exemplebox}

\begin{popfigure}
\begin{center}
\begin{tikzpicture}[
  level distance=1.5cm,
  sibling distance=2.2cm,
  hyp/.style={draw=popblue, fill=popblueL, rounded corners=3pt, inner sep=5pt, font=\small\bfseries},
  hypo/.style={draw=popgreen, fill=popgreenL, rounded corners=3pt, inner sep=4pt, font=\small},
  edge from parent/.style={draw=popline, thick},
  edge from parent path={(\tikzparentnode.south) -- (\tikzchildnode.north)}
]
\node[hyp] {animaux}
  child { node[hyp] {bétail}
    child { node[hypo] {chèvres} }
    child { node[hypo] {moutons} }
    child { node[hypo] {bovins} }
  }
  child { node[hypo] {volailles} }
  child { node[hypo] {poissons} };
\end{tikzpicture}
\end{center}
\legende{Arbre d'hypéronymie (exemple du marché camerounais) : « animaux » est l'hypéronyme de « bétail », « volailles », « poissons » ; « bétail » est lui-même l'hypéronyme de « chèvres », « moutons », « bovins ». Au résumé, « chèvres, moutons et bovins » peut devenir « bétail », voire « animaux ».}
\end{popfigure}

\begin{attention}[Quand l'hypéronyme est interdit]
Si le détail des hyponymes est \emph{signifiant}, la condensation trahit le texte. Dans « Il offrit des mangues à sa fiancée, car c'étaient ses fruits préférés », remplacer « mangues » par « fruits » fait disparaître une information essentielle. On ne condense que lorsque l'énumération n'a qu'une valeur d'ensemble.
\end{attention}

\courssub{Champs lexicaux et réseaux de mots dans \emph{Le Lion et la perle}}

L'hypéronymie s'articule avec la notion de \cle{champ lexical} : l'ensemble des mots qui se rapportent à une même idée. Dans \emph{Le Lion et la perle} de Wole Soyinka, deux grands champs lexicaux s'affrontent, et les repérer aide à résumer la pièce :

\begin{center}
\begin{tabularx}{\linewidth}{|l|X|X|}\hline
\rowcolor{popblueL} \textbf{Champ lexical} & \textbf{Mots spécifiques (hyponymes)} & \textbf{Personnage associé} \\ \hline
La tradition & coutume, dot, ancêtres, danse, chef, village, mariage & Baroka (le Lion) \\ \hline
La modernité & école, progrès, photographie, magazine, route, civilisation & Lakunle (l'instituteur) \\ \hline
\end{tabularx}
\end{center}

Pour résumer un passage, on peut remonter du mot spécifique au mot générique : « Lakunle parle sans cesse d'école, de progrès et de civilisation » se condense en « Lakunle vante la modernité ». C'est exactement le même mécanisme que pour les énumérations de fruits.

\begin{exoresolu}[Condenser une phrase complexe et une énumération]
\textbf{Énoncé.} Résumez la phrase suivante en une phrase simple d'au plus 8 mots, sans trahir le sens : « Baroka, qui est le chef du village d'Ilujinle et qui connaît parfaitement les coutumes, les danses et les traditions de ses ancêtres, réussit à épouser Sidi. »
\tcblower
\textbf{Solution détaillée.}
\begin{enumerate}
\item \emph{Verbes conjugués} : « est », « connaît », « réussit » $\rightarrow$ trois propositions, donc une phrase complexe.
\item \emph{Proposition principale} : « Baroka réussit à épouser Sidi » : c'est l'idée essentielle, on la garde.
\item \emph{Subordonnée 1} : « qui est le chef du village d'Ilujinle » se réduit en apposition : « le chef du village ».
\item \emph{Subordonnée 2} : l'énumération « les coutumes, les danses et les traditions » se condense par l'hypéronyme « la tradition ».
\item \emph{Phrase obtenue} : « Baroka, chef du village et fin connaisseur de la tradition, épousa Sidi. » (11 mots) ou, plus bref : « Le chef traditionnel du village épousa Sidi. » (\textbf{7 mots}).
\end{enumerate}
L'idée essentielle (l'identité de Baroka, son ancrage dans la tradition, son mariage avec Sidi) est entièrement conservée.
\end{exoresolu}

\begin{exercice}[À vous de jouer]
\begin{enumerate}
\item Indiquez si chaque phrase est simple, composée ou complexe, en justifiant par le nombre de verbes conjugués : a) « Sidi posa pour le photographe. » b) « Lakunle parla, mais Sidi ne l'écouta pas. » c) « Le magazine qui arriva de Lagos rendit Sidi célèbre. »
\item Condensez par l'hypéronyme : « Au marché, on vendait des chèvres, des moutons, des bovins et des volailles. »
\item Réduisez la subordonnée : « L'instituteur qui refusait la dot quitta la salle. »
\end{enumerate}
\end{exercice}

\begin{corrige}[Exercice : structure et condensation]
\begin{enumerate}
\item a) Phrase \textbf{simple} : un seul verbe conjugué (« posa »). b) Phrase \textbf{composée} : deux verbes conjugués (« parla », « écouta ») reliés par la coordination « mais ». c) Phrase \textbf{complexe} : deux verbes conjugués (« arriva », « rendit »), avec une subordonnée relative (« qui arriva de Lagos ») dans la principale.
\item « Au marché, on vendait des animaux » (ou « du bétail et des volailles » si l'on veut rester plus précis).
\item « L'instituteur hostile à la dot quitta la salle » ou « L'instituteur, refusant la dot, quitta la salle » : la subordonnée relative est réduite en groupe adjectival ou en proposition participiale, plus courte.
\end{enumerate}
\end{corrige}

\begin{aretenir}[L'essentiel de la section]
\begin{itemize}
\item Phrase \textbf{simple} : un seul verbe conjugué ; phrase \textbf{composée} : propositions juxtaposées ou coordonnées (\emph{mais, ou, et, donc, or, ni, car}) ; phrase \textbf{complexe} : une principale et des subordonnées (relative, conjonctive, participiale).
\item Pour résumer : on garde la proposition principale, on réduit les subordonnées en groupes nominaux, on ne supprime jamais une idée essentielle.
\item L'\textbf{hypéronyme} (mot générique) condense toute énumération d'\textbf{hyponymes} (mots spécifiques) : « mangues, oranges, ananas » $\rightarrow$ « fruits ».
\item Les champs lexicaux (tradition / modernité dans \emph{Le Lion et la perle}) permettent de généraliser et de raccourcir les formulations.
\end{itemize}
\end{aretenir}

\coursec{Les techniques de réduction du texte (Contraction de texte)}

Dans la section précédente, nous avons appris à démonter la phrase : distinguer phrase simple, phrase composée et phrase complexe, repérer les expansions du nom, et classer les mots selon leurs relations d'hypéronymie et d'hyponymie. Ces outils grammaticaux n'étaient pas une fin en soi : ils sont précisément les instruments qui vont nous permettre, maintenant, de \cle{réduire} un texte sans le trahir. Car faire une \cle{contraction de texte} — l'épreuve codifiée du Probatoire et du Baccalauréat technique —, ce n'est pas couper au hasard : c'est opérer des choix techniques précis, mot après mot, phrase après phrase, pour atteindre une longueur imposée. Cette section présente les quatre grandes techniques de réduction — la \cle{suppression}, la \cle{généralisation}, la \cle{reformulation} et la \cle{fusion} — ainsi que le travail de repérage et de hiérarchisation qui doit toujours les précéder.

\courssub{Le travail préalable : repérer et hiérarchiser}

Avant de couper le moindre mot, il faut savoir \emph{ce qui doit survivre}. Une contraction de texte n'est pas un texte raccourci au hasard : c'est la reproduction fidèle, en plus court, des \cle{idées essentielles} de l'auteur. Deux gestes préalables sont donc obligatoires.

\begin{methode}[Repérage préalable du texte]
\begin{enumerate}
\item Lire le texte entier une première fois, sans rien écrire, pour saisir le sens global et le ton.
\item À la deuxième lecture, \textbf{souligner les mots-clés} : les noms qui portent l'information (personnages, concepts, lieux), les verbes d'action principaux.
\item \textbf{Encadrer les connecteurs logiques} (\emph{car, mais, donc, cependant, en effet, d'abord, ensuite}) : ils signalent la charpente du raisonnement.
\item \textbf{Entourer la thèse} (l'idée que l'auteur défend) et numéroter les \textbf{arguments} dans la marge.
\item Vérifier : si l'on ne lit que les éléments soulignés, le sens général du texte doit déjà apparaître.
\end{enumerate}
\end{methode}

Une fois le texte balisé, il faut \cle{hiérarchiser} : toutes les phrases n'ont pas le même poids. On distingue trois niveaux.

\begin{definition}[Les trois niveaux d'information dans un texte]
\begin{itemize}
\item L'\textbf{idée principale} (ou thèse) : ce que l'auteur veut faire admettre. Elle doit \emph{toujours} figurer dans la contraction.
\item L'\textbf{argument} : la raison qui soutient l'idée principale. Il doit figurer dans la contraction, sous forme condensée.
\item L'\textbf{illustration} (exemple, anecdote, comparaison, citation, détail descriptif) : elle rend l'argument concret et vivant, mais elle est \emph{sacrifiable} dans une contraction.
\end{itemize}
\end{definition}

\begin{exemplebox}[Hiérarchiser une phrase issue de \emph{Le Lion et la perle}]
Phrase : « Baroka, le Baale d'Ilujinle, est un homme dangereux : ainsi, lors de la construction du chemin de fer, il a soudoyé l'arpenteur pour détourner la voie et empêcher le progrès d'entrer dans le village. »
\begin{itemize}
\item Idée principale : Baroka est un homme dangereux (il s'oppose au progrès).
\item Argument : il a fait détourner le chemin de fer.
\item Illustration : le détail du soudoiement de l'arpenteur.
\end{itemize}
Contraction possible : « Baroka, chef d'Ilujinle, s'oppose au progrès en détournant le chemin de fer. »
\end{exemplebox}

\courssub{Technique 1 : la suppression}

L'intuition est simple : dans un texte, beaucoup de mots sont de la \emph{décoration}. Ils plaisent, ils précisent, ils font vivre le récit — mais ils ne changent pas l'idée. La \cle{suppression} consiste à éliminer ces éléments accessoires.

\begin{definition}[La suppression]
La suppression est la technique qui consiste à éliminer du texte les éléments qui ne portent pas l'idée essentielle : les \textbf{répétitions}, les \textbf{exemples}, les \textbf{comparaisons}, les \textbf{citations}, les \textbf{détails descriptifs} et les \textbf{expansions du nom} (adjectifs épithètes, appositions, compléments du nom, propositions relatives non indispensables).
\end{definition}

\begin{exemplebox}[Supprimer les expansions du nom]
Original : « Sidi, la jeune et ravissante fille du village, coquette et fière de sa beauté, refusa avec mépris la proposition du vieux chef. » (22 mots)

Réduit : « Sidi refusa la proposition du chef. » (6 mots)

On a supprimé : l'apposition (« la jeune et ravissante fille du village »), les adjectifs épithètes (« coquette, fière, vieux ») et le complément de manière (« avec mépris »). L'action essentielle — le refus — demeure intacte.
\end{exemplebox}

\begin{attention}[Ne pas supprimer l'idée avec l'exemple]
Supprimer un exemple ne doit \textbf{jamais} faire disparaître l'idée qu'il illustrait. Si le texte dit : « Les villageois sont attachés aux traditions : ils célèbrent chaque année la danse du devin et consultent les anciens avant tout mariage », on peut supprimer les deux exemples (la danse, la consultation), mais il faut \emph{conserver} l'idée : « Les villageois sont attachés aux traditions. » Une contraction qui ne garderait que « ils célèbrent la danse du devin » trahirait le texte en remplaçant l'idée par son illustration.
\end{attention}

\courssub{Technique 2 : la généralisation}

La deuxième technique s'appuie directement sur la leçon de lexicologie (Sémantique) : les relations d'\cle{hypéronymie} et d'\cle{hyponymie}. Au lieu d'énumérer plusieurs termes précis (des hyponymes), on les remplace par le terme général qui les englobe (l'hypéronyme).

\begin{definition}[La généralisation]
La généralisation consiste à remplacer :
\begin{itemize}
\item une \textbf{énumération} par un \textbf{hypéronyme} : « des manguiers, des irokos et des baobabs » devient « de grands arbres » ;
\item une \textbf{anecdote} par \textbf{l'idée qu'elle illustre} : le récit détaillé d'une ruse de Baroka devient « Baroka est rusé ».
\end{itemize}
\end{definition}

\begin{exemplebox}[Généraliser une énumération]
Original : « Au marché d'Ilujinle, les femmes vendent des ignames, du manioc, des bananes plantains, des arachides, des piments et des feuilles de gombo. » (24 mots)

Réduit : « Au marché d'Ilujinle, les femmes vendent des produits vivriers. » (9 mots)

L'hypéronyme « produits vivriers » remplace à lui seul les six hyponymes énumérés. L'information essentielle (qui vend, où, quoi en général) est conservée.
\end{exemplebox}

\courssub{Technique 3 : la reformulation}

La \cle{reformulation} est la technique la plus fine : elle consiste à dire \emph{la même chose avec moins de mots}, en condensant les unités grammaticales les unes dans les autres. C'est ici que la grammaire de la phrase (phrase simple, composée, complexe ; expansions) devient un outil de compression.

\begin{propriete}[Les trois niveaux de condensation]
La reformulation suit une échelle de réduction :
\begin{enumerate}
\item une \textbf{phrase} peut se condenser en une \textbf{proposition} ;
\item une \textbf{proposition} peut se condenser en un \textbf{groupe de mots} ;
\item un \textbf{groupe de mots} peut se condenser en un \textbf{mot}.
\end{enumerate}
\end{propriete}

\begin{exemplebox}[Condenser progressivement]
\begin{itemize}
\item Phrase $\rightarrow$ proposition : « Lorsque Sadiku arriva, elle annonça la nouvelle » $\rightarrow$ « Sadiku annonça la nouvelle ».
\item Proposition $\rightarrow$ groupe de mots : « qui était plein de malice » $\rightarrow$ « malicieux ».
\item Groupe de mots $\rightarrow$ mot : « d'une manière rusée » $\rightarrow$ « rusément » ; « le fait de refuser » $\rightarrow$ « le refus ».
\end{itemize}
La nominalisation (transformer un verbe en nom : \emph{refuser} $\rightarrow$ \emph{le refus}, \emph{se marier} $\rightarrow$ \emph{le mariage}) est l'outil le plus puissant de cette technique.
\end{exemplebox}

\begin{exemplebox}[Un exemple chiffré de reformulation]
Original : « Baroka, le chef du village, un homme âgé mais vigoureux, rusé et plein de malice » (16 mots)

Réduit : « Baroka, chef rusé et vigoureux » (5 mots)

Gain : 11 mots économisés, soit une réduction de près de 70 \%, sans perte d'information essentielle : l'apposition « le chef du village » devient « chef » ; « plein de malice » est absorbé par « rusé » ; « âgé mais vigoureux » est réduit à « vigoureux », le trait le plus important pour la suite du récit.
\end{exemplebox}

\courssub{Technique 4 : la fusion}

Il arrive que deux phrases ou deux idées proches occupent chacune une phrase entière, alors qu'un seul énoncé bien construit pourrait les porter toutes les deux. La \cle{fusion} consiste à les regrouper en une seule phrase, à l'aide d'un \cle{connecteur logique} adapté qui rend explicite la relation entre les idées.

\begin{definition}[La fusion]
La fusion consiste à regrouper deux idées proches en une seule phrase, en choisissant le connecteur qui exprime exactement leur relation : la cause (\emph{car, parce que}), la conséquence (\emph{donc, si bien que}), l'opposition (\emph{mais, cependant}), l'addition (\emph{et, de plus}) ou la simultanéité (\emph{pendant que, tout en}).
\end{definition}

\begin{exemplebox}[Fusionner deux phrases]
Original : « Lakunle est l'instituteur du village. Il veut moderniser Ilujinle et abolir la dot. » (15 mots)

Réduit : « Lakunle, instituteur du village, veut moderniser Ilujinle en abolissant la dot. » (12 mots)

Autre exemple : « Sidi refuse d'épouser Baroka. Elle estime qu'il est trop vieux et trop vaniteux. » $\rightarrow$ « Sidi refuse d'épouser Baroka \textbf{car} elle le juge trop vieux et vaniteux. » Le connecteur \emph{car} rend la relation de cause visible : la fusion n'a rien perdu du raisonnement.
\end{exemplebox}

\begin{attention}[Bien choisir le connecteur]
Un mauvais connecteur déforme la pensée de l'auteur. Si le texte oppose deux idées, fusionner avec \emph{et} (addition) au lieu de \emph{mais} (opposition) est une faute de sens. Avant de fusionner, demandez-vous toujours : quelle est la relation logique entre ces deux idées ?
\end{attention}

\courssub{Vue d'ensemble : les quatre techniques en tableau}

\begin{popfigure}
\begin{center}
\begin{tabularx}{\linewidth}{|X|l|X|}
\hline
\rowcolor{popblueL} \textbf{Passage original} & \textbf{Technique} & \textbf{Version réduite} \\
\hline
« Sidi, la belle et fière jeune fille du village, marcha d'un pas lent et majestueux vers le marché. » & Suppression & « Sidi marcha vers le marché. » \\
\hline
« Il cultive des ignames, du manioc, du maïs et des bananes. » & Généralisation & « Il cultive des produits vivriers. » \\
\hline
« Baroka, le chef du village, un homme âgé mais vigoureux, rusé et plein de malice » & Reformulation & « Baroka, chef rusé et vigoureux » \\
\hline
« Lakunle aime Sidi. Il refuse de payer la dot. » & Fusion & « Lakunle aime Sidi mais refuse de payer la dot. » \\
\hline
\end{tabularx}
\end{center}
\legende{Les quatre techniques de réduction appliquées à des passages inspirés de \emph{Le Lion et la perle} : pour chaque passage original, la technique dominante et la version réduite.}
\end{popfigure}

\begin{methode}[Réduire un texte en trois passes]
Pour ne pas se perdre, travaillez toujours en trois lectures de travail successives :
\begin{enumerate}
\item \textbf{Première passe : supprimer l'accessoire.} Rayez répétitions, exemples, comparaisons, citations, détails descriptifs et expansions du nom.
\item \textbf{Deuxième passe : généraliser les énumérations.} Remplacez chaque liste par un hypéronyme et chaque anecdote par l'idée qu'elle illustre.
\item \textbf{Troisième passe : reformuler les tournures lourdes.} Condensez phrases en propositions, propositions en groupes de mots, groupes de mots en mots ; fusionnez les idées proches avec le bon connecteur.
\end{enumerate}
Comptez les mots à la fin de chaque passe pour vérifier que vous approchez de la longueur imposée.
\end{methode}

\courssub{Application guidée : réduire un paragraphe de 120 mots à 30 mots}

Appliquons la méthode complète à un paragraphe descriptif sur la vie du village d'Ilujinle, le cadre de \emph{Le Lion et la perle}.

\begin{exoresolu}[Réduction progressive d'un paragraphe sur Ilujinle]
\textbf{Énoncé.} Réduisez le paragraphe suivant (120 mots) à environ 30 mots, en appliquant les trois passes de la méthode.

« Le village d'Ilujinle s'éveille lentement aux premières lueurs de l'aube, dans une lumière douce et dorée qui baigne les toits de paille. Les femmes, comme chaque matin depuis des générations, se rendent à la rivière pour puiser l'eau nécessaire à la journée. Elles portent sur la tête des calebasses, des jarres de terre cuite et des seaux en aluminium. Les hommes, de leur côté, partent aux champs cultiver des ignames, du manioc, du maïs et des bananes plantains. Les enfants courent en riant vers l'école, où le jeune instituteur Lakunle les attend. Ainsi, comme on le voit, la vie du village suit depuis toujours un rythme régulier, immuable et profondément traditionnel. »

\tcblower

\textbf{Solution détaillée.}

\textbf{Passe 1 — suppression de l'accessoire.} On raye les détails descriptifs (« dans une lumière douce et dorée qui baigne les toits de paille »), les comparaisons et formules de remplissage (« comme chaque matin depuis des générations », « comme on le voit », « de leur côté »), les expansions (« nécessaire à la journée », « en riant », « jeune »). Reste environ 75 mots :

« Le village d'Ilujinle s'éveille à l'aube. Les femmes se rendent à la rivière pour puiser l'eau. Elles portent sur la tête des calebasses, des jarres de terre cuite et des seaux en aluminium. Les hommes partent aux champs cultiver des ignames, du manioc, du maïs et des bananes plantains. Les enfants courent vers l'école, où l'instituteur Lakunle les attend. La vie du village suit un rythme régulier, immuable et profondément traditionnel. »

\textbf{Passe 2 — généralisation des énumérations.} « des calebasses, des jarres de terre cuite et des seaux en aluminium » $\rightarrow$ « des récipients » ; « des ignames, du manioc, du maïs et des bananes plantains » $\rightarrow$ « des produits vivriers ». « régulier, immuable et profondément traditionnel » $\rightarrow$ « immuable et traditionnel ». Reste environ 45 mots :

« Le village d'Ilujinle s'éveille à l'aube. Les femmes se rendent à la rivière puiser l'eau dans des récipients. Les hommes partent aux champs cultiver des produits vivriers. Les enfants vont à l'école, où Lakunle les attend. La vie du village suit un rythme immuable et traditionnel. »

\textbf{Passe 3 — reformulation et fusion.} On condense et on fusionne avec des connecteurs : « Le village s'éveille à l'aube » + les trois activités $\rightarrow$ « À l'aube, les femmes puisent l'eau, les hommes cultivent et les enfants vont à l'école. » Version finale (29 mots) :

« À Ilujinle, la vie suit un rythme immuable et traditionnel : chaque matin, les femmes puisent l'eau, les hommes cultivent les champs et les enfants vont à l'école. »

L'idée principale (la permanence des traditions villageoises) est conservée en tête, et les trois activités essentielles sont résumées en une seule phrase rythmée.
\end{exoresolu}

\begin{popfigure}
\begin{center}
\begin{tikzpicture}
\begin{axis}[
    ybar,
    width=0.85\linewidth,
    height=6cm,
    bar width=0.9cm,
    ymin=0, ymax=140,
    ylabel={Nombre de mots},
    xlabel={Étapes de la réduction},
    symbolic x coords={Original, Passe1, Passe2, Passe3},
    xtick=data,
    nodes near coords,
    every node near coord/.append style={font=\small},
    ymajorgrids=true,
    grid style={popline, dashed},
]
\addplot[fill=popblue, draw=popdark] coordinates {(Original,120) (Passe1,75) (Passe2,45) (Passe3,29)};
\end{axis}
\end{tikzpicture}
\end{center}
\legende{Nombre de mots à chaque étape de la réduction du paragraphe sur Ilujinle : 120 mots au départ, 75 après suppression, 45 après généralisation, 29 après reformulation et fusion, soit une réduction totale de 76 \%.}
\end{popfigure}

\begin{aretenir}[L'essentiel des techniques de réduction]
\begin{itemize}
\item \textbf{Suppression} : éliminer répétitions, exemples, comparaisons, citations, détails et expansions du nom — sans jamais perdre l'idée que l'exemple illustrait.
\item \textbf{Généralisation} : remplacer une énumération par un hypéronyme, une anecdote par son idée.
\item \textbf{Reformulation} : condenser phrase $\rightarrow$ proposition $\rightarrow$ groupe de mots $\rightarrow$ mot ; la nominalisation est l'outil roi.
\item \textbf{Fusion} : regrouper deux idées proches avec le connecteur logique exact.
\item Toujours travailler en \textbf{trois passes} et compter les mots à chaque étape.
\end{itemize}
\end{aretenir}

\begin{savaistu}[La règle des 10 \% au Baccalauréat technique]
Au Baccalauréat technique, le sujet de contraction de texte impose toujours une longueur précise (par exemple « environ 60 mots »). Le correcteur pénalise tout résumé dépassant de plus de 10 \% la longueur imposée : pour 60 mots demandés, une contraction de plus de 66 mots perd des points, même si elle est excellente par ailleurs. Inversement, une contraction beaucoup trop courte est soupçonnée d'avoir omis des idées essentielles. D'où l'habitude à prendre dès maintenant : \textbf{comptez toujours vos mots} et indiquez leur nombre à la fin de votre contraction, entre parenthèses.
\end{savaistu}

\begin{exercice}[À vous de jouer]
Réduisez la phrase suivante (34 mots) à environ 12 mots, en indiquant pour chaque suppression la technique utilisée : « Sadiku, la première épouse de Baroka, une femme âgée, expérimentée et pleine d'expérience, se rendit chez Sidi, la plus belle fille du village d'Ilujinle, afin de lui transmettre la demande en mariage de son époux. »
\end{exercice}

\begin{corrige}[Exercice : réduction de la phrase sur Sadiku]
Version réduite possible (12 mots) : « Sadiku, première épouse de Baroka, annonce à Sidi la demande en mariage. »

Techniques utilisées :
\begin{itemize}
\item \textbf{Suppression} des expansions : « une femme âgée, expérimentée et pleine d'expérience » (en outre répétition : \emph{expérimentée} et \emph{pleine d'expérience} disent deux fois la même chose) ; « la plus belle fille du village d'Ilujinle ».
\item \textbf{Reformulation} : « se rendit chez Sidi afin de lui transmettre » $\rightarrow$ « annonce à Sidi » (verbe unique) ; « la demande en mariage de son époux » $\rightarrow$ « la demande en mariage » (le possessif « son » est implicite).
\end{itemize}
L'information essentielle (qui ? Sadiku. Qualité ? première épouse de Baroka. Action ? annonce. À qui ? Sidi. Objet ? demande en mariage) est intégralement conservée.
\end{corrige}

\coursec{Rédaction, amélioration et évaluation du résumé}

Dans la section précédente, nous avons appris à \cle{réduire} un texte : supprimer les exemples et les répétitions, généraliser les énumérations, condenser les formulations. Mais un texte réduit n'est pas encore un \cle{résumé} : il reste une suite d'idées décousues, souvent notées en vrac au brouillon. Il faut maintenant passer à l'étape décisive : \textbf{rédiger} un texte continu, fluide et fidèle, puis le \textbf{relire} pour l'améliorer, enfin l'\textbf{évaluer} avec le regard d'un correcteur d'examen. C'est cette triple étape — rédaction, relecture, évaluation — qui fait la différence entre une copie moyenne et une excellente copie au Probatoire et au Baccalauréat technique.

\courssub{Rédiger le résumé : un texte continu, lié et impersonnel}

\textbf{De la liste d'idées au texte rédigé.} Après la réduction, vous disposez d'un plan d'idées essentielles. Rédiger, c'est transformer ce plan en \cle{phrases complètes} (sujet, verbe, complément), enchaînées dans l'ordre logique du texte d'origine. Un résumé n'est jamais une liste de mots ni une suite de phrases télégraphiques : c'est un petit texte autonome qu'un lecteur peut comprendre sans avoir lu l'original.

\begin{propriete}[Les trois règles d'or de la rédaction du résumé]
\begin{enumerate}
\item \textbf{Des phrases complètes et liées} : chaque phrase est grammaticalement construite, et les phrases sont reliées par des \cle{connecteurs logiques} qui rendent l'enchaînement des idées visible.
\item \textbf{La troisième personne} : on écrit à la \cle{troisième personne} (il, elle, ils, on), jamais à la première personne. Le résumé ne contient ni \og je \fg{} du résuméiste, ni adresse au lecteur.
\item \textbf{Le temps du texte} : on conserve le temps dominant du texte d'origine ; à défaut, on utilise le \cle{présent de vérité générale}, qui convient à la plupart des textes argumentatifs.
\end{enumerate}
\end{propriete}

\textbf{Les connecteurs : la charpente du résumé.} Les connecteurs indiquent la relation logique entre les idées. Ils doivent être choisis en fonction de la structure du texte d'origine :

\begin{center}
\begin{tabularx}{\linewidth}{|l|X|}\hline
\rowcolor{popblueL} \textbf{Relation logique} & \textbf{Connecteurs utiles} \\ \hline
Ordre, énumération & d'abord, ensuite, puis, enfin \\ \hline
Opposition & cependant, pourtant, mais, en revanche \\ \hline
Cause & car, parce que, en effet \\ \hline
Conséquence & c'est pourquoi, ainsi, donc, par conséquent \\ \hline
Ajout & de plus, en outre, également \\ \hline
Conclusion & finalement, en somme, au total \\ \hline
\end{tabularx}
\end{center}

\begin{exemplebox}[Passage du plan au texte rédigé]
Plan après réduction : \og modernité apporte écoles, hôpitaux, routes / mais menace coutumes / il faut concilier les deux \fg{}.

Résumé rédigé : \og La modernité apporte au village des écoles, des hôpitaux et des routes qui améliorent la vie quotidienne. Cependant, elle menace les coutumes et la cohésion traditionnelles. C'est pourquoi l'auteur invite à concilier progrès et traditions. \fg{}

Chaque idée du plan devient une phrase complète ; les connecteurs \og cependant \fg{} et \og c'est pourquoi \fg{} restituent la logique du texte (constat, opposition, conclusion).
\end{exemplebox}

\begin{attention}[Erreur fréquente : la formule d'introduction lourde]
Beaucoup d'élèves commencent leur résumé par : \og Dans ce texte, l'auteur dit que... \fg{} ou \og Ce texte parle de... \fg{}. Ces formules sont \textbf{lourdes et inutiles} : elles consomment des mots sans transmettre d'idée. Le résumé doit \textbf{entrer directement dans les idées}, comme si le texte parlait lui-même. Écrivez : \og La modernité transforme le village... \fg{} et non \og L'auteur explique que la modernité transforme le village... \fg{}.
\end{attention}

\begin{attention}[La fidélité : ne rien ajouter, ne rien juger]
Le résumé est un exercice de \cle{fidélité absolue}. Sont interdits :
\begin{itemize}
\item toute \textbf{idée ajoutée} qui ne figure pas dans le texte, même si elle est vraie ;
\item tout \textbf{jugement personnel} (\og l'auteur a raison \fg{}, \og cette idée est fausse \fg{}, \og j'aime ce texte \fg{}) ;
\item toute \textbf{critique} ou prise de position : le résuméiste est transparent, il laisse parler l'auteur.
\end{itemize}
Ajouter ou juger, c'est sortir de l'exercice : le correcteur sanctionne lourdement ces écarts.
\end{attention}

\courssub{La relecture-amélioration : quatre points de contrôle}

Un résumé ne se rend jamais au premier jet. La \cle{relecture} est une étape obligatoire qui transforme un brouillon en copie d'examen. Elle porte sur quatre points précis, dans un ordre logique : d'abord le fond (longueur, fidélité), ensuite la forme (liaisons, langue).

\begin{methode}[Relire son résumé en 4 points]
\begin{enumerate}
\item \textbf{La longueur} : comptez les mots. Le nombre obtenu respecte-t-il la consigne (par exemple 60 mots, avec la marge de tolérance admise, généralement $\pm 10\,\%$) ? Si le résumé est trop long, supprimez encore des détails ; s'il est trop court, vérifiez qu'aucune idée essentielle n'a été oubliée.
\item \textbf{La fidélité} : relisez le texte d'origine et comparez. Chaque phrase de votre résumé correspond-elle à une idée du texte ? N'avez-vous rien ajouté, rien déformé, aucun jugement personnel ?
\item \textbf{Les liaisons} : les phrases sont-elles reliées par des connecteurs logiques corrects ? L'enchaînement des idées suit-il l'ordre du texte ? Lisez à voix haute : si vous devez vous arrêter pour comprendre, il manque une liaison.
\item \textbf{La langue} : vérifiez l'orthographe, les accords (sujet-verbe, adjectifs), la ponctuation et la grammaire. Variez les \cle{types de phrases} avec discernement (déclaratives surtout) et contrôlez les \cle{structures} : une phrase complexe mal construite (subordonnée sans verbe, relative mal introduite) coûte des points.
\end{enumerate}
\end{methode}

\begin{exemplebox}[Amélioration d'un résumé au brouillon]
Premier jet (45 mots, fautes soulignées par le correcteur) : \og La modernité apporte des écoles et des hôpitaux au village. Les coutumes sont menacés. L'auteur dit qu'il faut concilier tradition et progrès, ce qui est vrai. \fg{}

Problèmes : faute d'accord (\og menacés \fg{}), formule lourde (\og L'auteur dit que \fg{}), jugement personnel (\og ce qui est vrai \fg{}), absence de connecteur d'opposition.

Version améliorée : \og La modernité apporte des écoles et des hôpitaux au village. Cependant, les coutumes sont menacées. C'est pourquoi il faut concilier tradition et progrès. \fg{} (25 mots, fidèle, lié, correct.)
\end{exemplebox}

\courssub{La démarche complète : du texte au résumé évalué}

La figure suivante synthétise l'ensemble de la démarche de contraction de texte, de la première lecture à la relecture finale. C'est la feuille de route à mémoriser pour le jour de l'examen.

\begin{popfigure}
\begin{center}
\begin{tikzpicture}[
node distance=0.55cm,
etape/.style={rectangle, rounded corners=4pt, draw=popblue, fill=popblueL, thick, text width=3.1cm, align=center, minimum height=1.1cm, font=\small},
detail/.style={rectangle, draw=popmuted, fill=popcream, text width=3.1cm, align=center, font=\scriptsize, minimum height=0.9cm},
fleche/.style={-{Stealth[length=3mm]}, thick, popdark}]
\node[etape, align=center] (lire){\textbf{1. LIRE}\\ le texte 2 fois};
\node[detail, below=of lire] (dlire) {saisir le thème et la thèse};
\node[etape, below=of dlire, align=center] (reperer){\textbf{2. REPÉRER}\\ les idées essentielles};
\node[detail, below=of reperer] (dreperer) {souligner, numéroter, hiérarchiser};
\node[etape, below=of dreperer, align=center] (reduire){\textbf{3. RÉDUIRE}\\ le texte};
\node[detail, below=of reduire] (dreduire) {supprimer, généraliser, condenser};
\node[etape, below=of dreduire, align=center] (rediger){\textbf{4. RÉDIGER}\\ le résumé};
\node[detail, below=of rediger] (drediger) {phrases liées, 3\textsuperscript{e} personne, fidélité};
\node[etape, below=of drediger, draw=poporange, fill=poporangeL, align=center] (relire){\textbf{5. RELIRE}\\ et améliorer};
\node[detail, below=of relire] (drelire) {longueur ? fidélité ? liaisons ? langue ?};
\draw[fleche] (lire) -- (dlire);
\draw[fleche] (dlire) -- (reperer);
\draw[fleche] (reperer) -- (dreperer);
\draw[fleche] (dreperer) -- (reduire);
\draw[fleche] (reduire) -- (dreduire);
\draw[fleche] (dreduire) -- (rediger);
\draw[fleche] (rediger) -- (drediger);
\draw[fleche] (drediger) -- (relire);
\draw[fleche] (relire) -- (drelire);
\draw[fleche, poporange, dashed] (drelire.east) -- ++(1.2,0) |- (rediger.east) node[midway, right, font=\scriptsize, align=left, text=popdark]{si un point\\ est insuffisant,\\ retour à la\\ rédaction};
\end{tikzpicture}
\end{center}
\legende{Organigramme de la démarche complète du résumé : cinq étapes enchaînées. La relecture (étape 5) peut renvoyer à la rédaction (étape 4) si un des quatre points de contrôle n'est pas satisfait.}
\end{popfigure}

\courssub{Évaluer un résumé : la grille critériée de l'examen}

Pour progresser, il faut savoir \textbf{comment le correcteur note}. Au Probatoire et au Baccalauréat technique, le résumé est évalué selon des critères précis. Voici une grille type, sur 20 points, qui sert aussi d'outil d'\cle{auto-évaluation} : avant de rendre votre copie, notez-vous vous-même.

\begin{popfigure}
\begin{center}
\begin{tikzpicture}
\node[anchor=north west, font=\small, align=center] at (0,0){
\begin{tabular}{|p{2.6cm}|p{1.1cm}|p{6.2cm}|}
\hline
\rowcolor{popblueL} \textbf{Critère} & \textbf{Points} & \textbf{Indicateurs de réussite} \\
\hline
Fidélité au texte & /4 & Toutes les idées essentielles sont présentes ; aucune idée ajoutée, aucune déformation, aucun jugement personnel. \\
\hline
Respect de la longueur & /4 & Le nombre de mots est indiqué et respecte la consigne (marge de $\pm 10\,\%$ tolérée). \\
\hline
Qualité de la langue & /6 & Orthographe et accords corrects ; phrases bien construites (types et structures maîtrisés) ; ponctuation exacte ; vocabulaire précis. \\
\hline
Cohérence et liaisons & /6 & Ordre logique respecté ; connecteurs variés et pertinents ; texte fluide, compréhensible sans l'original. \\
\hline
\rowcolor{popgoldL} \textbf{Total} & \textbf{/20} & \\
\hline
\end{tabular}
};
\end{tikzpicture}
\end{center}
\legende{Grille d'évaluation critériée du résumé (barème type examen, sur 20 points). Utilisez-la en auto-évaluation : un résumé fidèle mais mal écrit perd jusqu'à 6 points ; un résumé bien écrit mais infidèle ou hors longueur en perd jusqu'à 8.}
\end{popfigure}

\begin{aretenir}[Ce que la grille nous apprend]
La langue et la cohérence pèsent à elles seules \textbf{12 points sur 20} : il ne suffit donc pas d'avoir compris le texte, il faut aussi bien écrire. Inversement, la fidélité et la longueur (8 points) se perdent très vite par négligence : une seule idée ajoutée ou un oubli du nombre de mots coûte cher. La relecture en 4 points couvre exactement ces quatre critères.
\end{aretenir}

\begin{savaistu}[L'exercice le plus complet de l'épreuve]
La contraction de texte est le seul exercice qui évalue \textbf{à la fois la réception et la production} : vous devez d'abord \emph{comprendre} parfaitement un texte (compétence de lecture), puis \emph{exprimer} son contenu dans une langue correcte et personnelle (compétence d'écriture). C'est pourquoi les correcteurs la considèrent comme l'indicateur le plus fiable du niveau global d'un candidat — et pourquoi elle rapporte des points précieux à qui s'y entraîne régulièrement.
\end{savaistu}

\courssub{Séance de production complète : un résumé guidé}

Mettons en pratique toute la démarche sur un texte argumentatif, comme le jour de l'examen.

\begin{exoresolu}[Résumé d'un texte sur la modernité au village (250 mots $\rightarrow$ 60 mots)]
\textbf{Texte d'origine (extrait représentatif, 250 mots).} \og Le village d'Abong-Mbang s'est profondément transformé en vingt ans. Autrefois, les cases de torchis s'alignaient le long de pistes de terre rouge ; aujourd'hui, des maisons en parpaings bordent une route goudronnée. L'école primaire, qui comptait deux classes, en accueille désormais douze, et un collège a ouvert ses portes en 2015. Le centre de santé, équipé d'une maternité, soigne chaque semaine des dizaines de patients qui n'ont plus à marcher jusqu'à la ville. L'électricité, arrivée récemment, permet aux commerçants de conserver leurs produits et aux élèves d'étudier le soir. Personne ne songerait à renoncer à ces progrès. Pourtant, un regard attentif mesure aussi ce qui se perd. Les veillées autour du feu, où les anciens transmettaient contes et proverbes, ont presque disparu : les jeunes préfèrent les écrans. Certaines cérémonies traditionnelles ne rassemblent plus que les vieillards. La langue locale recule devant le français, et même les plats d'autrefois délaissés témoignent d'un art de vivre qui s'efface. La solidarité de quartier, si forte jadis, s'affaiblit à mesure que chacun s'enferme dans sa clôture. Faut-il choisir entre le progrès et la tradition ? Ce serait une erreur. Le développement n'est pas l'ennemi de la culture : il devient dangereux seulement lorsqu'on l'accueille sans réflexion. Le village peut adopter l'école, l'hôpital et la route tout en préservant ses veillées, ses cérémonies et sa langue. Encore faut-il que les jeunes générations, fières de leur héritage, deviennent les gardiennes d'une modernité qui n'oublie pas ses racines. \fg{}

\textbf{Consigne :} résumez ce texte en 60 mots ($\pm 10\,\%$).
\tcblower
\textbf{Étape 1 — Lire.} Thème : la modernité au village. Thèse de l'auteur : il faut concilier progrès et traditions.

\textbf{Étape 2 — Repérer.} Trois idées essentielles : (1) la modernité apporte des progrès matériels (école, santé, route, électricité) ; (2) mais elle fait disparaître les traditions (veillées, cérémonies, langue, solidarité) ; (3) il faut donc concilier développement et culture.

\textbf{Étape 3 — Réduire.} On supprime les exemples précis (Abong-Mbang, la maternité, 2015), on généralise (\og infrastructures scolaires et sanitaires \fg{} pour école + collège + centre de santé), on condense.

\textbf{Étape 4 — Rédiger.} Résumé proposé (58 mots) :

\og La modernité transforme les villages en apportant écoles, soins, routes et électricité, qui améliorent la vie quotidienne. Cependant, elle fait reculer les traditions : veillées, cérémonies, langue locale et solidarité s'affaiblissent. L'auteur refuse de choisir entre progrès et culture : c'est pourquoi il invite les jeunes à construire une modernité fidèle à ses racines. \fg{} (58 mots)

\textbf{Étape 5 — Relire.} Longueur : 58 mots, consigne respectée. Fidélité : les trois idées du texte, rien d'ajouté, aucun jugement. Liaisons : \og cependant \fg{} (opposition), \og c'est pourquoi \fg{} (conséquence). Langue : accords et ponctuation vérifiés, troisième personne, présent de vérité générale. Auto-évaluation : fidélité 4/4, longueur 4/4, langue 6/6, cohérence 6/6.
\end{exoresolu}

\begin{exemplebox}[Un résumé réussi de 60 mots, commenté]
Reprenons le résumé ci-dessus et commentons ses points forts :
\begin{itemize}
\item \textbf{Entrée directe} : \og La modernité transforme les villages... \fg{} — aucune formule lourde, l'idée principale arrive dès le premier mot.
\item \textbf{Généralisation efficace} : \og écoles, soins, routes et électricité \fg{} condense une quinzaine de mots du texte en quatre termes.
\item \textbf{Connecteurs justes} : \og cependant \fg{} marque l'opposition entre progrès et pertes ; \og c'est pourquoi \fg{} introduit la conclusion de l'auteur.
\item \textbf{Fidélité totale} : la thèse (\og refuser de choisir \fg{}) est celle du texte, formulée autrement ; le résuméiste reste invisible.
\item \textbf{Chute forte} : l'image finale (\og une modernité fidèle à ses racines \fg{}) reprend la métaphore de l'auteur sans la copier mot à mot.
\end{itemize}
\end{exemplebox}

\courssub{Compte rendu collectif et remédiation}

La dernière séance de l'unité est consacrée au \cle{compte rendu collectif} : on apprend autant de ses erreurs que de ses réussites. Le professeur relève au tableau les erreurs les plus fréquentes de la classe, sans nommer personne, puis la classe les corrige ensemble avant une réécriture individuelle.

\textbf{Erreurs fréquemment observées et remédiations :}

\begin{center}
\begin{tabularx}{\linewidth}{|X|X|}\hline
\rowcolor{popblueL} \textbf{Erreur fréquente} & \textbf{Remédiation} \\ \hline
Formule lourde : \og l'auteur dit que... \fg{} & Entrer directement dans les idées ; le sujet de la phrase est l'idée, pas l'auteur. \\ \hline
Copie de longs passages du texte & Reformuler avec ses propres mots ; ne reprendre que les termes techniques indispensables. \\ \hline
Jugement personnel (\og je pense que... \fg{}) & Supprimer toute trace de \og je \fg{} ; rester fidèle et neutre. \\ \hline
Longueur non respectée ou non indiquée & Compter les mots et inscrire le total entre parenthèses à la fin. \\ \hline
Phrases juxtaposées sans connecteurs & Relire à voix haute ; ajouter les connecteurs de la relation logique voulue. \\ \hline
Fautes d'accord et de conjugaison & Appliquer le point 4 de la relecture ; vérifier chaque sujet et son verbe. \\ \hline
Idée essentielle oubliée (souvent la thèse) & Vérifier que le plan réduit contient bien la conclusion de l'auteur. \\ \hline
\end{tabularx}
\end{center}

\begin{methode}[Déroulement de la séance de remédiation]
\begin{enumerate}
\item \textbf{Correction mutualisée} : la classe corrige collectivement deux ou trois résumés anonymes projetés ou lus, en appliquant la grille critériée (fidélité /4, longueur /4, langue /6, cohérence /6).
\item \textbf{Auto-évaluation} : chaque élève note sa propre copie avec la même grille et repère ses deux points faibles prioritaires.
\item \textbf{Réécriture individuelle} : chacun réécrit son résumé en corrigeant ses points faibles, puis échange avec un voisin pour une dernière relecture croisée.
\item \textbf{Bilan} : le professeur fait formuler à voix haute les cinq étapes de la démarche (lire, repérer, réduire, rédiger, relire) et les trois règles d'or (phrases liées, troisième personne, fidélité).
\end{enumerate}
\end{methode}

\begin{exercice}[Entraînement à l'auto-évaluation]
Voici le résumé produit par un élève à partir du texte sur la modernité au village (consigne : 60 mots) :

\og Dans ce texte, l'auteur parle du village d'Abong-Mbang qui a beaucoup changé. Il y a maintenant une route goudronnée, un collège ouvert en 2015 et l'électricité. Moi je trouve que le progrès est une bonne chose. Mais les veillées ont disparu et les jeunes préfèrent les écrans, ce qui est dommage. Il faut préserver les traditions. \fg{} (67 mots)

Notez ce résumé avec la grille critériée (fidélité /4, longueur /4, langue /6, cohérence /6), justifiez chaque note, puis réécrivez-le en 60 mots maximum.
\end{exercice}

\begin{corrige}[Exercice : entraînement à l'auto-évaluation]
\textbf{Notation justifiée :}
\begin{itemize}
\item \textbf{Fidélité : 2/4.} Les idées principales sont présentes, mais deux jugements personnels (\og Moi je trouve que le progrès est une bonne chose \fg{}, \og ce qui est dommage \fg{}) violent la règle de neutralité.
\item \textbf{Respect de la longueur : 3/4.} 67 mots : hors de la marge de tolérance (54 à 66 mots) d'un mot seulement, mais le total n'est pas annoncé correctement et le dépassement vient de phrases inutiles.
\item \textbf{Qualité de la langue : 4/6.} Syntaxe correcte, mais la formule \og Dans ce texte, l'auteur parle de \fg{} est lourde et le \og Moi je \fg{} est une faute de registre dans un résumé.
\item \textbf{Cohérence : 4/6.} L'ordre est respecté et \og mais \fg{} marque l'opposition ; cependant les détails précis (2015, route goudronnée) encombrent le propos et la conclusion arrive abruptement.
\end{itemize}
\textbf{Total : 13/20.}

\textbf{Réécriture (57 mots) :} \og Le village s'est modernisé : routes, écoles, santé et électricité ont transformé la vie quotidienne. Cependant, ce progrès s'accompagne de pertes : les veillées disparaissent et les traditions reculent. C'est pourquoi il faut concilier développement et coutumes, afin que la modernité n'efface pas les racines culturelles du village. \fg{} (57 mots)
\end{corrige}

\begin{aretenir}[L'essentiel de la section]
\begin{itemize}
\item Rédiger, c'est transformer le plan réduit en \cle{phrases complètes et liées}, à la \cle{troisième personne}, au temps du texte, sans \og je \fg{}.
\item La \cle{fidélité} est absolue : aucune idée ajoutée, aucun jugement personnel, aucune formule lourde d'introduction.
\item La \cle{relecture} se fait en 4 points : longueur, fidélité, liaisons, langue.
\item La grille d'évaluation (fidélité /4, longueur /4, langue /6, cohérence /6) sert d'outil d'auto-évaluation avant de rendre la copie.
\item La remédiation collective, suivie d'une réécriture individuelle, transforme les erreurs en progrès durables.
\end{itemize}
\end{aretenir}

\newpage

\coursec{Activité d'intégration}

\courssub{Objectifs de l'activité}

Cette activité d'intégration clôt la leçon sur le résumé de texte. Elle invite l'élève à mobiliser, dans une situation authentique de communication scolaire, l'ensemble des savoirs et savoir-faire travaillés dans l'unité :

\begin{itemize}
\item repérer les \cle{idées essentielles} d'un texte informatif et hiérarchiser les informations ;
\item appliquer les quatre \cle{techniques de réduction} : suppression, sélection, généralisation (par \cle{hypéronymie}) et reformulation ;
\item reconnaître et utiliser correctement les \cle{types de phrases} (déclarative, interrogative, impérative, exclamative) et les \cle{structures de phrases} (simple, composée, complexe) ;
\item respecter une \cle{contrainte quantitative} précise (75 mots, soit le quart du texte de départ) ;
\item rédiger un résumé fidèle, cohérent et correct, puis le justifier et l'améliorer grâce à une grille critériée ;
\item développer l'expression orale par la justification des choix de réduction devant la classe.
\end{itemize}

\courssub{Situation}

Au lycée technique de Bafoussam, le délégué de la classe de Première technique a participé, avec les autres délégués, à la réunion trimestrielle organisée par le proviseur. À l'issue de cette réunion, il a rédigé le compte rendu suivant, destiné au journal de la classe.

\begin{exemplebox}[Le compte rendu du délégué (300 mots)]
Monsieur le proviseur a ouvert la réunion à dix heures précises, dans la salle polyvalente du lycée, en présence de vingt-quatre délégués et de trois membres de l'administration. Il a d'abord évoqué le problème de la cantine scolaire. En effet, les élèves se plaignent de la qualité des repas servis : le riz est parfois mal cuit, la sauce manque souvent d'assaisonnement et les portions sont jugées insuffisantes par les plus grands. Le proviseur a promis de rencontrer la cuisinière et le fournisseur dès la semaine prochaine afin d'améliorer la situation. Ensuite, la question de l'état des latrines a été longuement débattue. Certains blocs sont bouchés, d'autres n'ont plus de portes, et l'odeur qui s'en dégage incommode les classes voisines. Le proviseur a annoncé qu'une entreprise de vidange interviendrait avant la fin du mois et que les portes seraient réparées par le menuisier de la ville. Par ailleurs, les délégués ont demandé la modification des horaires des cours du soir, qui se terminent actuellement à dix-huit heures, heure à laquelle il fait déjà nuit et où les transporteurs deviennent rares. Après discussion, il a été décidé que les cours du soir finiraient désormais à dix-sept heures trente. Enfin, le délégué de la classe de Terminale F4 a signalé que le laboratoire de physique manque cruellement de matériel : les microscopes sont hors d'usage, les fils électriques sont dénudés et les produits chimiques sont périmés. Le proviseur a répondu qu'une commande de matériel avait été passée auprès d'un fournisseur de Douala et que la livraison était attendue pour la rentrée prochaine. La réunion s'est achevée à midi, dans une atmosphère cordiale et studieuse, par un mot de remerciement adressé à tous les participants.
\end{exemplebox}

Le proviseur, très satisfait de ce compte rendu, le juge toutefois \textbf{trop long pour être affiché} au tableau d'information du lycée. Il exige une \cle{synthèse} de \textbf{75 mots exactement} (à 5 mots près), fidèle aux décisions prises, rédigée dans un style neutre et administratif. C'est ce résumé que chaque groupe de quatre élèves doit produire. Le meilleur résumé collectif sera réellement affiché au tableau de la classe.

\courssub{Consignes}

Par groupes de quatre, vous travaillerez en suivant les étapes ci-dessous. Chaque étape doit laisser une trace écrite sur votre feuille de groupe.

\begin{enumerate}
\item \textbf{Lecture et repérage.} Lisez silencieusement le compte rendu, puis soulignez en rouge les \cle{idées essentielles} (les quatre problèmes et les quatre décisions) et en bleu les informations secondaires (exemples, détails, répétitions, circonstances).
\item \textbf{Analyse grammaticale.} Relevez dans le texte : une phrase simple, une phrase composée et une phrase complexe, et justifiez chaque classement en comptant les verbes conjugués et en nommant le mode de liaison. Constatez ensuite que toutes les phrases sont du type \cle{déclaratif} : expliquez pourquoi ce type domine dans un compte rendu et pourquoi les types interrogatif, impératif et exclamatif en sont absents.
\item \textbf{Repérage lexical.} Trouvez dans le texte deux énumérations qui peuvent être remplacées par un \cle{hypéronyme}. Proposez pour chacune le mot générique qui convient.
\item \textbf{Réduction.} Appliquez les quatre techniques de réduction : supprimez les informations secondaires, généralisez les énumérations, reformulez les phrases longues en phrases simples ou composées courtes.
\item \textbf{Rédaction.} Rédigez le résumé en \textbf{75 mots} (à 5 mots près), en style neutre, sans jugement personnel, sans reprendre les phrases du texte mot à mot.
\item \textbf{Évaluation croisée.} Échangez votre production avec un autre groupe et évaluez-la à l'aide de la grille critériée fournie (voir la boîte « Démarche et corrigé commenté »).
\item \textbf{Justification orale.} Un porte-parole par groupe explique à la classe deux choix de réduction effectués (une suppression, une généralisation) et les justifie.
\end{enumerate}

\courssub{Ressources à mobiliser}

\begin{aretenir}[Les quatre techniques de réduction]
\begin{enumerate}
\item \textbf{Suppression} : éliminer les exemples, les répétitions, les détails, les précisions de lieu, de temps ou de manière non indispensables.
\item \textbf{Sélection} : ne garder que les idées essentielles, celles qui portent le sens du texte (ici : les problèmes et les décisions).
\item \textbf{Généralisation} : remplacer une énumération ou un terme précis par un \cle{hypéronyme} (mot générique). Exemple : « le riz, la sauce, les portions » devient « les repas ».
\item \textbf{Reformulation} : dire autrement, avec ses propres mots, en privilégiant des phrases simples ou composées courtes, sans copier le texte.
\end{enumerate}
\end{aretenir}

\begin{aretenir}[Outils grammaticaux et lexicaux]
\begin{itemize}
\item \textbf{Types de phrases} : déclarative (informer), interrogative (questionner), impérative (ordonner), exclamative (exprimer un sentiment). Un compte rendu utilise presque exclusivement la phrase \cle{déclarative}, parce que sa fonction est d'\emph{informer}.
\item \textbf{Structures de phrases} : phrase \cle{simple} (une seule proposition, un seul verbe conjugué), phrase \cle{composée} (propositions juxtaposées ou coordonnées), phrase \cle{complexe} (proposition principale + proposition subordonnée).
\item \textbf{Hypéronymie} : relation d'un mot générique (\cle{hypéronyme}) à ses termes spécifiques (\cle{hyponymes}). Exemple : « matériel de laboratoire » est l'hypéronyme de « microscopes, fils électriques, produits chimiques ».
\end{itemize}
\end{aretenir}

\begin{attention}[Les interdits du résumé]
Dans un résumé, il est interdit de : donner son avis personnel ; ajouter une idée absente du texte ; recopier de longues phrases du texte ; dépasser ou sous-estimer fortement la limite de mots imposée ; déformer le sens (par exemple transformer une promesse en décision exécutée).
\end{attention}

\courssub{Démarche et corrigé commenté}

\begin{exoresolu}[Le rapport du délégué : résolution complète]
\textbf{Énoncé.} À partir du compte rendu de 300 mots reproduit ci-dessus, produire une synthèse de 75 mots (à 5 mots près) destinée au tableau d'information du lycée, en appliquant les techniques de contraction de texte.

\tcblower

\textbf{Étape 1 : repérage des idées essentielles.}

Le texte est informatif et suit l'ordre chronologique de la réunion. On identifie quatre thèmes, chacun associant un \cle{problème} et une \cle{décision} (ou une promesse) :

\begin{center}
\begin{tabularx}{\linewidth}{|l|X|X|}
\hline
\rowcolor{popblueL} \textbf{Thème} & \textbf{Problème signalé} & \textbf{Décision / promesse} \\
\hline
Cantine & repas de mauvaise qualité, portions insuffisantes & rencontre de la cuisinière et du fournisseur \\
\hline
Latrines & blocs bouchés, portes détruites, odeurs & vidange avant la fin du mois, réparation des portes \\
\hline
Cours du soir & fin à 18 h, nuit et transporteurs rares & fin avancée à 17 h 30 \\
\hline
Laboratoire & matériel hors d'usage ou périmé & commande passée, livraison à la rentrée \\
\hline
\end{tabularx}
\end{center}

Sont secondaires et donc supprimables : l'heure d'ouverture et de clôture, le lieu, le nombre de participants, l'atmosphère « cordiale et studieuse », le mot de remerciement, les exemples précis (riz mal cuit, sauce sans assaisonnement, microscopes, fils dénudés), le nom du fournisseur (Douala), le menuisier « de la ville ».

\textbf{Étape 2 : analyse grammaticale.}

\begin{itemize}
\item Phrase simple : « Monsieur le proviseur a ouvert la réunion à dix heures précises. » (une seule proposition, un seul verbe conjugué : « a ouvert »).
\item Phrase composée : « Certains blocs sont bouchés, d'autres n'ont plus de portes, et l'odeur incommode les classes voisines. » (deux propositions juxtaposées, puis une proposition coordonnée par « et » ; trois verbes conjugués).
\item Phrase complexe : « Le proviseur a annoncé qu'une entreprise de vidange interviendrait avant la fin du mois. » (proposition principale + subordonnée complétive introduite par « que »).
\item Type de phrase : toutes les phrases du texte sont \cle{déclaratives} (elles donnent une information et se terminent par un point). C'est normal : un compte rendu \emph{informe}, il ne questionne pas, n'ordonne pas et n'exprime pas de sentiment ; les types interrogatif, impératif et exclamatif y sont donc absents.
\end{itemize}

\textbf{Étape 3 : généralisation par hypéronymie.}

\begin{itemize}
\item « le riz mal cuit, la sauce sans assaisonnement, les portions insuffisantes » $\rightarrow$ hypéronyme : \textbf{les repas} (ou « la qualité et la quantité des repas »).
\item « les microscopes hors d'usage, les fils dénudés, les produits périmés » $\rightarrow$ hypéronyme : \textbf{le matériel de laboratoire}.
\end{itemize}

\textbf{Étape 4 : proposition de résumé.}

\emph{Lors de la réunion des délégués, quatre questions ont été traitées. La qualité des repas de la cantine sera améliorée après rencontre avec la cuisinière et le fournisseur. Les latrines seront vidangées et réparées avant la fin du mois. Les cours du soir se termineront désormais à dix-sept heures trente. Enfin, le matériel du laboratoire de physique, commandé à Douala, sera livré à la prochaine rentrée.}

Comptage : Lors(1) de(2) la(3) réunion(4) des(5) délégués(6), quatre(7) questions(8) ont(9) été(10) traitées(11). La(12) qualité(13) des(14) repas(15) de(16) la(17) cantine(18) sera(19) améliorée(20) après(21) rencontre(22) avec(23) la(24) cuisinière(25) et(26) le(27) fournisseur(28). Les(29) latrines(30) seront(31) vidangées(32) et(33) réparées(34) avant(35) la(36) fin(37) du(38) mois(39). Les(40) cours(41) du(42) soir(43) se(44) termineront(45) désormais(46) à(47) dix-sept(48) heures(49) trente(50). Enfin(51), le(52) matériel(53) du(54) laboratoire(55) de(56) physique(57), commandé(58) à(59) Douala(60), sera(61) livré(62) à(63) la(64) prochaine(65) rentrée(66). Soit \textbf{66 mots} : il reste une marge ; on peut ajouter « scolaire » après « cantine », « en présence du proviseur » après « délégués », etc., pour atteindre 70--75 mots sans ajouter d'idée nouvelle.

\textbf{Justification des choix de réduction.}

\begin{itemize}
\item \textbf{Suppressions} : les circonstances (heures d'ouverture et de clôture, salle polyvalente, vingt-quatre délégués, atmosphère cordiale) n'informent pas le lecteur sur les décisions ; elles disparaissent.
\item \textbf{Généralisations} : les plaintes détaillées sur les repas deviennent « la qualité des repas » ; les pannes du laboratoire deviennent « le matériel » (hypéronymie).
\item \textbf{Reformulations} : « il a été décidé que les cours du soir finiraient désormais à dix-sept heures trente » devient la phrase simple « Les cours du soir se termineront désormais à dix-sept heures trente », plus brève et de sens identique.
\item \textbf{Fidélité} : aucune promesse n'est présentée comme un fait accompli (« sera améliorée », « seront vidangées », « sera livré ») : le résumé respecte le sens du texte.
\end{itemize}

\textbf{Étape 5 : grille critériée d'évaluation (sur 10).}

\begin{center}
\begin{tabularx}{\linewidth}{|X|c|}
\hline
\rowcolor{popgreenL} \textbf{Critère} & \textbf{Points} \\
\hline
Respect de la limite de mots (75 $\pm$ 5) & 2 \\
\hline
Présence des quatre idées essentielles (problèmes et décisions) & 3 \\
\hline
Fidélité au sens (aucune idée ajoutée ni déformée, aucun avis personnel) & 2 \\
\hline
Reformulation (pas de recopie de phrases du texte) & 1 \\
\hline
Correction de la langue (orthographe, accords, ponctuation, phrases bien construites) & 2 \\
\hline
\textbf{Total} & \textbf{10} \\
\hline
\end{tabularx}
\end{center}

Un résumé qui oublie une décision (par exemple les latrines) perd au minimum 1 point au critère 2 ; un résumé de 110 mots perd les 2 points du critère 1, même s'il est excellent par ailleurs : la contrainte quantitative fait partie de la commande.
\end{exoresolu}

\courssub{Bilan de l'activité}

Cette activité a permis de mettre en évidence plusieurs acquis essentiels de la leçon.

\begin{itemize}
\item \textbf{Le résumé est une réécriture commandée.} Il répond à une consigne précise (75 mots, style neutre, destinataire : les lecteurs du tableau d'affichage). On ne résume jamais « en général » : on résume pour quelqu'un et pour un usage.
\item \textbf{La hiérarchisation est la clé de la réduction.} Dans un texte informatif, les idées essentielles sont celles qui engagent l'action (ici, les décisions du proviseur) ; les détails, exemples et circonstances peuvent disparaître sans perte de sens.
\item \textbf{La grammaire est un outil de contraction.} Transformer une phrase complexe en phrase simple, repérer les verbes pour identifier les propositions, remplacer une énumération par un hypéronyme : autant d'opérations qui gagnent des mots tout en préservant le sens.
\item \textbf{La fidélité prime sur l'élégance.} Un bon résumé ne juge pas, n'ajoute rien et ne transforme pas une promesse en fait accompli.
\item \textbf{L'évaluation critériée rend le progrès visible.} En évaluant la production d'un autre groupe puis en justifiant oralement ses propres choix, l'élève prend conscience que résumer est une compétence qui se contrôle, se discute et s'améliore — compétence directement mobilisable à l'épreuve de contraction de texte du Probatoire et du Baccalauréat technique.
\end{itemize}

\begin{savaistu}[Le résumé, une compétence professionnelle]
Au-delà de l'examen, savoir résumer est indispensable dans la vie professionnelle : compte rendu de réunion dans une entreprise de Bafoussam ou de Douala, note de synthèse pour un chef de service, rapport de chantier pour un technicien, message d'information affiché dans une administration. Celui qui sait dire l'essentiel en peu de mots se fait toujours écouter.
\end{savaistu}

\newpage

\coursec{Méthodes et exercices résolus}

\coursec{Leçon 20 : Rédiger un résumé de texte}

\courssub{Analyser le texte support : \emph{Le Lion et la Perle} (contexte et exposé)}

\begin{methode}[Situer une œuvre théâtrale africaine pour en faciliter la compréhension]
    \begin{enumerate}
        \item \textbf{Identifier les références bibliographiques :} titre, auteur (Wole Soyinka), nationalité (Nigérian), date de publication (1963), genre (comédie/théâtre).
        \item \textbf{Resituer le contexte historique et culturel :} Nigéria des années 1960, période post-indépendance, choc entre tradition yoruba et modernité occidentale.
        \item \textbf{Repérer la structure dramatique :} 3 actes (Matin, Midi, Soir) = unité de temps (une journée) ; lieu unique (Ilujinle) = unité de lieu ; intrigue centrée sur le mariage de Sidi.
        \item \textbf{Dégager les thèmes majeurs :} tradition vs modernité, pouvoir de la parole/image, condition féminine, ruse et intelligence.
        \item \textbf{Construire une fiche de lecture synthétique :} personnages principaux (Baroka, Lakunle, Sidi, Sadiku), schéma actanciel, ton (satirique, ironique).
    \end{enumerate}
\end{methode}

\begin{savaistu}[Wole Soyinka et le théâtre yoruba]
Wole \textbf{Soyinka} (né en 1934) est le premier Africain noir à recevoir le \textbf{prix Nobel de littérature} (1986). Il fonde le théâtre \emph{Orisun} et théorise un théâtre « total » mêlant danse, chant, mime et proverbes yoruba. \emph{Le Lion et la Perle} (1963) est sa pièce la plus jouée ; elle illustre la \cle{théorie de l'appendice} : la tradition n'est pas un musée mais une force vivante qui absorbe la modernité.
\end{savaistu}

\begin{exoresolu}[Fiche de lecture contextuelle : \emph{Le Lion et la Perle}]
    \textbf{Énoncé :} Vous préparez l'exposé oral sur \emph{Le Lion et la Perle}. Rédigez la fiche de lecture contextuelle (contexte, structure, thèmes, personnages) qui servira de base à votre présentation. Limite : 200 mots maximum.
    \tcblower
    \textbf{Solution détaillée :}
    
    \textbf{1. Références et contexte}
    \emph{Le Lion et la Perle} est une comédie de \textbf{Wole Soyinka} (Nigéria, prix Nobel 1986), publiée en 1963. L'action se déroule à \textbf{Ilujinle}, village yoruba fictif, au début des années 1960, juste après l'indépendance du Nigéria (1960). L'œuvre illustre le conflit entre la \textbf{tradition africaine} (coutumes, polygamie, pouvoir du chef) et la \textbf{modernité occidentale} (école, monogamie, technologie).
    
    \textbf{2. Structure dramatique}
    Respect des unités classiques : \textbf{unité de temps} (une journée : Matin, Midi, Soir), \textbf{unité de lieu} (le village et sa brousse), \textbf{unité d'action} (la conquête de Sidi par Baroka). La pièce mélange prose, chants, danses et mimes (théâtre total).
    
    \textbf{3. Personnages principaux}
    \begin{itemize}
        \item \textbf{Baroka (le Lion) :} Bale (chef) rusé, conservateur mais lucide, symbole de la tradition vivante.
        \item \textbf{Lakunle :} Instituteur, « moderniste » caricatural, pédant, rejette la dot et les coutumes.
        \item \textbf{Sidi (la Perle) :} Belle villageoise, vaniteuse, objet du conflit, choisit finalement la tradition.
        \item \textbf{Sadiku :} Première épouse de Baroka, manipulée puis manipulatrice, lien entre tradition et femmes.
    \end{itemize}
    
    \textbf{4. Thèmes majeurs}
    Tradition vs Modernité ; Pouvoir de l'image (photo) vs Parole ; Condition de la femme (objet de convoitise ou actrice ?) ; Ruse et sagesse du « vieux » face à l'arrogance du « jeune ».
    
    \textbf{Conclusion :} Cette fiche montre que l'œuvre n'est pas un simple folklore mais une satire sociale où la tradition (Baroka) triomphe par intelligence de la modernité creuse (Lakunle).
\end{exoresolu}

\courssub{Maîtriser la définition, les objectifs et les contraintes du résumé}

\begin{definition}[Résumé (Contraction de texte)]
Le résumé est un \textbf{texte court} qui restitue \textbf{fidèlement} les \textbf{idées essentielles} d'un texte source, sans interprétation personnelle, ni commentaire, ni exemple. Il respecte une \textbf{contrainte quantitative} (taux de compression) et des \textbf{contraintes qualitatives} (les 4 F).
\end{definition}

\begin{methode}[Définir le résumé et respecter ses contraintes formelles]
    \begin{enumerate}
        \item \textbf{Contrainte quantitative (taux de compression) :} Calculer le nombre de mots du texte source ($N_{source}$). Le résumé doit contenir environ \textbf{10\% à 25\%} de $N_{source}$ (souvent 1/4 ou 1/3 au Probatoire). \textit{Formule : $N_{résumé} \approx N_{source} \div 4$ (ou $\div 3$).}
        \item \textbf{Contraintes qualitatives (les 4 F) :}
            \begin{itemize}
                \item \textbf{Fidélité :} Respect du sens, de l'ordre des idées, de la thèse de l'auteur.
                \item \textbf{Concision (Fusion) :} Élimination du superflu (exemples, détails, répétitions, adjectifs qualificatifs).
                \item \textbf{Fluidité (Cohérence) :} Connecteurs logiques, temps verbaux adaptés (souvent présent de narration), phrases complètes.
                \item \textbf{Forme :} Un seul bloc (pas de paragraphes), pas de citations directes (style indirect), pas de « L'auteur dit que... ».
            \end{itemize}
        \item \textbf{Vérification finale :} Recompter les mots ($\pm 10\%$ de la cible), relire à voix haute pour la syntaxe.
    \end{enumerate}
\end{methode}

\begin{aretenir}[Règles d'or du résumé au Probatoire]
\begin{itemize}
    \item \textbf{Comptage obligatoire :} Source $\rightarrow$ Cible $\rightarrow$ Brouillon $\rightarrow$ Copie finale (3 comptages minimum).
    \item \textbf{Style indirect exclusif :} « Il dit : "Je viens" » $\rightarrow$ \emph{Il dit qu'il vient.}
    \item \textbf{Présent de narration :} Temps de base pour résumer un récit ou un texte argumentatif.
    \item \textbf{Un seul bloc :} Pas de retour à la ligne, pas de tirets, pas de numérotation.
\end{itemize}
\end{aretenir}

\begin{exoresolu}[Calcul de la longueur cible et identification des contraintes]
    \textbf{Énoncé :} Un texte source de 480 mots est proposé au Probatoire. La consigne demande un résumé d'environ \textbf{un quart} du texte.
    \begin{enumerate}
        \item Calculez le nombre de mots cible exact et l'intervalle de tolérance admis ($\pm 10\%$).
        \item Citez 3 erreurs formelles à éviter absolument dans la rédaction finale.
    \end{enumerate}
    \tcblower
    \textbf{Solution détaillée :}
    
    \textbf{1. Calcul quantitatif}
    \begin{itemize}
        \item Nombre de mots source : $N_s = 480$.
        \item Cible (1/4) : $N_c = 480 \div 4 = \textbf{120 mots}$.
        \item Tolérance $\pm 10\%$ : $10\% \times 120 = 12$ mots.
        \item \textbf{Intervalle admis :} de $120 - 12 = \textbf{108 mots}$ à $120 + 12 = \textbf{132 mots}$.
    \end{itemize}
    \textit{Note :} Si la consigne dit « environ 120 mots », viser 115-125 est plus sûr. Si elle dit « maximum 120 mots », ne pas dépasser 120.
    
    \textbf{2. Trois erreurs formelles interdites (sanctionnées au Bac)}
    \begin{enumerate}
        \item \textbf{Le « paragraphe » ou la liste à puces :} Le résumé est \textbf{un seul bloc de texte} continu.
        \item \textbf{Le style direct / Citation :} « Baroka dit : "Je suis le lion". » $\rightarrow$ \textbf{Interdit}. On écrit : \emph{Baroka affirme être le lion.}
        \item \textbf{L'énonciation méta-textuelle :} « L'auteur explique que... », « Le texte montre que... », « Dans le premier paragraphe... ». $\rightarrow$ \textbf{Interdit}. On résume le \textbf{contenu}, pas la structure du texte.
    \end{enumerate}
    
    \textbf{Conclusion :} Pour ce sujet, viser \textbf{115 à 125 mots} en un seul paragraphe, style indirect, sans méta-langage.
\end{exoresolu}

\courssub{Grammaire outil 1 : Exploiter les types de phrases pour la réduction}

\begin{methode}[Identifier et transformer les types de phrases pour alléger la syntaxe]
    \begin{enumerate}
        \item \textbf{Repérer le type de chaque phrase du source :}
            \begin{itemize}
                \item \textbf{Déclarative (affirmation/négation) :} Base de l'information $\rightarrow$ \textit{Garder le noyau (Sujet + Verbe + COD/COI essentiel)}.
                \item \textbf{Interrogative (directe/indirecte, rhétorique) :} Souvent argumentative $\rightarrow$ \textit{Transformer en déclarative indirecte}. Ex: « Pourquoi partir ? » $\rightarrow$ \emph{Il s'interroge sur le départ.}
                \item \textbf{Impérative (ordre, conseil, prière) :} $\rightarrow$ \textit{Transformer en infinitif ou subjonctif (style indirect)}. Ex: « Écoutez-moi ! » $\rightarrow$ \emph{Il demande qu'on l'écoute / Il exige d'être écouté.}
                \item \textbf{Exclamative :} Marque l'émotion $\rightarrow$ \textit{Supprimer ou réduire à un adverbe/verbe}. Ex: « Comme il est beau ! » $\rightarrow$ \emph{Elle admire sa beauté.}
            \end{itemize}
        \item \textbf{Supprimer les phrases « vides » :} Phrases de transition pure (« Or, ... », « En effet, ... »), incises (« il faut le dire »), formules de politesse.
        \item \textbf{Fusionner :} Regrouper plusieurs déclaratives courtes en une phrase complexe (subordination, coordination, participes).
    \end{enumerate}
\end{methode}

\begin{exoresolu}[Transformation de types de phrases en vue du résumé]
    \textbf{Énoncé :} Transformez les phrases suivantes en énoncés déclaratifs concis utilisables dans un résumé (style indirect, présent de narration).
    \begin{enumerate}
        \item « Baroka, écoute-moi bien ! Ne refuse pas ce cadeau ! » (Sadiku)
        \item « Est-ce que tu crois vraiment que je vais te croire ? » (Sidi, ironique)
        \item « Quel homme rusé que ce Baroka ! » (Narrateur / Didascalie)
        \item « Lakunle entre en scène. Il porte un costume usé. Il parle fort. »
    \end{enumerate}
    \tcblower
    \textbf{Solution détaillée :}
    
    \textbf{Analyse et transformation :}
    
\begin{enumerate}
    \item \textbf{Source :} « Baroka, écoute-moi bien ! Ne refuse pas ce cadeau ! » 
    \textit{Types :} 2 phrases \textbf{impératives} (ordre + interdiction/négation).
    \textit{Stratégie :} Fusion + Style indirect (subordonnée infinitive ou \emph{que} + subjonctif).
    \textbf{Résumé :} \emph{Sadiku ordonne à Baroka d'écouter et lui interdit de refuser le cadeau.} (16 mots $\rightarrow$ 14 mots, gain faible mais style unifié).
    
    \item \textbf{Source :} « Est-ce que tu crois vraiment que je vais te croire ? »
    \textit{Type :} Phrase \textbf{interrogative directe} (rhétorique = forte assertion négative).
    \textit{Stratégie :} Transformer en déclarative négative (le sens réel).
    \textbf{Résumé :} \emph{Sidi doute de la crédibilité de son interlocuteur / Sidi refuse de croire Lakunle.} (Suppression de l'ironie formelle, garde du sens).
    
    \item \textbf{Source :} « Quel homme rusé que ce Baroka ! »
    \textit{Type :} Phrase \textbf{exclamative} (jugement de valeur).
    \textit{Stratégie :} Transformer en attribut du sujet (verbe d'état ou verbe d'opinion).
    \textbf{Résumé :} \emph{Baroka apparaît rusé / La ruse caractérise Baroka.}
    
    \item \textbf{Source :} « Lakunle entre en scène. Il porte un costume usé. Il parle fort. »
    \textit{Types :} 3 phrases \textbf{déclaratives} simples, parataxe (juxtaposition).
    \textit{Stratégie :} \textbf{Subordination} (participes présents / relatifs) + \textbf{Généralisation} (costume usé = détail ?).
    \textbf{Résumé :} \emph{Lakunle entre, vêtu d'un costume usé, et s'exprime avec force.} (Ou plus court : \emph{Lakunle fait une entrée remarquée dans un costume usé.})
    
    \textbf{Bilan :} Le passage de l'interrogatif/exclamatif/impératif au déclaratif indirect est la clé de la \textbf{fidélité énonciative} du résumé.
\end{enumerate}
\end{exoresolu}

\courssub{Grammaire outil 2 : Structure de la phrase et relations lexicales (Hyperonymie/Hyponymie)}

\begin{methode}[Utiliser la structure syntaxique et l'hyperonymie pour généraliser et condenser]
    \begin{enumerate}
        \item \textbf{Analyser la structure de la phrase source :}
            \begin{itemize}
                \item \textbf{Phrase simple :} 1 seul verbe conjugué $\rightarrow$ Souvent une info unique $\rightarrow$ Garder le noyau (SV + compléments essentiels).
                \item \textbf{Phrase composée :} Verbes coordonnés (mais, ou, et, donc, or, ni, car) $\rightarrow$ \textit{Garder la relation logique (concession, cause, conséquence) via subordination}. Ex: « Il est vieux \textbf{mais} il est fort » $\rightarrow$ \emph{Malgré son âge, il est fort.}
                \item \textbf{Phrase complexe :} Proposition principale + subordonnées (relative, conjonctive, infinitive, participiale) $\rightarrow$ \textit{Identifier la proposition principale (info majeure) et réduire les subordonnées en participes, GN, ou les supprimer si accessoires}.
            \end{itemize}
        \item \textbf{Appliquer l'Hyperonymie (Généralisation) :} Remplacer une liste d'hyponymes (termes spécifiques) par leur \textbf{hyperonyme} (terme générique).
            \begin{itemize}
                \item \textit{Liste :} « ignames, maïs, riz, haricots » $\rightarrow$ \textbf{Hyperonyme :} \cle{vivres} / \cle{denrées} / \cle{nourriture}.
                \item \textit{Liste :} « père, mère, oncle, tante, cousins » $\rightarrow$ \textbf{Hyperonyme :} \cle{famille} / \cle{proches} / \cle{parenté}.
                \item \textit{Verbes d'action précis :} « courir, sauter, grimper, nager » $\rightarrow$ \textbf{Hyperonyme :} \cle{se déplacer} / \cle{faire de l'exercice}.
            \end{itemize}
        \item \textbf{Attention :} Ne pas généraliser au point de perdre l'information cruciale (ex: « Baroka boit de l'eau, du vin, du lait » $\nrightarrow$ « Baroka boit des liquides » mais $\rightarrow$ « Baroka boit » si le liquide n'importe pas, ou garder « vin » si c'est symbolique).
    \end{enumerate}
\end{methode}

\begin{exoresolu}[Réduction syntaxique et généralisation lexicale]
    \textbf{Énoncé :} Réduisez le passage suivant (62 mots) en une phrase unique (cible $\approx$ 20 mots) en utilisant la subordination et l'hyperonymie.
    
    \textit{« Le village d'Ilujinle se réveille lentement. Les femmes allument les feux. Elles préparent le foufou, la sauce gombo, le riz et les bananes plantains. Les hommes sortent dans la brousse. Ils vérifient les pièges, coupent du bois et ramassent des champignons. Les enfants courent vers l'école en chantant. »}
    \tcblower
    \textbf{Solution détaillée :}
    
    \textbf{Étape 1 : Segmentation et analyse structurelle}
    \begin{itemize}
        \item P1 (Simple) : Le village se réveille. (Cadre spatio-temporel)
        \item P2 (Composée/Complexe) : Femmes $\rightarrow$ allument + préparent [liste alimentaire].
        \item P3 (Composée) : Hommes $\rightarrow$ sortent + vérifient + coupent + ramassent [liste actions].
        \item P4 (Simple) : Enfants $\rightarrow$ courent [circonstance but/manière].
    \end{itemize}
    
    \textbf{Étape 2 : Généralisation lexicale (Hyperonymie)}
    \begin{itemize}
        \item \textit{Foufou, sauce gombo, riz, bananes plantains} $\rightarrow$ \textbf{le repas} / \textbf{la nourriture} / \textbf{le déjeuner}.
        \item \textit{Vérifient pièges, coupent bois, ramassent champignons} $\rightarrow$ \textbf{vaquent à leurs occupations} / \textbf{travaillent à la brousse} / \textbf{effectuent leurs tâches}.
        \item \textit{Femmes, hommes, enfants} $\rightarrow$ \textbf{Les villageois} / \textbf{La population} (si distinction sexes non cruciale ici).
    \end{itemize}
    
    \textbf{Étape 3 : Fusion syntaxique (Subordination / Participes)}
    \textit{Structure cible :} [Sujet global] + [Verbe principal] + [Participes présents pour actions simultanées].
    
    \textbf{Proposition de résumé (21 mots) :}
    \emph{Au lever du jour, les villageois d'Ilujinle s'activent : les femmes préparent le repas, les hommes travaillent à la brousse, les enfants gagnent l'école en chantant.}
    
    \textbf{Vérification :}
    \begin{itemize}
        \item Fidélité : Actions, acteurs, lieu, temps conservés.
        \item Concision : 62 $\rightarrow$ 21 mots (taux 34\%, acceptable pour un exercice de phrase unique).
        \item Fluidité : Deux-points + structure parallèle (GN + V + Compl).
        \item Hyperonymie : Listes remplacées par « le repas », « travaillent à la brousse ».
    \end{itemize}
\end{exoresolu}

\courssub{Appliquer les 4 techniques de réduction (Suppression, Généralisation, Substitution, Fusion)}

\begin{methode}[Maîtriser les quatre opérations de condensation textuelle]
    \begin{enumerate}
        \item \textbf{La Suppression (Élagage) :} Rayez tout ce qui n'est pas l'idée force.
            \begin{itemize}
                \item Exemples, illustrations, chiffres précis (sauf si thèse), citations.
                \item Adjectifs qualificatifs subjectifs / détails descriptifs.
                \item Répétitions, reformulations, connecteurs de pure liaison (« or », « en effet », « par ailleurs »).
                \item Incises, parenthèses, formules d'atténuation (« il semble que », « probablement »).
            \end{itemize}
        \item \textbf{La Généralisation (Hyperonymie) :} Voir fiche précédente. Remplacer énumérations $\rightarrow$ terme générique. Actions précises $\rightarrow$ verbe global.
        \item \textbf{La Substitution (Réécriture) :} Remplacer un segment long par un équivalent court.
            \begin{itemize}
                \item \textbf{Nominalisation :} « Il a décidé de partir » $\rightarrow$ \emph{Son départ / Sa décision de partir}.
                \item \textbf{Verbalisation :} « Il a fait une analyse » $\rightarrow$ \emph{Il a analysé}.
                \item \textbf{Pronominalisation :} Reprendre les antécédents (Baroka $\rightarrow$ il / le chef / le Bale).
                \item \textbf{Style indirect :} Transformer le dialogue en récit.
            \end{itemize}
        \item \textbf{La Fusion (Intégration syntaxique) :} Relier les idées par la subordination (participes, relatifs, conjonctions) au lieu de la coordination/juxtaposition.
            \begin{itemize}
                \item « Baroka est vieux. Il est rusé. » $\rightarrow$ \emph{Vieux, Baroka n'en est pas moins rusé.} / \emph{Malgré son âge, Baroka est rusé.}
                \item « Sidi voit la photo. Elle est fière. » $\rightarrow$ \emph{Voyant sa photo, Sidi est fière.}
            \end{itemize}
        \item \textbf{Ordre d'application recommandé :} 1. Suppression (nettoyage) $\rightarrow$ 2. Généralisation (vocabulaire) $\rightarrow$ 3. Substitution (syntaxe) $\rightarrow$ 4. Fusion (assemblage final).
    \end{enumerate}
\end{methode}

\begin{exoresolu}[Application séquentielle des 4 techniques sur un extrait]
    \textbf{Énoncé :} Appliquez les 4 techniques pas à pas sur ce paragraphe (58 mots) pour obtenir un résumé de $\approx$ 18 mots.
    
    \textit{« Lakunle, l'instituteur du village, porte un vieux costume trop grand pour lui. Il marche de long en large devant l'école. Il récite des poèmes anglais à haute voix. Les villageois le regardent avec moquerie. Ils trouvent son comportement ridicule et incompréhensible. Lakunle, lui, se croit supérieur aux traditions locales. »}
    \tcblower
    \textbf{Solution détaillée :}
    
    \textbf{Texte source (58 mots) :} « Lakunle, l'instituteur du village, porte un vieux costume trop grand pour lui. Il marche de long en large devant l'école. Il récite des poèmes anglais à haute voix. Les villageois le regardent avec moquerie. Ils trouvent son comportement ridicule et incompréhensible. Lakunle, lui, se croit supérieur aux traditions locales. »
    
    \textbf{1. Suppression (Élagage) :} \textit{~35 mots restants}
    On supprime : « l'instituteur du village », « trop grand pour lui », « de long en large », « devant l'école », « anglais », « à haute voix », « avec moquerie », « ridicule et », « lui ».
    $\rightarrow$ \emph{Lakunle porte un costume inadapté, parade, récite des poèmes. Les villageois le regardent, jugent son attitude incompréhensible. Lakunle se croit supérieur aux traditions.}
    
    \textbf{2. Généralisation (Hyperonymie) :} \textit{~25 mots}
    \emph{Lakunle \textbf{affecte des manières occidentales} (remplace : porte costume inadapté, parade, récite poèmes). Les villageois \textbf{le raillent} (remplace : regardent, jugent attitude incompréhensible). Il se croit \textbf{au-dessus des coutumes} (remplace : supérieur aux traditions locales).}
    
    \textbf{3. Substitution (Nominalisation / Style indirect) :} \textit{~20 mots}
    \emph{L'affectation de manières occidentales par Lakunle provoque la raillerie des villageois. L'instituteur se juge pourtant au-dessus des coutumes.}
    
    \textbf{4. Fusion (Subordination) :} \textbf{Cible 18 mots}
    \emph{\textbf{Moqué pour son affectation occidentale, Lakunle se juge pourtant supérieur aux traditions du village.}} (16 mots)
    
    \textbf{Bilan :} 58 $\rightarrow$ 16 mots (Taux 27\%). Information conservée : Acteur, Action (affectation), Réaction (moquerie), État d'esprit (supériorité), Objet (traditions). Syntaxe fluide (participe passé adjectival « Moqué », coordination « pourtant »).
\end{exoresolu}

\courssub{Rédiger, améliorer et auto-évaluer son résumé (Méthode complète)}

\begin{methode}[Procédé complet en 5 étapes pour réussir l'épreuve de contraction]
    \begin{enumerate}
        \item \textbf{ÉTAPE 1 : Lecture active \& Comptage (10-15 min)}
            \begin{itemize}
                \item Lire 2 fois : 1. globale (sujet, thèse, plan), 2. détaillée (surligner idées forces \textbf{[ID]} et mots-clés, rayer superflu).
                \item Compter mots source ($N_s$). Calculer cible ($N_c = N_s / 4$ ou $/3$). Noter $N_{min}$ et $N_{max}$.
            \end{itemize}
        \item \textbf{ÉTAPE 2 : Rédaction du « Brouillon brut » (15-20 min)}
            \begin{itemize}
                \item Réécrire \textbf{uniquement les idées forces} dans l'ordre, au style indirect, présent de narration.
                \item Appliquer les 4 techniques (Suppression $\rightarrow$ Généralisation $\rightarrow$ Substitution $\rightarrow$ Fusion).
                \item Ne pas compter les mots encore. Viser la concision maximale.
            \end{itemize}
        \item \textbf{ÉTAPE 3 : Ajustement quantitatif (5-10 min)}
            \begin{itemize}
                \item Compter mots brouillon ($N_b$).
                \item Si $N_b > N_{max}$ : \textbf{Resserer} (couper adjectifs, fusionner phrases, généraliser plus).
                \item Si $N_b < N_{min}$ : \textbf{Développer légèrement} (rajouter un détail signifiant, un connecteur logique, préciser un agent).
            \end{itemize}
        \item \textbf{ÉTAPE 4 : Copie propre \& Vérification qualitative (10 min)}
            \begin{itemize}
                \item Copier au net (stylo bleu/noir, lisible).
                \item \textbf{Check-list « 4F » :} \textbf{F}idélité (sens ? ordre ?), \textbf{F}ormation (1 bloc ? pas de « je » ? style indirect ?), \textbf{F}luidité (connecteurs ? syntaxe correcte ?), \textbf{F}ormat (mots comptés ? marge ?).
                \item \textbf{Recompter impérativement} sur la copie finale. Écrire le nombre de mots en bas à droite.
            \end{itemize}
        \item \textbf{ÉTAPE 5 : Gestion du temps (Total $\approx$ 45-50 min)} Respecter les créneaux. Ne pas dépasser 50 min pour garder du temps pour les autres exercices.
    \end{enumerate}
\end{methode}

\begin{exoresolu}[Simulation d'épreuve : De la lecture à la copie finale]
    \textbf{Énoncé :} Voici un extrait (112 mots). Produisez le résumé final (cible 1/4 = 28 mots $\pm$ 10\% $\rightarrow$ \textbf{25 à 31 mots}). Montrez les étapes clés (idées forces, brouillon ajusté, copie finale).
    
    \textit{« Baroka, le Bale d'Ilujinle, reçoit Lakunle dans sa chambre. Le chef est en pleine lutte avec un lutteur professionnel. Il gagne facilement malgré son âge. Lakunle, choqué par cette scène barbare, critique la violence. Baroka répond calmement que la force physique garantit l'autorité. Il offre ensuite une boisson à son hôte. Lakunle refuse avec dédain, affirmant ne pas boire d'alcool. Baroka boit seul, satisfait, et demande à Lakunle le but de sa visite. L'instituteur parle alors du projet de route goudronnée. Baroka écoute sans s'engager, feignant l'ignorance. »}
    \tcblower
    \textbf{Solution détaillée :}
    
    \textbf{Étape 1 : Analyse \& Idées Forces (IF) + Comptage}
    $N_s = 112$ mots. Cible $= 28$. Intervalle $[25; 31]$.
    
    \textit{Segmentation :}
    \begin{enumerate}
        \item \textbf{IF1 :} Baroka (Bale) reçoit Lakunle. Lutte + Victoire (force/âge).
        \item \textbf{IF2 :} Lakunle critique (barbarie). Baroka justifie (force = autorité).
        \item \textbf{IF3 :} Boisson : Lakunle refuse (dédain/alcool). Baroka boit seul.
        \item \textbf{IF4 :} Sujet visite : Route goudronnée. Baroka écoute, feint ignorance (ruse).
    \end{enumerate}
    
    \textbf{Étape 2 : Brouillon Brut (Application techniques)}
    \emph{Baroka, le Bale, reçoit Lakunle pendant qu'il gagne une lutte, prouvant que sa force assure son autorité. Lakunle critique cette violence et refuse l'alcool offert. Baroka boit seul, puis écoute sans s'engager la proposition de route goudronnée de l'instituteur, feignant l'ignorance.} (42 mots) $\rightarrow$ Trop long.
    
    \textbf{Étape 3 : Ajustement (Resserrage + Fusion)}
    \begin{itemize}
        \item « pendant qu'il gagne une lutte » $\rightarrow$ \emph{vainquant un lutteur}.
        \item « prouvant que sa force assure son autorité » $\rightarrow$ \emph{affirmant son autorité par la force}.
        \item « Lakunle critique... et refuse » $\rightarrow$ \emph{Lakunle récuse violence et alcool}.
        \item « Baroka boit seul, puis écoute sans s'engager... feignant l'ignorance » $\rightarrow$ \emph{Baroka, ruseur, feint d'ignorer le projet de route}.
    \end{itemize}
    \textbf{Brouillon ajusté :}
    \emph{Baroka, victorieux d'une lutte, affirme son autorité par la force face à Lakunle qui récuse violence et alcool. Le chef, ruseur, feint ensuite d'ignorer le projet de route goudronnée de l'instituteur.} (30 mots). \textbf{Dans l'intervalle [25-31] !}
    
    \textbf{Étape 4 : Copie Finale \& Vérification}
    \begin{center}
    \framebox[\linewidth]{\parbox{0.9\linewidth}{
    \textbf{Résumé :} Baroka, victorieux d'une lutte, affirme son autorité par la force face à Lakunle qui récuse violence et alcool. Le chef, ruseur, feint ensuite d'ignorer le projet de route goudronnée de l'instituteur.
    }}
    \end{center}
    \textbf{Comptage final :} 30 mots. ✓ Fidélité (ordre respecté, thèse Baroka/Lakunle). ✓ Forme (1 bloc, style indirect, présent). ✓ Fluidité (connecteurs logiques).
\end{exoresolu}

\courssub{Exercice intégré : Résumé d'un extrait de \emph{Le Lion et la Perle} (Matin)}

\begin{methode}[Synthèse : Traiter un extrait littéraire au Probatoire]
    \begin{enumerate}
        \item \textbf{Contextualiser l'extrait :} Identifier l'acte, la scène, les personnages présents, l'enjeu immédiat (ex: Acte 1, Matin : Lakunle/Sidi, débat dot/mariage).
        \item \textbf{Segmenter en séquences narratives/argumentatives :} Découper le texte en 3-4 macro-épisodes (ex: 1. Rencontre/Seau d'eau, 2. Débat dot/modernité, 3. Séduction/Photos, 4. Arrivée du messager/Baroka).
        \item \textbf{Extraire le noyau de chaque séquence :} Qui ? Fait quoi ? Pourquoi ? Résultat ? (Supprimer didascalies, chants, mimes, détails gestuels sauf si signifiants).
        \item \textbf{Rédiger en « Présent de narration » :} Standard pour le résumé d'œuvre. \emph{Baroka reçoit... Lakunle propose... Sidi hésite...}
        \item \textbf{Gérer les dialogues :} Transformer en style indirect libre ou rapporté. \emph{« Je ne paierai pas » $\rightarrow$ Lakunle refuse de payer la dot.}
    \end{enumerate}
\end{methode}

\begin{exoresolu}[Résumée d'extrait : Acte 1, Scène 1 (Lakunle et Sidi) - Simulation Bac]
    \textbf{Énoncé :} Résumez le texte suivant (extrait Acte 1, \emph{Le Lion et la Perle}, env. 180 mots) en \textbf{45 mots maximum} ($\pm 10\%$ $\rightarrow$ \textbf{40 à 50 mots}).
    
    \textit{[Extrait simulé : Lakunle aide Sidi à porter son seau d'eau. Il critique la tradition de la dot (« acheter une femme »). Il propose le mariage « chrétien » sans dot. Sidi refuse : « Je ne serai pas la risée du village ». Lakunle jure de la rendre « civilisée ». Sidi s'impatiente, veut son seau. Lakunle parle de machines, de Lagos. Sidi se moque : « Tu parles trop ». Arrivée des filles : elles dansent le « voyage du voyageur ». Lakunle essaie de les chasser. Sidi part voir la danse.]}
    \tcblower
    \textbf{Solution détaillée :}
    
    \textbf{1. Comptage \& Cible}
    Source $\approx$ 180 mots. Cible $= 45$ mots. Intervalle $[40; 50]$.
    
    \textbf{2. Idées Forces (Séquences)}
    \begin{enumerate}
        \item Lakunle aide Sidi $\rightarrow$ Débat dot (Lakunle contre / Sidi pour tradition).
        \item Lakunle veut « civiliser » Sidi $\rightarrow$ Sidi rejette modernité creuse.
        \item Danse « voyageur » $\rightarrow$ Sidi choisit spectacle traditionnel / Lakunle exclu.
    \end{enumerate}
    
    \textbf{3. Brouillon \& Réduction}
    \emph{Lakunle aide Sidi et refuse la dot, proposant un mariage moderne. Sidi rejette cette modernité qui la déshonorerait. Lakunle prétend la civiliser mais Sidi moque son verbiage. Arrivée des danseuses : Sidi rejoint la tradition, Lakunle reste à l'écart.} (38 mots) $\rightarrow$ Un peu court.
    
    \textbf{4. Ajustement Final (44 mots)}
    \emph{Lakunle aide Sidi et lui propose un mariage sans dot, refusant la tradition. Sidi rejette cette modernité humiliante. Tandis que Lakunle vante la civilisation, Sidi se moque de ses paroles. L'arrivée des danseuses exécutant la « danse du voyageur » scinde le couple : Sidi rejoint la fête traditionnelle, laissant Lakunle seul.}
    
    \textbf{5. Vérification « 4F »}
    \begin{itemize}
        \item \textbf{Fidélité :} Conflit central (dot), caractères (Lakunle moderne/Sidi fière/tradition), chute (danse/séparation) présents.
        \item \textbf{Concision :} 44 mots $\in [40; 50]$. ✓
        \item \textbf{Fluidité :} Connecteurs (Tandis que, Laissant). Présent narration. ✓
        \item \textbf{Forme :} 1 bloc. Pas de « Je », pas de citations. ✓
    \end{itemize}
    \textbf{Nombre de mots : 44}
\end{exoresolu}

\begin{attention}[Les 5 erreurs qui coûtent des points au Probatoire (Résumé)]
    \begin{enumerate}
        \item \textbf{Non-respect de la longueur (Hors tolérance) :} Ne pas compter les mots source $\rightarrow$ cible fausse. Ne pas recompter la copie finale $\rightarrow$ 0 point sur la forme / pénalité lourde. \textit{Réflexe :} Compter 3 fois (Source, Brouillon, Copie finale). Écrire le total en bas.
        \item \textbf{Le « Commentaire » déguisé :} Écrire « L'auteur montre que... », « Dans le premier paragraphe... », « Le texte est divisé en trois parties... ». \textbf{Le résumé n'analyse pas, il restitue.} $\rightarrow$ Perte de points « Fidélité/Contenu ».
        \item \textbf{Rupture de l'énonciation (Style direct / « Je ») :} Garder les guillemets (« Baroka dit : "Je suis le Lion" ») ou écrire « Je pense que... ». \textbf{Obligation :} Style indirect (\emph{Baroka affirme être le Lion}), 3ème personne, présent de narration.
        \item \textbf{Perte de la logique / Ordre des idées :} Mélanger la fin et le début, inverser cause/conséquence (ex: dire que Baroka boit \emph{après} avoir feint l'ignorance alors que c'est avant). \textit{Astuce :} Numéroter les IF sur le brouillon (1, 2, 3...) et les suivre.
        \item \textbf{Syntaxe télégraphique / Phrases nominales uniquement :} Rendre une liste de noms : « Baroka. Lutte. Victoire. Lakunle. Refus. Route. » $\rightarrow$ \textbf{Zéro point « Fluidité/Langue »}. \textbf{Exigence :} Phrases verbales complètes, connecteurs logiques (\emph{car, mais, puis, alors que, bien que}).
    \end{enumerate}
\end{attention}

\newpage

\coursec{Exercices}

\courssub{Je vérifie mes connaissances}

\begin{exercice}[QCM : le résumé et ses règles]
Pour chaque question, choisis la bonne réponse.
\begin{enumerate}
\item Un résumé de texte doit contenir :
\begin{enumerate}
\item des phrases copiées mot à mot dans le texte ;
\item les idées essentielles du texte, reformulées avec ses propres mots ;
\item l'opinion personnelle du rédacteur sur le sujet.
\end{enumerate}
\item Le nombre de mots d'un résumé est généralement :
\begin{enumerate}
\item libre, selon l'envie du rédacteur ;
\item égal à celui du texte de départ ;
\item fixé par la consigne, souvent le quart du texte, avec une marge de tolérance.
\end{enumerate}
\item Dans un résumé, les exemples et les répétitions du texte doivent être :
\begin{enumerate}
\item conservés intégralement ;
\item supprimés ou condensés ;
\item développés davantage.
\end{enumerate}
\item Un résumé se rédige :
\begin{enumerate}
\item à la première personne du singulier ;
\item dans un style neutre, sans jugement personnel ;
\item sous forme de dialogue.
\end{enumerate}
\end{enumerate}
\end{exercice}

\begin{exercice}[Vrai ou faux, avec justification]
Dis si chaque affirmation est vraie ou fausse, puis justifie ta réponse en une phrase.
\begin{enumerate}
\item Une phrase complexe contient au moins une proposition subordonnée.
\item Une phrase composée contient obligatoirement une subordonnée relative.
\item « Viens ici ! » est une phrase de type impératif.
\item Dans un résumé, on peut transformer une phrase exclamative en phrase déclarative.
\item L'hyperonyme de « moto » est « véhicule ».
\end{enumerate}
\end{exercice}

\begin{exercice}[Définitions à compléter]
Recopie et complète chaque définition avec le mot qui convient, choisi dans la liste suivante : \emph{résumé, hyperonyme, hyponyme, phrase simple, phrase complexe, contraction}.
\begin{enumerate}
\item Le \ldots\ldots\ldots{} est un texte bref qui rend compte, de manière objective et fidèle, des idées essentielles d'un texte plus long.
\item Un \ldots\ldots\ldots{} est un mot au sens général qui englobe le sens de plusieurs autres mots (exemple : « fruit » pour « mangue », « banane »).
\item Un \ldots\ldots\ldots{} est un mot dont le sens est inclus dans celui d'un terme plus général (exemple : « okok » est un \ldots\ldots\ldots{} de « plat »).
\item Une \ldots\ldots\ldots{} ne contient qu'une seule proposition, c'est-à-dire un seul verbe conjugué.
\item La \ldots\ldots\ldots{} de texte est l'opération qui consiste à réduire un texte tout en conservant son sens.
\end{enumerate}
\end{exercice}

\begin{exercice}[Comptage et respect de la contrainte]
Un texte argumentatif compte 320 mots. La consigne demande un résumé de 80 mots, avec une marge de tolérance de 10 \%.
\begin{enumerate}
\item Calcule le nombre minimal et le nombre maximal de mots autorisés pour ce résumé.
\item Un élève a rédigé un résumé de 71 mots. Respecte-t-il la consigne ? Justifie.
\item Un autre élève a rédigé 90 mots. Respecte-t-il la consigne ? Justifie.
\item Quelle proportion du texte initial représente un résumé de 80 mots ? Exprime le résultat en fraction puis en pourcentage.
\end{enumerate}
\end{exercice}

\courssub{Je m'entraîne}

\begin{exercice}[Identifier les types de phrases]
Indique le type de chacune des phrases suivantes (déclarative, interrogative, exclamative, impérative), puis précise si elle est de forme affirmative ou négative.
\begin{enumerate}
\item Le vieux chef accueillit l'étranger avec méfiance.
\item Pourquoi les villageois refusent-ils la modernité ?
\item Quelle sagesse dans les paroles de l'aîné !
\item Écoutez attentivement le conteur.
\item La perle ne brille pas pour tout le monde.
\end{enumerate}
\end{exercice}

\begin{exercice}[Transformer pour le résumé]
Les phrases suivantes sont tirées d'un récit inspiré de \emph{Le Lion et la perle}. Transforme chaque phrase exclamative ou interrogative en un énoncé déclaratif neutre, adapté à un résumé.
\begin{enumerate}
\item Comme ce village est paisible au lever du soleil !
\item Le chef acceptera-t-il de vendre la terre des ancêtres ?
\item Quelle arrogance chez ce jeune fonctionnaire !
\item Faut-il vraiment abandonner nos traditions pour être moderne ?
\item Ah ! que la décision du conseil des anciens fut sage !
\end{enumerate}
\end{exercice}

\begin{exercice}[Classer et condenser les phrases]
\textbf{A.} Classe les dix phrases suivantes en trois catégories : phrase simple, phrase composée (coordination ou juxtaposition), phrase complexe (subordination).
\begin{enumerate}
\item Le marchand arriva tôt au marché.
\item Le soleil se levait et les vendeuses installaient leurs étals.
\item L'enfant qui portait un panier courait vers l'école.
\item La pluie tomba ; les clients se réfugièrent sous les hangars.
\item Le chef convoqua les anciens parce que la situation était grave.
\item La rivière coulait paisiblement entre les collines.
\item Bien qu'il fût fatigué, le messager poursuivit sa route.
\item Les femmes chantaient, les hommes dansaient, les enfants riaient.
\item Je pense que la tradition mérite respect.
\item Le tam-tam résonna dans la nuit.
\end{enumerate}
\textbf{B.} Condense les phrases 3, 5, 7 et 9 en réduisant la subordonnée (participe, apposition ou groupe nominal), sans changer le sens.
\end{exercice}

\begin{exercice}[Hyperonymie et hyponymie]
\textbf{A.} À partir des listes de mots ci-dessous, construis trois arbres hyperonyme/hyponymes (un terme général au sommet, au moins trois termes spécifiques en dessous).
\begin{itemize}
\item Liste 1 : moto, moyens de transport, taxi, pirogue, bus ;
\item Liste 2 : produits du marché, igname, huile de palme, arachide, manioc ;
\item Liste 3 : membres de la famille, oncle, grand-mère, cousin, tante.
\end{itemize}
\textbf{B.} Condense chacune des phrases suivantes en remplaçant l'énumération par l'hyperonyme qui convient.
\begin{enumerate}
\item Au marché de Mokolo, on vend des ignames, du manioc, des arachides et de l'huile de palme.
\item Pour se rendre à Douala, les voyageurs prennent le bus, le taxi, la moto ou le train.
\item À la fête du village étaient présents mes oncles, mes tantes, mes cousins et ma grand-mère.
\end{enumerate}
\end{exercice}

\courssub{Je me prépare au Bac}

\begin{exercice}[Type Baccalauréat technique : résumé d'un texte argumentatif]
Lis attentivement le texte suivant, puis réponds aux questions.

\textbf{Texte : L'école traditionnelle et l'école moderne au Cameroun}

\emph{Longtemps, dans nos villages, l'éducation des enfants ne passait pas par les salles de classe. Elle se faisait au quotidien, auprès des parents et des anciens. La grand-mère transmettait les contes et les proverbes le soir, autour du feu. Le père apprenait à ses fils la pêche, la chasse ou la culture des champs. La mère enseignait à ses filles la cuisine, le respect des aînés et les chants des cérémonies. Cette école traditionnelle formait des hommes et des femmes utiles à leur communauté, profondément enracinés dans leur culture.}

\emph{Aujourd'hui, l'école moderne occupe une place centrale dans la vie des jeunes Camerounais. Dès l'âge de six ans, l'enfant apprend à lire, à écrire et à compter. Plus tard, il découvre les sciences, les langues étrangères et les techniques qui lui ouvriront les portes d'un métier. Grâce à l'école moderne, un jeune du village peut devenir médecin, ingénieur ou enseignant. Elle offre ainsi des chances de promotion sociale que l'éducation traditionnelle ne pouvait pas donner.}

\emph{Cependant, l'école moderne présente des limites. Elle éloigne parfois les jeunes de leur culture : certains bacheliers ignorent les proverbes de leur langue et méprisent les savoir-faire de leurs parents. De plus, les diplômes ne garantissent plus un emploi, et beaucoup de jeunes diplômés restent sans travail.}

\emph{Il ne faut donc pas opposer les deux écoles, mais les associer. L'école moderne doit préparer à la vie professionnelle, tandis que la famille et la communauté continuent de transmettre les valeurs et la culture. Un jeune Camerounais épanoui est celui qui maîtrise les savoirs modernes sans renier ses racines.} (320 mots)

\begin{enumerate}
\item Résume ce texte en 80 mots (plus ou moins 10 \%), en respectant les règles du résumé : fidélité aux idées, reformulation, neutralité, suppression des exemples. Indique à la fin le nombre exact de mots utilisés.
\item Justifie deux choix de réduction que tu as effectués (exemple supprimé, énumération condensée par un hyperonyme, phrase complexe simplifiée).
\item Releve dans le texte deux phrases complexes et montre comment tu les as condensées dans ton résumé.
\end{enumerate}
\end{exercice}

\begin{exercice}[Type Probatoire : grammaire et résumé à partir de \emph{Le Lion et la perle}]
Voici cinq énoncés inspirés de l'univers de \emph{Le Lion et la perle}.
\begin{enumerate}
\item Le lion règne sur la savane depuis toujours.
\item Qui osera contester la force du lion ?
\item Quelle ruse chez ce petit animal qui détient la perle !
\item Observez comment la sagesse triomphe de la force.
\item Bien que le lion soit puissant, la perle échappe à ses griffes.
\end{enumerate}
\textbf{Questions :}
\begin{enumerate}
\item Identifie le type de chacun des cinq énoncés (déclaratif, interrogatif, exclamatif, impératif).
\item Classe ces cinq phrases en phrases simples, composées ou complexes, en justifiant par le nombre et la nature des propositions.
\item Transforme les phrases exclamatives, interrogatives et impératives en énoncés déclaratifs neutres, propres à figurer dans un résumé.
\item Rédige, à partir de tes transformations, un mini-résumé de trois phrases maximum qui rend compte de l'opposition entre la force et la sagesse dans l'œuvre.
\end{enumerate}
\end{exercice}

\begin{exercice}[Type Baccalauréat technique : remédiation d'un résumé fautif]
Un élève devait résumer en 60 mots un texte de 240 mots sur l'importance des contes dans l'éducation des enfants. Voici sa production :

\emph{« Le conte est très important dans l'éducation des enfants car il transmet les valeurs de la société, comme le courage, la solidarité, le respect des anciens, la patience et l'honnêteté. » (phrase copiée mot à mot dans le texte). Moi je trouve que les contes sont magnifiques et que les parents d'aujourd'hui ont tort de ne plus raconter des histoires à leurs enfants. Le conteur est assis près du feu, la lune brille, les enfants écoutent en silence, quel beau spectacle ! Il faut absolument que tout le monde lise ce texte formidable. En conclusion, le conte éduque et divertit les enfants.} (95 mots)

\textbf{Questions :}
\begin{enumerate}
\item Ce résumé comporte cinq défauts majeurs. Repère-les et cite pour chacun le passage concerné. (Indices : longueur, copie, jugement personnel, description inutile, type de phrase inadapté.)
\item Pour chaque défaut repéré, explique la règle du résumé qui a été enfreinte.
\item Propose une version corrigée du résumé, en 60 mots (plus ou moins 10 \%), qui respecte toutes les règles de la contraction de texte. Indique le nombre de mots à la fin.
\end{enumerate}
\end{exercice}

\newpage

\coursec{Corrigés des exercices}

\begin{corrige}[Exercice 1 — QCM : le résumé et ses règles]
\begin{enumerate}
\item Réponse \textbf{b} : les idées essentielles du texte, reformulées avec ses propres mots. Le résumé est une \cle{reformulation} : copier des phrases (réponse a) est interdit, et y glisser son opinion (réponse c) viole la règle de neutralité.
\item Réponse \textbf{c} : le nombre de mots est fixé par la consigne, souvent le quart du texte, avec une marge de tolérance (généralement 10\,\%). C'est la \cle{contrainte de longueur}, première règle de la contraction de texte.
\item Réponse \textbf{b} : les exemples et les répétitions doivent être supprimés ou condensés. Ce sont des procédés d'illustration et d'insistance qui ne font pas partie des idées essentielles.
\item Réponse \textbf{b} : un résumé se rédige dans un style neutre, sans jugement personnel, à la troisième personne. Le rédacteur doit rester invisible : il rapporte la pensée de l'auteur, pas la sienne.
\end{enumerate}
\end{corrige}

\begin{corrige}[Exercice 2 — Vrai ou faux, avec justification]
\begin{enumerate}
\item \textbf{Vrai.} Par définition, la phrase complexe contient une proposition principale et au moins une proposition subordonnée (relative, conjonctive ou participiale).
\item \textbf{Faux.} La phrase composée ne contient que des propositions \emph{indépendantes}, reliées par coordination (\emph{et, mais, ou}) ou par juxtaposition (virgule, point-virgule) ; la subordonnée relative est le signe de la phrase complexe.
\item \textbf{Vrai.} « Viens ici ! » exprime un ordre avec un verbe à l'impératif : c'est une phrase de type impératif.
\item \textbf{Vrai.} Dans un résumé, on neutralise les énoncés : « Quelle sagesse ! » peut devenir « L'auteur souligne la sagesse des anciens », car le résumé privilégie la phrase déclarative.
\item \textbf{Vrai.} « Véhicule » est le terme général (hyperonyme) qui englobe « moto », terme spécifique (hyponyme), comme « bus » ou « pirogue ».
\end{enumerate}
\end{corrige}

\begin{corrige}[Exercice 3 — Définitions à compléter]
\begin{enumerate}
\item Le \textbf{résumé} est un texte bref qui rend compte, de manière objective et fidèle, des idées essentielles d'un texte plus long.
\item Un \textbf{hyperonyme} est un mot au sens général qui englobe le sens de plusieurs autres mots (exemple : « fruit » pour « mangue », « banane »).
\item Un \textbf{hyponyme} est un mot dont le sens est inclus dans celui d'un terme plus général (exemple : « okok » est un \textbf{hyponyme} de « plat »).
\item Une \textbf{phrase simple} ne contient qu'une seule proposition, c'est-à-dire un seul verbe conjugué.
\item La \textbf{contraction} de texte est l'opération qui consiste à réduire un texte tout en conservant son sens.
\end{enumerate}
\end{corrige}

\begin{corrige}[Exercice 4 — Comptage et respect de la contrainte]
\begin{enumerate}
\item La marge de tolérance est de 10\,\% de 80 mots, soit :
\[ 80 \times \frac{10}{100} = 8 \text{ mots.} \]
Nombre minimal : $80 - 8 = 72$ mots. Nombre maximal : $80 + 8 = 88$ mots. Le résumé doit donc contenir \textbf{entre 72 et 88 mots}.
\item Le résumé de 71 mots \textbf{ne respecte pas} la consigne : $71 < 72$, il est en dessous du minimum autorisé.
\item Le résumé de 90 mots \textbf{ne respecte pas} la consigne : $90 > 88$, il dépasse le maximum autorisé.
\item Proportion du texte initial :
\[ \frac{80}{320} = \frac{1}{4} = 0,25 = 25\,\%. \]
Le résumé représente \textbf{le quart} du texte, soit \textbf{25\,\%}.
\end{enumerate}
\end{corrige}

\begin{attention}[Point de vigilance]
La tolérance se calcule toujours sur le nombre de mots \emph{demandé} (ici 80), jamais sur la longueur du texte de départ. Au Baccalauréat technique, indique le nombre de mots à la fin de ton résumé : c'est une exigence de la consigne.
\end{attention}

\begin{corrige}[Exercice 5 — Identifier les types de phrases]
\begin{enumerate}
\item « Le vieux chef accueillit l'étranger avec méfiance. » : phrase \textbf{déclarative}, de forme \textbf{affirmative} (elle donne une information, ponctuée d'un point).
\item « Pourquoi les villageois refusent-ils la modernité ? » : phrase \textbf{interrogative}, de forme \textbf{affirmative} (point d'interrogation, inversion du sujet).
\item « Quelle sagesse dans les paroles de l'aîné ! » : phrase \textbf{exclamative}, de forme \textbf{affirmative} (point d'exclamation, marqueur « quelle »).
\item « Écoutez attentivement le conteur. » : phrase \textbf{impérative}, de forme \textbf{affirmative} (verbe à l'impératif exprimant un ordre).
\item « La perle ne brille pas pour tout le monde. » : phrase \textbf{déclarative}, de forme \textbf{négative} (présence de la négation « ne\ldots pas »).
\end{enumerate}
\end{corrige}

\begin{corrige}[Exercice 6 — Transformer pour le résumé]
Voici des transformations possibles (d'autres formulations correctes sont acceptables si elles restent neutres et fidèles) :
\begin{enumerate}
\item Comme ce village est paisible au lever du soleil ! $\rightarrow$ \emph{Le village est paisible au lever du soleil.}
\item Le chef acceptera-t-il de vendre la terre des ancêtres ? $\rightarrow$ \emph{Le chef hésite à vendre la terre des ancêtres.} (ou : \emph{On se demande si le chef acceptera de vendre la terre des ancêtres.})
\item Quelle arrogance chez ce jeune fonctionnaire ! $\rightarrow$ \emph{Le jeune fonctionnaire fait preuve d'arrogance.}
\item Faut-il vraiment abandonner nos traditions pour être moderne ? $\rightarrow$ \emph{L'auteur s'interroge sur la nécessité d'abandonner les traditions pour être moderne.}
\item Ah ! que la décision du conseil des anciens fut sage ! $\rightarrow$ \emph{La décision du conseil des anciens fut sage.}
\end{enumerate}
\end{corrige}

\begin{methode}[Rappel de la démarche]
Pour neutraliser un énoncé : (1) supprimer la marque d'exclamation ou d'interrogation ; (2) remplacer la ponctuation par un point ; (3) si nécessaire, introduire un verbe déclaratif (\emph{souligner, s'interroger sur, constater}) qui rapporte l'idée sans la commenter.
\end{methode}

\begin{corrige}[Exercice 7 — Classer et condenser les phrases]
\textbf{A. Classement.} On compte les verbes conjugués et on observe le lien entre les propositions.

\begin{center}
\begin{tabularx}{\linewidth}{|c|X|l|}
\hline
\rowcolor{popblueL} N° & Phrase & Catégorie \\
\hline
1 & Le marchand arriva tôt au marché. & Phrase simple (1 verbe) \\
\hline
2 & Le soleil se levait \textbf{et} les vendeuses installaient leurs étals. & Phrase composée (coordination) \\
\hline
3 & L'enfant \textbf{qui} portait un panier courait vers l'école. & Phrase complexe (subordonnée relative) \\
\hline
4 & La pluie tomba \textbf{;} les clients se réfugièrent sous les hangars. & Phrase composée (juxtaposition) \\
\hline
5 & Le chef convoqua les anciens \textbf{parce que} la situation était grave. & Phrase complexe (subordonnée conjonctive de cause) \\
\hline
6 & La rivière coulait paisiblement entre les collines. & Phrase simple (1 verbe) \\
\hline
7 & \textbf{Bien qu'}il fût fatigué, le messager poursuivit sa route. & Phrase complexe (subordonnée conjonctive de concession) \\
\hline
8 & Les femmes chantaient, les hommes dansaient, les enfants riaient. & Phrase composée (juxtaposition) \\
\hline
9 & Je pense \textbf{que} la tradition mérite respect. & Phrase complexe (subordonnée conjonctive COD) \\
\hline
10 & Le tam-tam résonna dans la nuit. & Phrase simple (1 verbe) \\
\hline
\end{tabularx}
\end{center}

\textbf{B. Condensation des subordonnées.}
\begin{enumerate}
\item[3.] L'enfant qui portait un panier courait vers l'école. $\rightarrow$ \emph{L'enfant \textbf{portant un panier} courait vers l'école.} (relative $\rightarrow$ participe présent)
\item[5.] Le chef convoqua les anciens parce que la situation était grave. $\rightarrow$ \emph{Le chef convoqua les anciens \textbf{en raison de la gravité de la situation}.} (subordonnée $\rightarrow$ groupe nominal)
\item[7.] Bien qu'il fût fatigué, le messager poursuivit sa route. $\rightarrow$ \emph{\textbf{Malgré sa fatigue}, le messager poursuivit sa route.} (subordonnée $\rightarrow$ groupe nominal introduit par « malgré »)
\item[9.] Je pense que la tradition mérite respect. $\rightarrow$ \emph{Selon l'auteur, \textbf{la tradition mérite respect}.} (suppression de la principale « je pense que », qui marque un point de vue, et conservation de l'idée essentielle)
\end{enumerate}
\end{corrige}

\begin{corrige}[Exercice 8 — Hyperonymie et hyponymie]
\textbf{A. Les trois arbres.}

\begin{popfigure}
\begin{tikzpicture}[level distance=1.2cm, sibling distance=2.6cm, every node/.style={draw, rounded corners, fill=popblueL, font=\small}]
\node {moyens de transport}
  child { node {moto} }
  child { node {taxi} }
  child { node {pirogue} }
  child { node {bus} };
\end{tikzpicture}
\legende{Arbre 1 : l'hyperonyme « moyens de transport » et quatre hyponymes.}
\end{popfigure}

\begin{popfigure}
\begin{tikzpicture}[level distance=1.2cm, sibling distance=2.8cm, every node/.style={draw, rounded corners, fill=popgreenL, font=\small}]
\node {produits du marché}
  child { node {igname} }
  child { node {huile de palme} }
  child { node {arachide} }
  child { node {manioc} };
\end{tikzpicture}
\legende{Arbre 2 : l'hyperonyme « produits du marché » et quatre hyponymes.}
\end{popfigure}

\begin{popfigure}
\begin{tikzpicture}[level distance=1.2cm, sibling distance=2.6cm, every node/.style={draw, rounded corners, fill=poporangeL, font=\small}]
\node {membres de la famille}
  child { node {oncle} }
  child { node {grand-mère} }
  child { node {cousin} }
  child { node {tante} };
\end{tikzpicture}
\legende{Arbre 3 : l'hyperonyme « membres de la famille » et quatre hyponymes.}
\end{popfigure}

\textbf{B. Condensation par l'hyperonyme.}
\begin{enumerate}
\item Au marché de Mokolo, on vend \textbf{des produits du marché} (ou : \textbf{des denrées alimentaires}). L'énumération « ignames, manioc, arachides, huile de palme » est remplacée par le terme général.
\item Pour se rendre à Douala, les voyageurs utilisent \textbf{différents moyens de transport}. L'énumération « bus, taxi, moto, train » est condensée par l'hyperonyme.
\item À la fête du village étaient présents \textbf{tous les membres de ma famille}. L'énumération « oncles, tantes, cousins, grand-mère » est remplacée par l'hyperonyme.
\end{enumerate}
\end{corrige}

\begin{aretenir}[L'intérêt pour le résumé]
Remplacer une énumération par son hyperonyme est une technique de réduction très rentable : on gagne plusieurs mots sans perdre l'idée essentielle.
\end{aretenir}

\begin{corrige}[Exercice 9 — Type Baccalauréat technique : résumé d'un texte argumentatif]
\textbf{1. Résumé proposé (74 mots, dans la fourchette 72--88) :}

\emph{Autrefois, l'éducation traditionnelle transmettait aux enfants les savoir-faire et les valeurs de la communauté. Aujourd'hui, l'école moderne apporte des connaissances et offre des chances de promotion sociale, mais elle éloigne parfois les jeunes de leur culture et ses diplômes ne garantissent plus l'emploi. Il faut donc associer les deux écoles : la première prépare à la vie professionnelle, la seconde enracine dans la culture. Le jeune épanoui maîtrise les savoirs modernes sans renier ses racines.} \textbf{(74 mots)}

Vérification de la contrainte : $80 \times 10\,\% = 8$, donc la fourchette est $[72 ; 88]$ ; $72 \leq 74 \leq 88$ : la consigne est respectée.

\textbf{2. Justification de deux choix de réduction :}
\begin{itemize}
\item \emph{Exemples supprimés} : les détails de la transmission traditionnelle (contes de la grand-mère, pêche du père, cuisine de la mère) et les métiers cités (médecin, ingénieur, enseignant) sont des illustrations ; je les ai condensés en « savoir-faire et valeurs » et en « chances de promotion sociale ».
\item \emph{Énumération condensée} : « lire, écrire, compter, sciences, langues étrangères, techniques » est résumé par l'hyperonyme « connaissances ».
\end{itemize}

\textbf{3. Deux phrases complexes condensées :}
\begin{itemize}
\item « Elle se faisait au quotidien, auprès des parents et des anciens » et la longue phrase « La mère enseignait à ses filles la cuisine\ldots » (phrases à subordonnées et énumérations) $\rightarrow$ condensées en : « l'éducation traditionnelle transmettait les savoir-faire et les valeurs ».
\item « Un jeune Camerounais épanoui est celui qui maîtrise les savoirs modernes sans renier ses racines » (subordonnée relative « qui maîtrise\ldots ») $\rightarrow$ condensée en : « Le jeune épanoui maîtrise les savoirs modernes sans renier ses racines » (suppression de la relative par nominalisation).
\end{itemize}

\textbf{Barème indicatif (sur 10) :} respect de la longueur et nombre de mots indiqué (2 pts) ; fidélité aux quatre idées essentielles : école traditionnelle, apports de l'école moderne, limites, complémentarité (4 pts) ; reformulation sans copie (2 pts) ; neutralité et correction de la langue (2 pts).
\end{corrige}

\begin{attention}[Points de vigilance]
Ne recopie aucune phrase du texte ; n'ajoute aucun avis personnel ; conserve l'ordre des idées de l'auteur ; vérifie le comptage final mot par mot.
\end{attention}

\begin{corrige}[Exercice 10 — Type Probatoire : grammaire et résumé à partir de \emph{Le Lion et la perle}]
\textbf{1. Types d'énoncés :}
\begin{enumerate}
\item « Le lion règne sur la savane depuis toujours. » : \textbf{déclaratif}.
\item « Qui osera contester la force du lion ? » : \textbf{interrogatif}.
\item « Quelle ruse chez ce petit animal qui détient la perle ! » : \textbf{exclamatif}.
\item « Observez comment la sagesse triomphe de la force. » : \textbf{impératif}.
\item « Bien que le lion soit puissant, la perle échappe à ses griffes. » : \textbf{déclaratif}.
\end{enumerate}

\textbf{2. Classement selon la structure :}
\begin{itemize}
\item \emph{Phrase simple} : phrase 1 (un seul verbe conjugué : « règne »).
\item \emph{Phrases complexes} : phrase 3 (« détient » : subordonnée relative « qui détient la perle ») ; phrase 4 (« triomphe » : subordonnée conjonctive introduite par « comment ») ; phrase 5 (subordonnée de concession « Bien que le lion soit puissant » + principale « la perle échappe »).
\item \emph{Phrase simple} également : phrase 2 (un seul verbe conjugué : « osera »).
\item Aucune phrase composée : il n'y a ni coordination ni juxtaposition.
\end{itemize}

\textbf{3. Transformations en énoncés déclaratifs neutres :}
\begin{itemize}
\item (2) $\rightarrow$ \emph{Personne n'ose contester la force du lion.}
\item (3) $\rightarrow$ \emph{Le petit animal qui détient la perle fait preuve de ruse.}
\item (4) $\rightarrow$ \emph{La sagesse triomphe de la force.}
\end{itemize}

\textbf{4. Mini-résumé (3 phrases) :}

\emph{Le lion règne sur la savane et personne n'ose contester sa force. Pourtant, le petit animal détient la perle grâce à sa ruse : bien que le lion soit puissant, la perle échappe à ses griffes. Ainsi, dans l'œuvre, la sagesse triomphe de la force.}

\textbf{Barème indicatif (sur 10) :} types d'énoncés (2,5 pts) ; classement justifié par le nombre de propositions (2,5 pts) ; transformations neutres (3 pts) ; mini-résumé fidèle et correct (2 pts).
\end{corrige}

\begin{attention}[Points de vigilance]
Pour justifier le classement, cite toujours le nombre de verbes conjugués et la nature du lien (coordination, juxtaposition, subordination). Dans le mini-résumé, aucune exclamation ni question ne doit subsister.
\end{attention}

\begin{corrige}[Exercice 11 — Type Baccalauréat technique : remédiation d'un résumé fautif]
\textbf{1. et 2. Les cinq défauts et les règles enfreintes :}

\begin{center}
\begin{tabularx}{\linewidth}{|c|X|X|}
\hline
\rowcolor{popblueL} N° & Passage fautif & Règle enfreinte \\
\hline
1 & 95 mots au lieu de 60 ($\pm 10\,\%$, soit 54 à 66 mots) & Non-respect de la contrainte de longueur \\
\hline
2 & « Le conte est très important\ldots l'honnêteté. » (phrase copiée mot à mot) & Interdiction de copier : le résumé exige la reformulation \\
\hline
3 & « Moi je trouve que les contes sont magnifiques\ldots ont tort\ldots » & Jugement personnel : le résumé doit rester neutre et objectif \\
\hline
4 & « Le conteur est assis près du feu, la lune brille\ldots » & Description pittoresque inutile : seules les idées essentielles sont conservées \\
\hline
5 & « quel beau spectacle ! Il faut absolument que tout le monde lise ce texte formidable. » & Phrases exclamative et injonctive inadaptées : le résumé n'utilise que la phrase déclarative \\
\hline
\end{tabularx}
\end{center}

\textbf{3. Version corrigée (54 mots, dans la fourchette 54--66) :}

\emph{Le conte joue un rôle essentiel dans l'éducation des enfants : il transmet les valeurs de la société et développe leur imagination. En réunissant la famille autour de la parole des anciens, il crée aussi un lien entre les générations. Ainsi, le conte éduque tout en divertissant, et il mérite d'être transmis aux jeunes générations.} \textbf{(54 mots)}

Vérification : $60 \times 10\,\% = 6$, fourchette $[54 ; 66]$ ; $54 \leq 54 \leq 66$ : la contrainte est respectée. Aucune phrase copiée, aucun jugement personnel, uniquement des phrases déclaratives.

\textbf{Barème indicatif (sur 10) :} repérage des cinq défauts avec citations (5 pts) ; règles énoncées correctement (2,5 pts) ; résumé corrigé conforme (longueur, neutralité, reformulation) (2,5 pts).
\end{corrige}

\begin{attention}[Points de vigilance]
Chaque défaut doit être \emph{cité} précisément, pas seulement décrit. Dans la version corrigée, indique le nombre de mots : l'oublier fait perdre des points même si le résumé est bon.
\end{attention}

\newpage

\coursec{Fiche bilan}

\begin{popfigure}
\begin{tikzpicture}[mindmap, grow cyclic, every node/.style={concept, font=\small},
  root concept/.append style={concept color=popblue, font=\bfseries},
  level 1/.append style={level distance=3.4cm, sibling angle=51.4},
  level 1 concept/.append style={font=\small},
  text width=2.6cm, align=center]
\node[root concept, text width=3cm]{Rédiger un résumé de texte}
  child[concept color=poporange]{node[align=center]{L'œuvre support\\ \emph{Le Lion et la perle} : exposé, contexte, bilan}}
  child[concept color=popgreen]{node[align=center]{Le résumé\\ définition, objectifs, contraintes (1/4 du texte)}}
  child[concept color=poppurple]{node[align=center]{Types de phrases\\ déclarative, interrogative, impérative, exclamative}}
  child[concept color=poppink]{node[align=center]{Structure\\ phrase simple, composée, complexe}}
  child[concept color=popgold]{node[align=center]{Lexique\\ hypéronymie et hyponymie}}
  child[concept color=popblue!60]{node[align=center]{Réduction\\ suppression, généralisation, condensation}}
  child[concept color=popmuted]{node[align=center]{Rédaction\\ reformuler, enchaîner, corriger, évaluer}};
\end{tikzpicture}
\legende{Carte mentale de la leçon : les sept savoirs à mobiliser pour réussir un résumé au Probatoire et au Baccalauréat technique.}
\end{popfigure}

\begin{aretenir}[L'essentiel de la leçon — définitions clés]
\begin{itemize}
  \item \cle{Résumé} : texte bref (environ le quart du texte de départ) qui en restitue les idées essentielles, avec des mots personnels, sans commentaire ni jugement.
  \item \cle{Contraction de texte} : exercice qui consiste à réduire un texte en respectant son sens, sa progression et la longueur imposée.
  \item \cle{Hypéronyme} : mot au sens général qui englobe d'autres mots (ex. \emph{animal}) ; \cle{hyponyme} : mot au sens plus précis (ex. \emph{lion}). La généralisation par l'hypéronyme est un outil majeur de réduction.
\end{itemize}
\end{aretenir}

\begin{aretenir}[Grammaire à connaître par cœur]
\begin{center}
\begin{tabularx}{\linewidth}{|l|X|X|}
\hline
\rowcolor{popblueL} \textbf{Notion} & \textbf{Définition} & \textbf{Exemple} \\
\hline
Phrase déclarative & donne une information (point final) & Le lion rugit. \\
\hline
Phrase interrogative & pose une question (totale ou partielle) & Qui a tué le lion ? \\
\hline
Phrase impérative & ordre, conseil, prière & Écoutez ce récit ! \\
\hline
Phrase exclamative & exprime une émotion & Quel courage ! \\
\hline
Phrase simple & un seul verbe conjugué & Sadike chante. \\
\hline
Phrase composée & plusieurs propositions juxtaposées ou coordonnées (\emph{et, ou, mais}) & Il chante et elle danse. \\
\hline
Phrase complexe & proposition principale + subordonnée (relative, conjonctive, circonstancielle) & Je sais qu'il viendra. \\
\hline
\end{tabularx}
\end{center}
\end{aretenir}

\begin{aretenir}[Les techniques de réduction]
\begin{itemize}
  \item \cle{Suppression} : éliminer répétitions, exemples, citations, comparaisons, détails, énumérations, expansions du nom.
  \item \cle{Généralisation} : remplacer plusieurs hyponymes par un hypéronyme (\emph{pagnes, boubous, chemises} $\rightarrow$ \emph{vêtements}).
  \item \cle{Condensation} : fondre plusieurs phrases en une seule ; transformer une subordonnée en groupe nominal (\emph{quand il arriva} $\rightarrow$ \emph{à son arrivée}).
\end{itemize}
\end{aretenir}

\begin{methode}[Méthode express en 5 étapes]
\begin{enumerate}
  \item Lire deux fois le texte ; souligner les mots-clés et repérer la progression des idées.
  \item Dégager l'idée essentielle de chaque paragraphe (plan du texte).
  \item Réduire : supprimer, généraliser, condenser.
  \item Rédiger avec ses propres mots, en phrases claires reliées par des connecteurs (\emph{d'abord, ensuite, enfin}).
  \item Vérifier : nombre de mots respecté, fidélité au texte, orthographe et ponctuation.
\end{enumerate}
\end{methode}

\begin{aretenir}[L'œuvre support]
\begin{itemize}
  \item \emph{Le Lion et la perle} (Wole Soyinka) : pièce nigériane qui met en scène le conflit entre la tradition (Baroka, le lion) et la modernité (Lakunle) autour de Sidi, la perle du village.
  \item L'œuvre s'inscrit dans le contexte de l'Afrique postcoloniale : affrontement des valeurs, satire des prétentions modernistes, éloge nuancé de la tradition.
\end{itemize}
\end{aretenir}

\begin{aretenir}[Checklist avant le Bac]
Avant de rendre votre résumé, vérifiez chaque point :
\begin{itemize}
  \item[$\square$] J'ai lu le texte au moins deux fois et repéré son thème.
  \item[$\square$] J'ai dégagé une idée essentielle par paragraphe.
  \item[$\square$] J'ai supprimé exemples, citations, répétitions et détails.
  \item[$\square$] J'ai généralisé les énumérations par des hypéronymes.
  \item[$\square$] J'ai condensé les subordonnées en groupes nominaux.
  \item[$\square$] Mon résumé respecte la longueur imposée (environ le quart du texte).
  \item[$\square$] J'ai reformulé avec mes propres mots, sans copier de phrases entières.
  \item[$\square$] Je n'ai ajouté ni opinion personnelle ni commentaire.
  \item[$\square$] Mes phrases sont reliées par des connecteurs logiques.
  \item[$\square$] Je sais distinguer phrase simple, composée et complexe dans mes propres phrases.
  \item[$\square$] J'ai relu pour corriger orthographe, accords et ponctuation.
  \item[$\square$] Je connais l'intrigue et le contexte du \emph{Lion et la perle}.
\end{itemize}
\end{aretenir}

\findecours

\end{document}
